%==============================================================================
%  Decisiones de tecnología (SAD v1.2). Archivo incluido desde main.tex con %==============================================================================
%  Decisiones de tecnología (SAD v1.2). Archivo incluido desde main.tex con %==============================================================================
%  Decisiones de tecnología (SAD v1.2). Archivo incluido desde main.tex con %==============================================================================
%  Decisiones de tecnología (SAD v1.2). Archivo incluido desde main.tex con \input{tecnologia}.
%  Generado a partir de la especificación de la versión 1.2. No carga paquetes.
%==============================================================================
\section{Decisiones de tecnología}
\label{sec:tecnologia}

El enunciado del proyecto (§8) pide seleccionar las tecnologías en la segunda entrega.
CON-8 exigía independencia tecnológica en la primera entrega.
Esta sección sustituye esa condición para la segunda entrega.
El Anexo~\ref{sec:anexo-a} deja de ser una excepción a CON-8 y describe la implementación de FP-06 con las tecnologías elegidas aquí.

Las decisiones se numeran de AD-50 a AD-66.
AD-66 responde a CON-9, que se agregó después de AD-50 a AD-65.
Ninguna modifica una decisión estructural.
Cada una materializa uno o varios mecanismos de AD-01 a AD-49, y su justificación parte de los escenarios QA-nn que ese mecanismo debe cumplir (CON-7).

\subsection{Método de decisión}
\label{sec:tec-metodo}

Cada decisión sigue cinco pasos:
\begin{enumerate}
  \item Se identifica el mecanismo del SAD que la tecnología materializa y los escenarios QA-nn que ese mecanismo atiende.
  \item Se eligen entre tres y cinco alternativas con uso real para ese mecanismo.
  \item Se definen los criterios. Cada criterio cita el QA-nn, AD-nn o CON-n del que proviene. La experiencia del equipo y la licencia o el costo se incluyen con un peso menor que los criterios de calidad.
  \item Se asigna un peso a cada criterio. Los pesos de una decisión suman 100\,\%.
  \item Se califica cada alternativa en cada criterio y se calcula el total ponderado.
\end{enumerate}

La escala de calificación tiene cinco niveles:
\begin{itemize}
  \item 1: la alternativa no satisface el criterio o lo impide.
  \item 2: lo satisface con desarrollo propio considerable.
  \item 3: lo satisface con configuración adicional o con complementos de terceros.
  \item 4: lo satisface de forma nativa con limitaciones menores.
  \item 5: lo satisface de forma nativa y la implementación lo verifica.
\end{itemize}

El total ponderado de una alternativa es $\sum_i p_i \, c_i / 100$, donde $p_i$ es el peso en porcentaje y $c_i$ la calificación del criterio $i$.
El máximo es 5.
Las cifras de verificación citadas provienen del perfil local de un computador, con 20 clientes, y no reemplazan la prueba con 10\,000 usuarios (L-04).

\subsection{Resumen de las decisiones}
\label{sec:tec-resumen}

El \cref{tab:tecnologia-resumen} relaciona cada decisión con su tecnología, el mecanismo que materializa y los escenarios que favorece.
Las versiones son las de los archivos de proyecto.

\begingroup\small
\begin{longtable}{L{1.15cm} L{2.2cm} L{6.0cm} L{2.6cm} L{2.6cm}}
\cabecera{5}{tab:tecnologia-resumen}{Resumen de las decisiones de tecnología}{\textbf{ID} & \textbf{Categoría} &
\textbf{Tecnología y versión} & \textbf{Mecanismo} & \textbf{QA que favorece}}
AD-50 & Lenguaje y plataforma & C\# sobre .NET 10 (SDK 10.0.100, \texttt{net10.0}) & AD-01, AD-08, AD-09, AD-45 & QA-06, QA-07, QA-14 \\
AD-51 & Marco web & ASP.NET Core 10 (controladores y \emph{minimal APIs}) & AD-08, AD-22, AD-24 & QA-06, QA-07, QA-11 \\
AD-52 & Persistencia & PostgreSQL 18 con PostGIS 3.6 (\texttt{postgis/postgis:18-3.6}), una base por microservicio y réplica \texttt{db-replica} & AD-01, AD-17, AD-20, AD-24, AD-41, AD-45 & QA-03, QA-04, QA-06, QA-13, QA-15 \\
AD-53 & Acceso a datos & Entity Framework Core 10.0.14, Npgsql 10.0.3 y NetTopologySuite 10.0.3, con contenedor migrador & AD-08, AD-17, AD-41, AD-45 & QA-06, QA-10, QA-13 \\
AD-54 & Mensajería & RabbitMQ 4.4 (3 nodos, colas quorum acotadas, DLQ) con RabbitMQ.Client 7.2.2 & AD-02, AD-05, AD-36, AD-45 & QA-01, QA-04, QA-08 \\
AD-55 & Caché & Redis 8.8 (\texttt{redis:8.8.4-alpine}) con StackExchange.Redis 3.4.29 & AD-16, AD-17 & QA-05, QA-06 \\
AD-56 & Almacén de objetos & MinIO con API S3 (paquete Minio 7.0.0) & AD-24, AD-31 & QA-02, QA-03, QA-14 \\
AD-57 & Puerta de entrada & YARP 2.3.0 con catálogo \texttt{rutas.json} (56 rutas) & AD-04, AD-19, AD-22 & QA-05, QA-06, QA-11 \\
AD-58 & Identidad & Keycloak 26.5.2 (realm \texttt{plataforma}), AppAuth-Android 0.11.1, angular-auth-oidc-client 22.0.2, cliente OpenID Connect de ASP.NET Core para el tablero y validador JWT propio en \texttt{Plataforma.Comun.Seguridad} & AD-04, AD-22, AD-49 & QA-11, QA-14 \\
AD-59 & Cliente web & Angular 22.2 (TypeScript 6.0) servido por nginx 1.29 detrás del gateway & AD-04, AD-08 & QA-06, QA-10, QA-11 \\
AD-60 & Mapas & Leaflet 1.9.4 con teselas de OpenStreetMap & AD-16, AD-17 & QA-06 \\
AD-61 & Cliente móvil & Android nativo: Kotlin 2.0.21, Jetpack Compose (BOM 2024.09.00), Hilt 2.60.1, minSdk 26, targetSdk 36 & AD-31, AD-44, AD-48 & QA-07, QA-14 \\
AD-62 & Persistencia móvil & Room 2.8.5 sobre SQLite, migraciones explícitas 1$\to$2$\to$3 & AD-31, AD-44, AD-48 & QA-14 \\
AD-63 & Sincronización móvil & WorkManager 2.12.0, Retrofit 3.0.0, OkHttp 5.4.0 y kotlinx.serialization 1.9.0 & AD-31, AD-44, AD-48 & QA-11, QA-14 \\
AD-64 & Resiliencia & Microsoft.Extensions.Http.Resilience 10.11.0 (Polly v8) & AD-05, AD-37 & QA-01, QA-04 \\
AD-65 & Despliegue & Contenedores Docker (\texttt{mcr.microsoft.com/\br{}dotnet/\br{}aspnet:10.0}), Docker Compose para un computador y Kubernetes con HPA y KEDA para escalar & AD-18, AD-24, AD-25 & QA-01, QA-02, QA-05 \\
AD-66 & Adaptador de SE-8 & Jakarta EE 11 sobre WildFly 42.0.0.Final (Java 21): Jakarta Messaging con el cliente JMS de RabbitMQ 3.9.1, EJB, Jakarta Persistence, Jakarta REST y MicroProfile Config & AD-03, AD-05, AD-09, AD-25, AD-45 & QA-08, QA-09, CON-9 \\
\end{longtable}
\endgroup

%---------------------------------------------------------------------- AD-50
\subsection{AD-50. Lenguaje y plataforma de los microservicios}
\label{sec:tec-ad50}

\paragraph{Contexto.} Esta decisión materializa los microservicios con almacén propio (AD-01), las cuatro capas de cada microservicio (AD-08), el contrato común de los adaptadores (AD-09) y la bandeja de salida transaccional (AD-45).
El lenguaje determina la latencia del camino síncrono, que exige QA-06 (p95 $<$ 2\,s) y QA-07 (p95 $<$ 3\,s).
También determina si la base común \texttt{comun} y el prototipo de FP-06 del Anexo~\ref{sec:anexo-a} se reutilizan, lo que afecta la verificación de QA-14.
CON-2 exige integrarse con sistemas de tecnologías diferentes, y esa heterogeneidad la absorben los adaptadores (AD-03), no el lenguaje de los microservicios.

\paragraph{Alternativas.} Se consideran C\# sobre .NET 10, Java 21, Node.js 22 con TypeScript y Go.
Las cuatro tienen licencia libre y ejecución en contenedores.
El \cref{tab:tec-ad50} presenta la matriz de decisión ponderada.
La alternativa elegida obtiene 4,90 y la segunda mejor, Java 21, obtiene 3,85.

\begingroup\small
\begin{longtable}{L{4.96cm} C{1.10cm} C{2.25cm} C{2.25cm} C{2.25cm} C{2.25cm}}
\cabecera{6}{tab:tec-ad50}{Matriz de decisión ponderada de AD-50}{\textbf{Criterio (driver)} & \textbf{Peso} & \hspace{0pt}\textbf{.NET 10} & \hspace{0pt}\textbf{Java 21} & \hspace{0pt}\textbf{Node.js 22} & \hspace{0pt}\textbf{Go}}
Rendimiento del camino síncrono (QA-06, QA-07) & 25\,\% & 5 & 4 & 3 & 5 \\
Bibliotecas maduras para bus, resiliencia, datos espaciales y JWT (AD-02, AD-05, AD-17, AD-22) & 25\,\% & 5 & 5 & 3 & 3 \\
Reutilización de \texttt{comun} y del prototipo de FP-06 (QA-14, AD-09) & 20\,\% & 5 & 2 & 2 & 2 \\
Memoria y arranque de cada réplica (QA-05, AD-18) & 10\,\% & 4 & 3 & 4 & 5 \\
Experiencia del equipo & 10\,\% & 5 & 4 & 3 & 1 \\
Licencia y costo & 10\,\% & 5 & 5 & 5 & 5 \\
\midrule
\textbf{Total ponderado} & 100\,\% & \textbf{4,90} & 3,85 & 3,10 & 3,50 \\
\end{longtable}
\endgroup

\paragraph{Alternativas descartadas.}
\begin{itemize}
  \item Java 21 pierde en la reutilización de \texttt{comun} y del prototipo de FP-06 (QA-14), que habría que reescribir.
  \item Node.js 22 pierde en el rendimiento del camino síncrono (QA-06, QA-07) y en la madurez de las bibliotecas para datos espaciales y bus.
  \item Go gana en memoria y arranque de cada réplica (QA-05), pero pierde en bibliotecas para los mecanismos del SAD y en experiencia del equipo.
\end{itemize}

\paragraph{Compromiso aceptado.} Una sola pila reduce el costo de mantener \texttt{comun} (AD-09) y acepta que todos los microservicios dependen del ciclo de soporte de .NET.
La heterogeneidad que exige CON-2 queda confinada en los adaptadores (AD-03).
El núcleo sigue en .NET aun con CON-9: Jakarta EE se confina a la frontera según AD-03 y AD-66.

\paragraph{Riesgo.} El fin del soporte de .NET 10 obligaría a migrar todos los microservicios a la vez.
El archivo \texttt{global.json} fija el SDK con avance solo de la versión de características, lo que acota este riesgo.
No corresponde a un riesgo del ATAM.

\paragraph{Dónde se ve.} Vista de implementación: proyectos \texttt{*.Api}, \texttt{*.Aplicacion}, \texttt{*.Dominio} e \texttt{*.Infraestructura} de cada microservicio.
Vista física de la implementación: artefactos \texttt{plataforma/<servicio>:local} sobre \texttt{mcr.microsoft.com/\br{}dotnet/\br{}aspnet:10.0}.

\paragraph{Evidencia en el código.} La suite de .NET suma 1\,008 pruebas en 19 proyectos en verde. QA-07 se midió con p95 de 1\,334\,ms en el perfil local (L-04).
Archivos de evidencia:
{\raggedright
\begin{itemize}
  \item \texttt{proyectoparte2/\br{}global.json}
  \item \texttt{proyectoparte2/\br{}Directory.Build.props}
\end{itemize}
\par}

%---------------------------------------------------------------------- AD-51
\subsection{AD-51. Marco de las API de los microservicios}
\label{sec:tec-ad51}

\paragraph{Contexto.} Esta decisión materializa la capa de acceso de cada microservicio (AD-08), el nivel 2 de la autorización (AD-22) y el sondeo de salud de la réplica activa (AD-24).
El marco atiende la latencia de QA-06 y QA-07 y el rechazo del 100\,\% de las modificaciones no autorizadas de QA-11.
La elección está condicionada por AD-50, porque el marco debe pertenecer a la pila elegida.

\paragraph{Alternativas.} Se consideran ASP.NET Core 10, Spring Boot 3 y NestJS.
Spring Boot 3 y NestJS se evalúan como el marco que se habría usado con Java y con Node.js.
El \cref{tab:tec-ad51} presenta la matriz de decisión ponderada.
La alternativa elegida obtiene 4,85 y la segunda mejor, Spring Boot 3, obtiene 3,75.

\begingroup\small
\begin{longtable}{L{6.74cm} C{1.10cm} C{2.50cm} C{2.50cm} C{2.50cm}}
\cabecera{5}{tab:tec-ad51}{Matriz de decisión ponderada de AD-51}{\textbf{Criterio (driver)} & \textbf{Peso} & \hspace{0pt}\textbf{ASP.NET Core 10} & \hspace{0pt}\textbf{Spring Boot 3} & \hspace{0pt}\textbf{NestJS}}
Latencia por petición (QA-06, QA-07) & 25\,\% & 5 & 4 & 3 \\
Validación JWT y autorización de nivel 2 en la tubería (QA-11, AD-22) & 20\,\% & 5 & 5 & 4 \\
Sondeos de salud vivo y listo (QA-02, QA-04, AD-24) & 15\,\% & 4 & 5 & 4 \\
Coherencia con AD-50 y con el gateway (QA-10, AD-57) & 20\,\% & 5 & 1 & 1 \\
Experiencia del equipo & 10\,\% & 5 & 3 & 3 \\
Licencia y costo & 10\,\% & 5 & 5 & 5 \\
\midrule
\textbf{Total ponderado} & 100\,\% & \textbf{4,85} & 3,75 & 3,15 \\
\end{longtable}
\endgroup

\paragraph{Alternativas descartadas.}
\begin{itemize}
  \item Spring Boot 3 gana en sondeos de salud con Actuator (QA-02), pero pierde en coherencia con AD-50, porque introduce una segunda pila.
  \item NestJS pierde en latencia por petición (QA-06, QA-07) y en coherencia con AD-50.
\end{itemize}

\paragraph{Compromiso aceptado.} Los sondeos de salud se componen con \texttt{Plataforma.Comun.Salud} en lugar de un módulo integrado como Actuator (AD-24).

\paragraph{Riesgo.} Una actualización mayor del marco afecta a todos los microservicios a la vez.
No corresponde a un riesgo del ATAM.

\paragraph{Dónde se ve.} Vista de implementación: proyectos \texttt{*.Api} con \texttt{Sdk.Web}.
Vista física de la implementación: artefactos \texttt{plataforma/<servicio>:local}.

\paragraph{Evidencia en el código.} Los proyectos \texttt{*.Api.csproj} usan \texttt{Microsoft.NET.Sdk.Web}. QA-07 se cumple con 100 de 100 registros respondidos con 201.
Archivos de evidencia:
{\raggedright
\begin{itemize}
  \item \texttt{proyectoparte2/\br{}servicios/\br{}gestion-desastres/\br{}src/\br{}GestionDesastres.Api}
  \item \texttt{proyectoparte2/\br{}comun/\br{}Plataforma.Comun.Salud}
\end{itemize}
\par}

%---------------------------------------------------------------------- AD-52
\subsection{AD-52. Motor de base de datos}
\label{sec:tec-ad52}

\paragraph{Contexto.} Esta decisión materializa el almacén propio de cada microservicio (AD-01), la geometría precalculada de las zonas (AD-17), los historiales versionados y de solo inserción (AD-20, AD-41), la réplica activa (AD-24) y la bandeja de salida (AD-45).
El motor atiende la consulta espacial de QA-06, la publicación sin pérdida de QA-04, la inmutabilidad de QA-13 y QA-15 y el RPO $\le$ 60\,s de QA-03.

\paragraph{Alternativas.} Se consideran PostgreSQL con PostGIS, MySQL 8, SQL Server y MongoDB.
El \cref{tab:tec-ad52} presenta la matriz de decisión ponderada.
La alternativa elegida obtiene 4,85 y la segunda mejor, SQL Server, obtiene 4,10.

\begingroup\small
\begin{longtable}{L{4.96cm} C{1.10cm} C{2.25cm} C{2.25cm} C{2.25cm} C{2.25cm}}
\cabecera{6}{tab:tec-ad52}{Matriz de decisión ponderada de AD-52}{\textbf{Criterio (driver)} & \textbf{Peso} & \hspace{0pt}\textbf{PostgreSQL + PostGIS} & \hspace{0pt}\textbf{MySQL 8} & \hspace{0pt}\textbf{SQL Server} & \hspace{0pt}\textbf{MongoDB}}
Consulta espacial con índice (QA-06, AD-17) & 25\,\% & 5 & 3 & 4 & 3 \\
Transacción ACID para la bandeja y el historial (QA-04, AD-20, AD-45) & 20\,\% & 5 & 4 & 5 & 3 \\
Inmutabilidad garantizada por permisos del almacén (QA-13, QA-15, AD-41) & 15\,\% & 5 & 4 & 5 & 3 \\
Réplica continua para la réplica activa y el RPO (QA-03, AD-24, AD-27) & 15\,\% & 4 & 4 & 5 & 5 \\
Experiencia del equipo & 10\,\% & 5 & 4 & 3 & 3 \\
Licencia y costo & 15\,\% & 5 & 4 & 2 & 3 \\
\midrule
\textbf{Total ponderado} & 100\,\% & \textbf{4,85} & 3,75 & 4,10 & 3,30 \\
\end{longtable}
\endgroup

\paragraph{Alternativas descartadas.}
\begin{itemize}
  \item MySQL 8 pierde en consulta espacial con índice (QA-06), porque sus funciones espaciales son menos completas que PostGIS.
  \item SQL Server gana en réplica continua (QA-03), pero pierde en licencia y costo, porque su licencia comercial se cobra por núcleo.
  \item MongoDB gana en réplica continua (QA-03), pero pierde en transacción ACID para la bandeja (QA-04) y en inmutabilidad por permisos (QA-13).
\end{itemize}

\paragraph{Compromiso aceptado.} En el entorno de un computador las nueve bases comparten la instancia \texttt{db}.
La separación por microservicio se conserva por base y por rol con mínimo privilegio (AD-01), no por servidor.

\paragraph{Riesgo.} La instancia \texttt{db} es un punto único de fallo en el entorno de un computador.
La réplica \texttt{db-replica} reduce la pérdida de datos, pero la promoción no es automática y no hay segunda región (L-01).
QA-02 y QA-03 no se midieron.

\paragraph{Dónde se ve.} Vista física de la implementación: nodos \texttt{db} y \texttt{db-replica}.
Vista de implementación: paquete \texttt{Infraestructura/Persistencia} de cada microservicio y migradores.

\paragraph{Evidencia en el código.} Los servicios \texttt{db} y \texttt{db-replica} usan la imagen \texttt{postgis/postgis:18-3.6}. QA-06 se midió con p95 de 34\,ms en el consolidado del desastre en el perfil local (L-04).
Archivos de evidencia:
{\raggedright
\begin{itemize}
  \item \texttt{proyectoparte2/\br{}despliegue/\br{}compose/\br{}docker-compose.yml}
\end{itemize}
\par}

%---------------------------------------------------------------------- AD-53
\subsection{AD-53. Acceso a datos y migraciones}
\label{sec:tec-ad53}

\paragraph{Contexto.} Esta decisión materializa la capa de infraestructura de cada microservicio (AD-08), el mapeo de las geometrías (AD-17), los permisos del historial de solo inserción (AD-41) y la escritura atómica con la bandeja (AD-45).
El acceso a datos afecta la latencia de QA-06, la evolución del esquema que pide QA-10 y la inmutabilidad de QA-13.

\paragraph{Alternativas.} Se consideran Entity Framework Core, Dapper, NHibernate y SQL directo con Npgsql sin mapeador.
El \cref{tab:tec-ad53} presenta la matriz de decisión ponderada.
La alternativa elegida obtiene 4,80 y la segunda mejor, Dapper, obtiene 3,60.

\begingroup\small
\begin{longtable}{L{4.96cm} C{1.10cm} C{2.25cm} C{2.25cm} C{2.25cm} C{2.25cm}}
\cabecera{6}{tab:tec-ad53}{Matriz de decisión ponderada de AD-53}{\textbf{Criterio (driver)} & \textbf{Peso} & \hspace{0pt}\textbf{EF Core 10} & \hspace{0pt}\textbf{Dapper} & \hspace{0pt}\textbf{NHibernate} & \hspace{0pt}\textbf{SQL directo}}
Migraciones versionadas con los permisos del historial (QA-10, QA-13, AD-41) & 25\,\% & 5 & 2 & 3 & 1 \\
Tipos espaciales mapeados (QA-06, AD-17) & 20\,\% & 5 & 3 & 2 & 3 \\
Control del SQL en las lecturas críticas (QA-06) & 20\,\% & 4 & 5 & 3 & 5 \\
Aislamiento en la capa de infraestructura y pruebas (QA-10, AD-08) & 15\,\% & 5 & 4 & 4 & 3 \\
Experiencia del equipo & 10\,\% & 5 & 4 & 2 & 4 \\
Licencia y costo & 10\,\% & 5 & 5 & 5 & 5 \\
\midrule
\textbf{Total ponderado} & 100\,\% & \textbf{4,80} & 3,60 & 3,05 & 3,20 \\
\end{longtable}
\endgroup

\paragraph{Alternativas descartadas.}
\begin{itemize}
  \item Dapper gana en control del SQL (QA-06), pero pierde en migraciones versionadas (QA-10, QA-13), que exigirían otra herramienta.
  \item NHibernate pierde en tipos espaciales mapeados (QA-06) y en experiencia del equipo.
  \item SQL directo gana en control del SQL (QA-06), pero pierde en migraciones versionadas y en aislamiento de la capa de infraestructura (QA-10).
\end{itemize}

\paragraph{Compromiso aceptado.} El mapeador agrega una capa sobre el SQL.
Las lecturas críticas del modelo de lectura (AD-15) pueden usar SQL explícito dentro del mismo contexto cuando QA-06 lo exija.
El esquema lo aplica un contenedor migrador con el rol dueño, y el microservicio corre con un rol sin permisos de DDL (AD-41).

\paragraph{Riesgo.} Una consulta generada por el mapeador puede degradar la latencia del mapa de daños, que se midió con p95 de 1\,628\,ms frente al límite de 2\,s de QA-06.
No corresponde a un riesgo del ATAM.

\paragraph{Dónde se ve.} Vista de implementación: paquete \texttt{Infraestructura/Persistencia} y migradores (\texttt{*-migrador}).
Vista física de la implementación: nodos \texttt{db} y \texttt{db-replica}.

\paragraph{Evidencia en el código.} El servicio \texttt{reconstruccion-migrador} del archivo de composición crea los roles y aplica las migraciones antes de arrancar el microservicio.
Archivos de evidencia:
{\raggedright
\begin{itemize}
  \item \texttt{proyectoparte2/\br{}Directory.Packages.props}
  \item \texttt{proyectoparte2/\br{}comun/\br{}Plataforma.Comun.Migraciones}
  \item \texttt{proyectoparte2/\br{}servicios/\br{}evaluacion-danos/\br{}src/\br{}EvaluacionDanos.Infraestructura/\br{}Persistencia}
\end{itemize}
\par}

%---------------------------------------------------------------------- AD-54
\subsection{AD-54. Bus de eventos}
\label{sec:tec-ad54}

\paragraph{Contexto.} Esta decisión materializa el bus de eventos con colas por consumidor (AD-02), las colas acotadas de la frontera (AD-05), la cola de salida con reintento (AD-36) y la entrega desde la bandeja de salida (AD-45).
El bus atiende QA-04, cuya respuesta exige 0 eventos perdidos ante el fallo de un componente interno, y QA-01, que exige procesar el 100\,\% de lo encolado cuando el sistema externo vuelve.

\paragraph{Alternativas.} Se consideran RabbitMQ, Apache Kafka, NATS JetStream y MassTransit sobre RabbitMQ.
MassTransit se evalúa como capa de abstracción sobre el mismo intermediario.
El \cref{tab:tec-ad54} presenta la matriz de decisión ponderada.
La alternativa elegida obtiene 4,90 y la segunda mejor, NATS JetStream, obtiene 3,85.

\begingroup\small
\begin{longtable}{L{4.96cm} C{1.10cm} C{2.25cm} C{2.25cm} C{2.25cm} C{2.25cm}}
\cabecera{6}{tab:tec-ad54}{Matriz de decisión ponderada de AD-54}{\textbf{Criterio (driver)} & \textbf{Peso} & \hspace{0pt}\textbf{RabbitMQ 4.4} & \hspace{0pt}\textbf{Kafka} & \hspace{0pt}\textbf{NATS JetStream} & \hspace{0pt}\textbf{MassTransit}}
Cola acotada por consumidor con rechazo al productor (QA-04, AD-05) & 25\,\% & 5 & 2 & 4 & 4 \\
Confirmación por mensaje, límite de entregas y DLQ (QA-01, QA-04, AD-36) & 20\,\% & 5 & 2 & 4 & 5 \\
Replicación de la cola ante la caída de un nodo (QA-02, QA-04) & 15\,\% & 5 & 5 & 5 & 5 \\
Topología explícita y declarada al desplegar (QA-09, QA-10) & 15\,\% & 5 & 3 & 3 & 2 \\
Experiencia del equipo & 10\,\% & 4 & 2 & 1 & 3 \\
Licencia y costo & 15\,\% & 5 & 5 & 5 & 2 \\
\midrule
\textbf{Total ponderado} & 100\,\% & \textbf{4,90} & 3,05 & 3,85 & 3,65 \\
\end{longtable}
\endgroup

\paragraph{Alternativas descartadas.}
\begin{itemize}
  \item Kafka iguala en replicación (QA-04), pero pierde en cola acotada por consumidor y en confirmación por mensaje con DLQ, que no ofrece de forma nativa.
  \item NATS JetStream pierde en topología declarada (QA-09) y en experiencia del equipo.
  \item MassTransit gana en confirmación y DLQ (QA-04), pero pierde en topología explícita, porque la oculta, y en licencia, porque sus autores anunciaron licencia comercial para la versión 9.
\end{itemize}

\paragraph{Compromiso aceptado.} La retención de un registro de eventos, que Kafka ofrece, no existe en RabbitMQ.
La recuperación de los mensajes pendientes descansa en la bandeja de salida de cada productor (AD-45).
La política \texttt{plataforma-colas} acota cada cola a 100\,000 mensajes y rechaza la publicación al llenarse.

\paragraph{Riesgo.} Una cola llena rechaza la publicación y el relevo de la bandeja reintenta.
Si la saturación persiste, la bandeja crece en la base del productor.
La alarma de la bandeja no se registra en Auditoría (L-02).

\paragraph{Dónde se ve.} Vista de procesos: exchange \texttt{plataforma.eventos} y colas \texttt{q.*} con su \texttt{dlq.*}.
Vista física de la implementación: nodos \texttt{bus1}, \texttt{bus2} y \texttt{bus3}.

\paragraph{Evidencia en el código.} Las 71 pruebas de \texttt{Plataforma.Comun.Pruebas} con RabbitMQ real verifican consumidor detenido sin pérdida, mensaje envenenado apartado y cola llena sin descarte (QA-04).
Archivos de evidencia:
{\raggedright
\begin{itemize}
  \item \texttt{proyectoparte2/\br{}despliegue/\br{}bus/\br{}definiciones.json}
  \item \texttt{proyectoparte2/\br{}comun/\br{}Plataforma.Comun.Mensajeria}
\end{itemize}
\par}

%---------------------------------------------------------------------- AD-55
\subsection{AD-55. Motor de caché}
\label{sec:tec-ad55}

\paragraph{Contexto.} Esta decisión materializa la caché de geometrías y de respuestas del mapa (AD-17) y la marca de tiempo del dato consultado (AD-16).
El motor atiende QA-06, que exige p95 $<$ 2\,s y antigüedad del dato $\le$ 2\,s en el p95.
También atiende QA-05, porque la caché separa la carga de lectura de la base de datos.

\paragraph{Alternativas.} Se consideran Redis, Valkey, Memcached y una vista materializada en PostgreSQL.
El \cref{tab:tec-ad55} presenta la matriz de decisión ponderada.
La alternativa elegida obtiene 4,70 y la segunda mejor, Valkey, obtiene 4,45.

\begingroup\small
\begin{longtable}{L{4.96cm} C{1.10cm} C{2.25cm} C{2.25cm} C{2.25cm} C{2.25cm}}
\cabecera{6}{tab:tec-ad55}{Matriz de decisión ponderada de AD-55}{\textbf{Criterio (driver)} & \textbf{Peso} & \hspace{0pt}\textbf{Redis 8.8} & \hspace{0pt}\textbf{Valkey} & \hspace{0pt}\textbf{Memcached} & \hspace{0pt}\textbf{Vista materializada}}
Lectura con TTL e invalidación por clave (QA-06, AD-17) & 30\,\% & 5 & 5 & 4 & 3 \\
Antigüedad del dato controlada (QA-06, AD-16) & 20\,\% & 4 & 4 & 4 & 2 \\
Componente separado de la base que escala aparte (QA-05, AD-18) & 15\,\% & 5 & 5 & 5 & 1 \\
Cliente maduro en la pila de AD-50 & 15\,\% & 5 & 4 & 3 & 5 \\
Experiencia del equipo & 10\,\% & 5 & 3 & 2 & 4 \\
Licencia y costo & 10\,\% & 4 & 5 & 5 & 5 \\
\midrule
\textbf{Total ponderado} & 100\,\% & \textbf{4,70} & 4,45 & 3,90 & 3,10 \\
\end{longtable}
\endgroup

\paragraph{Alternativas descartadas.}
\begin{itemize}
  \item Valkey gana en licencia, porque conserva la licencia BSD, pero pierde en experiencia del equipo y en madurez del cliente probado con la pila.
  \item Memcached pierde en invalidación por clave (QA-06), porque no ofrece estructuras ni invalidación por patrón.
  \item La vista materializada gana en cliente y licencia, pero pierde en componente separado (QA-05), porque concentra la lectura en la misma base.
\end{itemize}

\paragraph{Compromiso aceptado.} Redis 8 se distribuye con licencias RSALv2, SSPL y AGPLv3.
La AGPLv3 permite su uso en este proyecto académico.
Valkey es un sustituto compatible con el mismo protocolo si la licencia cambia.

\paragraph{Riesgo.} Riesgo R-14 del ATAM: la antigüedad del dato presentado depende de la carga.
En el perfil local se midió p95 de 1\,485\,ms con máximo de 11,8\,s.
Una corrida anterior dio p95 de 4,2\,s con 10 muestras.

\paragraph{Dónde se ve.} Vista física de la implementación: nodo \texttt{cache}.

\paragraph{Evidencia en el código.} QA-06 se midió con p95 de 142\,ms en el mapa de desastres y con la marca \texttt{datoAl} en el 100\,\% de las respuestas.
Archivos de evidencia:
{\raggedright
\begin{itemize}
  \item \texttt{proyectoparte2/\br{}servicios/\br{}consulta-danos/\br{}src/\br{}ConsultaDanos.Infraestructura/\br{}Cache.cs}
  \item \texttt{proyectoparte2/\br{}servicios/\br{}consulta/\br{}src/\br{}Consulta.Infraestructura/\br{}Adaptadores.cs}
\end{itemize}
\par}

%---------------------------------------------------------------------- AD-56
\subsection{AD-56. Almacén de archivos}
\label{sec:tec-ad56}

\paragraph{Contexto.} Esta decisión materializa el almacén de archivos de la vista física, donde se guardan las evidencias de Evaluación de Daños y los documentos de Reconstrucción.
El almacén debe mantener la evidencia disponible aunque falle una réplica del microservicio (AD-24, QA-02).
También debe conservar el 100\,\% de las fotos sincronizadas desde el cliente móvil (AD-31, QA-14).

\paragraph{Alternativas.} Se consideran MinIO, el disco local del contenedor, una columna binaria en PostgreSQL y SeaweedFS.
El \cref{tab:tec-ad56} presenta la matriz de decisión ponderada.
La alternativa elegida obtiene 4,20 y la segunda mejor, SeaweedFS, obtiene 4,05.

\begingroup\small
\begin{longtable}{L{4.96cm} C{1.10cm} C{2.25cm} C{2.25cm} C{2.25cm} C{2.25cm}}
\cabecera{6}{tab:tec-ad56}{Matriz de decisión ponderada de AD-56}{\textbf{Criterio (driver)} & \textbf{Peso} & \hspace{0pt}\textbf{MinIO} & \hspace{0pt}\textbf{Disco local} & \hspace{0pt}\textbf{Columna binaria} & \hspace{0pt}\textbf{SeaweedFS}}
Evidencia disponible ante la pérdida de una réplica (QA-02, AD-24) & 25\,\% & 4 & 1 & 4 & 4 \\
Integridad de la evidencia sincronizada (QA-14) & 20\,\% & 5 & 3 & 5 & 5 \\
Tamaño de la base y de su respaldo (QA-03, QA-06) & 15\,\% & 5 & 4 & 1 & 5 \\
Interfaz S3 estándar que permite sustituir el producto (QA-10) & 15\,\% & 5 & 1 & 1 & 4 \\
Experiencia del equipo & 10\,\% & 4 & 5 & 4 & 1 \\
Licencia y continuidad del producto & 15\,\% & 2 & 5 & 5 & 4 \\
\midrule
\textbf{Total ponderado} & 100\,\% & \textbf{4,20} & 2,85 & 3,45 & 4,05 \\
\end{longtable}
\endgroup

\paragraph{Alternativas descartadas.}
\begin{itemize}
  \item El disco local gana en experiencia y licencia, pero pierde en disponibilidad ante la pérdida de una réplica (QA-02), que AD-24 descarta.
  \item La columna binaria iguala en integridad (QA-14), pero pierde en tamaño de la base y de su respaldo (QA-03).
  \item SeaweedFS gana en continuidad del producto, pero pierde en experiencia del equipo y en compatibilidad S3 verificada.
\end{itemize}

\paragraph{Compromiso aceptado.} El proyecto original de MinIO dejó de publicar imágenes y el archivo de composición usa el fork comunitario \texttt{pgsty/minio} (DT-11).
El código depende solo de la API S3, lo que permite sustituir el producto sin modificar los microservicios.

\paragraph{Riesgo.} El fork comunitario puede dejar de mantenerse.
En el entorno de un computador \texttt{archivos} es un solo nodo, y su caída deja sin evidencias a las réplicas de \texttt{evaluacion-danos}.
No corresponde a un riesgo del ATAM.

\paragraph{Dónde se ve.} Vista física de la implementación: nodo \texttt{archivos}.

\paragraph{Evidencia en el código.} La corrida de QA-14 conservó 120 de 120 fotos a través del gateway.
Archivos de evidencia:
{\raggedright
\begin{itemize}
  \item \texttt{proyectoparte2/\br{}servicios/\br{}evaluacion-danos/\br{}src/\br{}EvaluacionDanos.Infraestructura/\br{}Archivos/\br{}AlmacenEvidenciasMinio.cs}
  \item \texttt{proyectoparte2/\br{}servicios/\br{}reconstruccion/\br{}src/\br{}Reconstruccion.Infraestructura/\br{}Archivos/\br{}AlmacenDocumentosMinio.cs}
\end{itemize}
\par}

%---------------------------------------------------------------------- AD-57
\subsection{AD-57. API Gateway}
\label{sec:tec-ad57}

\paragraph{Contexto.} Esta decisión materializa el punto único de entrada (AD-04), la priorización y degradación por tipo de cliente (AD-19) y el nivel 1 de la autorización (AD-22).
El gateway atiende QA-11, que exige el rechazo del 100\,\% de las modificaciones de un organismo no responsable.
También atiende QA-05 y QA-06, porque prioriza el tráfico crítico cuando la carga crece.

\paragraph{Alternativas.} Se consideran YARP, Kong, Traefik, Ocelot y nginx.
El \cref{tab:tec-ad57} presenta la matriz de decisión ponderada.
La alternativa elegida obtiene 4,85 y la segunda mejor, Kong, obtiene 3,45.

\begingroup\small
\begin{longtable}{L{3.93cm} C{1.10cm} C{1.95cm} C{1.95cm} C{1.95cm} C{1.95cm} C{1.95cm}}
\cabecera{7}{tab:tec-ad57}{Matriz de decisión ponderada de AD-57}{\textbf{Criterio (driver)} & \textbf{Peso} & \hspace{0pt}\textbf{YARP 2.3} & \hspace{0pt}\textbf{Kong} & \hspace{0pt}\textbf{Traefik} & \hspace{0pt}\textbf{Ocelot} & \hspace{0pt}\textbf{nginx}}
Priorización por clase de tráfico y degradación (QA-05, QA-06, AD-19) & 25\,\% & 5 & 3 & 2 & 3 & 2 \\
Validación JWT y regla de dominio por ruta (QA-11, AD-22) & 25\,\% & 5 & 4 & 3 & 4 & 2 \\
Catálogo de rutas validado al arrancar (QA-11) & 15\,\% & 5 & 4 & 3 & 3 & 2 \\
Latencia agregada por el intermediario (QA-06) & 15\,\% & 4 & 4 & 4 & 3 & 5 \\
Experiencia del equipo en la pila de AD-50 & 10\,\% & 5 & 2 & 3 & 4 & 3 \\
Licencia y mantenimiento & 10\,\% & 5 & 3 & 5 & 3 & 5 \\
\midrule
\textbf{Total ponderado} & 100\,\% & \textbf{4,85} & 3,45 & 3,10 & 3,35 & 2,85 \\
\end{longtable}
\endgroup

\paragraph{Alternativas descartadas.}
\begin{itemize}
  \item Kong pierde en priorización con degradación (QA-05), que exigiría complementos propios en otro lenguaje.
  \item Traefik pierde en priorización y en validación de la regla de dominio por ruta (QA-11).
  \item Ocelot pierde en catálogo validado al arrancar (QA-11) y en mantenimiento del proyecto.
  \item nginx gana en latencia agregada (QA-06), pero pierde en validación JWT con regla de dominio (QA-11), que no ofrece sin módulos adicionales.
\end{itemize}

\paragraph{Compromiso aceptado.} La priorización y la regla de dominio son código propio en el gateway.
Ese código se prueba con la misma suite de .NET, pero lo mantiene el equipo.

\paragraph{Riesgo.} Riesgo R-12 del ATAM: el gateway recibe tráfico sin TLS en el entorno de laboratorio (L-07).
Riesgo R-13 del ATAM: la validación depende de las llaves y del emisor de Keycloak (L-08).
Los rechazos del gateway no se registran en Auditoría (L-02).

\paragraph{Dónde se ve.} Vista lógica: componente API Gateway.
Vista física de la implementación: nodo \texttt{gateway} con dos réplicas y el artefacto \texttt{rutas.json}.

\paragraph{Evidencia en el código.} QA-11 se verificó con 218 de 218 intentos rechazados (100\,\%), 0 falsos negativos y 0 datos modificados sobre 12 rutas con regla de dominio.
Archivos de evidencia:
{\raggedright
\begin{itemize}
  \item \texttt{proyectoparte2/\br{}gateway/\br{}src/\br{}Plataforma.Gateway/\br{}rutas.json}
  \item \texttt{proyectoparte2/\br{}gateway/\br{}src/\br{}Plataforma.Gateway/\br{}Prioridad/\br{}Priorizacion.cs}
  \item \texttt{proyectoparte2/\br{}gateway/\br{}src/\br{}Plataforma.Gateway/\br{}Seguridad/\br{}AutorizacionNivel1.cs}
\end{itemize}
\par}

%---------------------------------------------------------------------- AD-58
\subsection{AD-58. Proveedor de identidad y clientes OIDC}
\label{sec:tec-ad58}

\paragraph{Contexto.} Esta decisión materializa el proveedor de identidad federado con OIDC y PKCE (AD-49) y la validación del token en el gateway y en cada microservicio (AD-04, AD-22).
Sus clientes son la consola, el cliente móvil y el tablero de daños. El tablero usa el cliente confidencial \texttt{consulta-\br{}danos-\br{}tablero} y una sesión con cookie que valida la propia página, porque su ruta es pública en el gateway (E-41).
El proveedor atiende QA-11, porque emite los roles y el atributo de organismo que usa la regla de dominio.
También atiende QA-14, porque el refresco del token sin conexión permite trabajar días en campo y sincronizar al volver.

\paragraph{Alternativas.} Se consideran Keycloak, Auth0, Duende IdentityServer y una identidad propia.
El \cref{tab:tec-ad58} presenta la matriz de decisión ponderada.
La alternativa elegida obtiene 4,90 y la segunda mejor, Duende IdentityServer, obtiene 3,75.

\begingroup\small
\begin{longtable}{L{4.96cm} C{1.10cm} C{2.25cm} C{2.25cm} C{2.25cm} C{2.25cm}}
\cabecera{6}{tab:tec-ad58}{Matriz de decisión ponderada de AD-58}{\textbf{Criterio (driver)} & \textbf{Peso} & \hspace{0pt}\textbf{Keycloak 26} & \hspace{0pt}\textbf{Auth0} & \hspace{0pt}\textbf{Duende IdentityServer} & \hspace{0pt}\textbf{Identidad propia}}
OIDC con PKCE para la consola y el móvil, con refresco sin conexión (QA-11, QA-14, AD-49) & 25\,\% & 5 & 5 & 4 & 2 \\
Roles y organismo como \emph{claims} para la regla de dominio (QA-11, AD-22) & 20\,\% & 5 & 4 & 4 & 3 \\
Despliegue dentro de la plataforma, sin servicio externo en la emergencia (QA-02, CON-3) & 15\,\% & 5 & 2 & 5 & 5 \\
Configuración declarada en un archivo versionado (QA-10) & 10\,\% & 5 & 3 & 4 & 3 \\
Experiencia del equipo & 10\,\% & 4 & 2 & 2 & 3 \\
Licencia y costo & 20\,\% & 5 & 2 & 3 & 5 \\
\midrule
\textbf{Total ponderado} & 100\,\% & \textbf{4,90} & 3,25 & 3,75 & 3,45 \\
\end{longtable}
\endgroup

\paragraph{Alternativas descartadas.}
\begin{itemize}
  \item Auth0 iguala en OIDC con PKCE (QA-11), pero pierde en despliegue dentro de la plataforma (CON-3) y en costo por usuario activo.
  \item Duende IdentityServer pierde en licencia, porque exige licencia comercial por encima de un umbral de ingresos, y en experiencia del equipo.
  \item La identidad propia gana en licencia y despliegue, pero pierde en OIDC con PKCE (QA-11), porque obligaría a construir la seguridad desde cero.
\end{itemize}

\paragraph{Compromiso aceptado.} El token se valida dos veces, en el gateway y en el microservicio, con un validador propio en \texttt{Plataforma.Comun.Seguridad}.
La doble validación agrega trabajo por petición a cambio de la defensa en profundidad de AD-22.

\paragraph{Riesgo.} Riesgo R-13 del ATAM: el emisor del token es \texttt{http://localhost:8180} y debe coincidir en el gateway y en los microservicios (L-08).
La recarga de las llaves JWT está pendiente (L-09).
Riesgo R-12 del ATAM: Keycloak corre en \texttt{start-dev} sin TLS (L-07).

\paragraph{Dónde se ve.} Vista lógica: componente externo Proveedor de identidad.
Vista física de la implementación: nodo \texttt{keycloak}.
Vista de implementación móvil: paquete \texttt{infraestructura/identidad}.

\paragraph{Evidencia en el código.} Siete pruebas instrumentadas de la app corren contra el gateway y Keycloak, incluido el rechazo por rol.
Archivos de evidencia:
{\raggedright
\begin{itemize}
  \item \texttt{proyectoparte2/\br{}despliegue/\br{}keycloak/\br{}realm-plataforma.json}
  \item \texttt{proyectoparte2/\br{}comun/\br{}Plataforma.Comun.Seguridad/\br{}IdentidadPorJwt.cs}
  \item \texttt{danos/\br{}app/\br{}src/\br{}main/\br{}java/\br{}com/\br{}example/\br{}danos/\br{}infraestructura/\br{}identidad/\br{}ConexionOidcBuilder.kt}
\end{itemize}
\par}

%---------------------------------------------------------------------- AD-59
\subsection{AD-59. Consola de coordinación}
\label{sec:tec-ad59}

\paragraph{Contexto.} Esta decisión materializa la capa de presentación web (AD-08) de la Consola de coordinación, con los módulos \texttt{mapa}, \texttt{evaluaciones} y \texttt{reconstruccion}.
La consola entra por el gateway (AD-04) y atiende QA-06 en la carga del mapa, QA-11 en el inicio de sesión y QA-10 en la incorporación de módulos.

\paragraph{Alternativas.} Se consideran Angular, React 19, Vue 3 y Blazor.
El \cref{tab:tec-ad59} presenta la matriz de decisión ponderada.
La alternativa elegida obtiene 4,80 y la segunda mejor, Vue 3, obtiene 4,10.

\begingroup\small
\begin{longtable}{L{4.96cm} C{1.10cm} C{2.25cm} C{2.25cm} C{2.25cm} C{2.25cm}}
\cabecera{6}{tab:tec-ad59}{Matriz de decisión ponderada de AD-59}{\textbf{Criterio (driver)} & \textbf{Peso} & \hspace{0pt}\textbf{Angular 22} & \hspace{0pt}\textbf{React 19} & \hspace{0pt}\textbf{Vue 3} & \hspace{0pt}\textbf{Blazor}}
Módulos con rutas, inyección y guardas (QA-10, AD-08) & 25\,\% & 5 & 3 & 4 & 4 \\
Cliente OIDC con PKCE certificado (QA-11, AD-49) & 20\,\% & 5 & 4 & 4 & 4 \\
Carga inicial y respuesta de la interfaz (QA-06) & 20\,\% & 4 & 4 & 5 & 2 \\
Archivos estáticos servidos detrás del gateway (QA-02, AD-04) & 10\,\% & 5 & 5 & 5 & 4 \\
Experiencia del equipo & 15\,\% & 5 & 3 & 2 & 3 \\
Licencia y costo & 10\,\% & 5 & 5 & 5 & 5 \\
\midrule
\textbf{Total ponderado} & 100\,\% & \textbf{4,80} & 3,80 & 4,10 & 3,55 \\
\end{longtable}
\endgroup

\paragraph{Alternativas descartadas.}
\begin{itemize}
  \item React 19 pierde en estructura por módulos con inyección y guardas (QA-10), que exige elegir y combinar bibliotecas.
  \item Vue 3 gana en carga inicial (QA-06), pero pierde en experiencia del equipo y en estructura por módulos.
  \item Blazor pierde en carga inicial (QA-06), por el tamaño del entorno WebAssembly que descarga el navegador.
\end{itemize}

\paragraph{Compromiso aceptado.} Angular produce un paquete inicial mayor que Vue.
La carga se mitiga con módulos cargados bajo demanda por ruta.

\paragraph{Riesgo.} Angular publica una versión mayor cada seis meses y obliga a actualizaciones frecuentes.
No corresponde a un riesgo del ATAM.

\paragraph{Dónde se ve.} Vista física de la implementación: nodo \texttt{web}.
Vista lógica: Consola de coordinación con los módulos \texttt{mapa}, \texttt{evaluaciones} y \texttt{reconstruccion}.

\paragraph{Evidencia en el código.} El \texttt{Dockerfile} compila con \texttt{node:22-alpine} y sirve los archivos con \texttt{nginx:1.29-alpine}.
Archivos de evidencia:
{\raggedright
\begin{itemize}
  \item \texttt{proyectoparte2/\br{}consola/\br{}package.json}
  \item \texttt{proyectoparte2/\br{}consola/\br{}Dockerfile}
  \item \texttt{proyectoparte2/\br{}consola/\br{}nginx.conf}
  \item \texttt{proyectoparte2/\br{}consola/\br{}src/\br{}app/\br{}modulos}
\end{itemize}
\par}

%---------------------------------------------------------------------- AD-60
\subsection{AD-60. Biblioteca de mapas}
\label{sec:tec-ad60}

\paragraph{Contexto.} Esta decisión materializa la presentación de las geometrías precalculadas (AD-17) y de su marca de tiempo (AD-16) en el módulo \texttt{mapa} de la consola.
La biblioteca atiende QA-06, porque el mapa es la consulta frecuente del coordinador.

\paragraph{Alternativas.} Se consideran Leaflet, OpenLayers, MapLibre GL y Google Maps.
El \cref{tab:tec-ad60} presenta la matriz de decisión ponderada.
La alternativa elegida obtiene 4,80 y la segunda mejor, OpenLayers, obtiene 3,90.

\begingroup\small
\begin{longtable}{L{4.96cm} C{1.10cm} C{2.25cm} C{2.25cm} C{2.25cm} C{2.25cm}}
\cabecera{6}{tab:tec-ad60}{Matriz de decisión ponderada de AD-60}{\textbf{Criterio (driver)} & \textbf{Peso} & \hspace{0pt}\textbf{Leaflet 1.9} & \hspace{0pt}\textbf{OpenLayers} & \hspace{0pt}\textbf{MapLibre GL} & \hspace{0pt}\textbf{Google Maps}}
Capas GeoJSON agregadas por zona (QA-06, AD-17) & 25\,\% & 5 & 5 & 5 & 4 \\
Peso de la biblioteca y tiempo de carga (QA-06) & 20\,\% & 5 & 3 & 2 & 3 \\
Independencia de un proveedor y de una llave (QA-02, CON-3) & 20\,\% & 4 & 4 & 4 & 1 \\
Integración con Angular y con el mapa del prototipo (QA-10) & 10\,\% & 5 & 3 & 3 & 3 \\
Experiencia del equipo & 10\,\% & 5 & 2 & 2 & 3 \\
Licencia y costo & 15\,\% & 5 & 5 & 5 & 2 \\
\midrule
\textbf{Total ponderado} & 100\,\% & \textbf{4,80} & 3,90 & 3,70 & 2,70 \\
\end{longtable}
\endgroup

\paragraph{Alternativas descartadas.}
\begin{itemize}
  \item OpenLayers iguala en capas GeoJSON (QA-06), pero pierde en peso de la biblioteca y en experiencia del equipo.
  \item MapLibre GL iguala en capas GeoJSON, pero pierde en peso de la biblioteca (QA-06), porque su ventaja con millones de puntos no hace falta con zonas agregadas.
  \item Google Maps pierde en independencia de un proveedor (CON-3) y en costo por carga de mapa.
\end{itemize}

\paragraph{Compromiso aceptado.} Leaflet dibuja en el navegador sin aceleración WebGL.
El volumen se controla con la agregación por zona del servidor (AD-17).

\paragraph{Riesgo.} Las teselas de OpenStreetMap no tienen acuerdo de servicio y su política de uso prohíbe el uso intensivo.
No corresponde a un riesgo del ATAM.

\paragraph{Dónde se ve.} Vista física de la implementación: nodo \texttt{web}.
Vista lógica: módulo \texttt{mapa} de la Consola de coordinación.

\paragraph{Evidencia en el código.} QA-06 se midió con p95 de 1\,628\,ms y p99 de 1\,935\,ms en el mapa de daños en el perfil local (L-04).
Archivos de evidencia:
{\raggedright
\begin{itemize}
  \item \texttt{proyectoparte2/\br{}consola/\br{}src/\br{}app/\br{}modulos/\br{}mapa/\br{}lienzo-mapa.ts}
  \item \texttt{proyectoparte2/\br{}consola/\br{}package.json}
\end{itemize}
\par}

%---------------------------------------------------------------------- AD-61
\subsection{AD-61. Plataforma del cliente móvil}
\label{sec:tec-ad61}

\paragraph{Contexto.} Esta decisión materializa la réplica local del cliente móvil (AD-31), la cola local de operaciones (AD-44) y el registro del desastre en campo (AD-48).
La plataforma atiende QA-14, cuya medida es 0 evaluaciones perdidas o sobrescritas en silencio al sincronizar 120 evaluaciones, y cuya respuesta exige 0 duplicados.

\paragraph{Alternativas.} Se consideran Android nativo con Kotlin, Flutter, React Native y Kotlin Multiplatform.
El \cref{tab:tec-ad61} presenta la matriz de decisión ponderada.
La alternativa elegida obtiene 4,60 y la segunda mejor, Kotlin Multiplatform, obtiene 4,00.

\begingroup\small
\begin{longtable}{L{4.96cm} C{1.10cm} C{2.25cm} C{2.25cm} C{2.25cm} C{2.25cm}}
\cabecera{6}{tab:tec-ad61}{Matriz de decisión ponderada de AD-61}{\textbf{Criterio (driver)} & \textbf{Peso} & \hspace{0pt}\textbf{Kotlin nativo} & \hspace{0pt}\textbf{Flutter} & \hspace{0pt}\textbf{React Native} & \hspace{0pt}\textbf{Kotlin Multiplatform}}
Trabajo en segundo plano y almacenamiento del sistema (QA-14, AD-44) & 25\,\% & 5 & 3 & 3 & 4 \\
Base local con transacciones y migraciones (QA-14, AD-31) & 20\,\% & 5 & 4 & 3 & 4 \\
Cámara y ubicación para las evidencias (QA-14) & 10\,\% & 5 & 4 & 4 & 4 \\
Reutilización del prototipo de FP-06 y de sus pruebas (QA-14) & 15\,\% & 5 & 1 & 1 & 4 \\
Portabilidad a otro sistema operativo móvil (QA-10) & 10\,\% & 1 & 5 & 5 & 4 \\
Experiencia del equipo & 10\,\% & 5 & 2 & 3 & 3 \\
Licencia y costo & 10\,\% & 5 & 5 & 5 & 5 \\
\midrule
\textbf{Total ponderado} & 100\,\% & \textbf{4,60} & 3,30 & 3,20 & 4,00 \\
\end{longtable}
\endgroup

\paragraph{Alternativas descartadas.}
\begin{itemize}
  \item Flutter gana en portabilidad (QA-10), pero pierde en reutilización del prototipo de FP-06 (QA-14), que habría que reescribir.
  \item React Native gana en portabilidad, pero pierde en base local transaccional (QA-14) y en reutilización del prototipo.
  \item Kotlin Multiplatform pierde en trabajo en segundo plano (QA-14), porque esa parte sigue siendo específica de cada plataforma.
\end{itemize}

\paragraph{Compromiso aceptado.} El cliente funciona solo en Android.
Un cliente para iOS exigiría otra implementación de la capa de presentación y de la sincronización.

\paragraph{Riesgo.} La app de entrega es la variante \texttt{debug}, que admite \texttt{http://} (L-07).
La evidencia solo admite fotos y no guarda un autor propio (L-10).

\paragraph{Dónde se ve.} Vista de implementación móvil: paquete \texttt{presentacion}.
Vista física de la implementación: \texttt{<<device>>} Teléfono Android con \texttt{DanosCampo.apk}.

\paragraph{Evidencia en el código.} La app \texttt{co.plataforma.danos} 2.1.0 tiene 81 pruebas unitarias y 16 instrumentadas en verde.
Archivos de evidencia:
{\raggedright
\begin{itemize}
  \item \texttt{danos/\br{}gradle/\br{}libs.versions.toml}
  \item \texttt{danos/\br{}app/\br{}build.gradle.kts}
\end{itemize}
\par}

%---------------------------------------------------------------------- AD-62
\subsection{AD-62. Base de datos local del cliente móvil}
\label{sec:tec-ad62}

\paragraph{Contexto.} Esta decisión materializa la réplica local (AD-31), la cola local de operaciones inmutables con copia completa del contenido (AD-44) y la tabla \texttt{registro\_desastre} del registro en campo (AD-48).
La base atiende QA-14, porque la escritura y su operación en la cola deben confirmarse en la misma transacción.

\paragraph{Alternativas.} Se consideran Room, Realm, SQLDelight y DataStore.
El \cref{tab:tec-ad62} presenta la matriz de decisión ponderada.
La alternativa elegida obtiene 4,90 y la segunda mejor, SQLDelight, obtiene 4,35.

\begingroup\small
\begin{longtable}{L{4.96cm} C{1.10cm} C{2.25cm} C{2.25cm} C{2.25cm} C{2.25cm}}
\cabecera{6}{tab:tec-ad62}{Matriz de decisión ponderada de AD-62}{\textbf{Criterio (driver)} & \textbf{Peso} & \hspace{0pt}\textbf{Room 2.8} & \hspace{0pt}\textbf{Realm} & \hspace{0pt}\textbf{SQLDelight} & \hspace{0pt}\textbf{DataStore}}
Transacción atómica entre la escritura y la cola local (QA-14, AD-44) & 30\,\% & 5 & 4 & 5 & 1 \\
Migraciones explícitas y verificables (QA-10, QA-14) & 20\,\% & 5 & 3 & 4 & 2 \\
Integración con Compose y corrutinas (AD-61) & 15\,\% & 5 & 4 & 4 & 4 \\
Continuidad del producto & 15\,\% & 5 & 1 & 5 & 5 \\
Experiencia del equipo & 10\,\% & 5 & 2 & 2 & 4 \\
Control del SQL verificado en compilación (QA-14) & 10\,\% & 4 & 2 & 5 & 1 \\
\midrule
\textbf{Total ponderado} & 100\,\% & \textbf{4,90} & 2,95 & 4,35 & 2,55 \\
\end{longtable}
\endgroup

\paragraph{Alternativas descartadas.}
\begin{itemize}
  \item Realm pierde en continuidad del producto, porque su fabricante retiró el SDK de dispositivos.
  \item SQLDelight gana en control del SQL verificado en compilación, pero pierde en migraciones con esquema exportado y en experiencia del equipo.
  \item DataStore pierde en transacción atómica con la cola (QA-14), porque guarda pares clave y valor sin relaciones.
\end{itemize}

\paragraph{Compromiso aceptado.} Cada cambio de esquema exige una migración escrita a mano y su esquema exportado.
Ese costo es el precio de no perder datos locales al actualizar la app (QA-14).

\paragraph{Riesgo.} Una migración defectuosa puede impedir abrir la base local con operaciones pendientes.
No corresponde a un riesgo del ATAM.

\paragraph{Dónde se ve.} Vista de implementación móvil: paquete \texttt{infraestructura/local}.
Vista física de la implementación: artefacto \texttt{danos.db} en el Teléfono Android.

\paragraph{Evidencia en el código.} \texttt{AppDatabase} declara la versión 3 con \texttt{MIGRACION\_1\_2} y \texttt{MIGRACION\_2\_3}, y la carpeta de esquemas guarda las versiones 1, 2 y 3.
Archivos de evidencia:
{\raggedright
\begin{itemize}
  \item \texttt{danos/\br{}app/\br{}src/\br{}main/\br{}java/\br{}com/\br{}example/\br{}danos/\br{}infraestructura/\br{}local/\br{}AppDatabase.kt}
  \item \texttt{danos/\br{}app/\br{}schemas}
\end{itemize}
\par}

%---------------------------------------------------------------------- AD-63
\subsection{AD-63. Sincronización del cliente móvil}
\label{sec:tec-ad63}

\paragraph{Contexto.} Esta decisión materializa la sincronización diferida (AD-31), el reenvío de la cola local (AD-44) y el envío idempotente del registro del desastre por el identificador del cliente (AD-48).
La sincronización atiende QA-14, cuya respuesta exige reenvío sin duplicados y 120 evaluaciones en $\le$ 10\,min.
También atiende QA-11, porque todo envío pasa por el gateway con el token OIDC.

\paragraph{Alternativas.} Se consideran WorkManager, SyncAdapter, un servicio en primer plano y Firebase.
El \cref{tab:tec-ad63} presenta la matriz de decisión ponderada.
La alternativa elegida obtiene 4,85 y la segunda mejor, SyncAdapter, obtiene 4,10.

\begingroup\small
\begin{longtable}{L{4.96cm} C{1.10cm} C{2.25cm} C{2.25cm} C{2.25cm} C{2.25cm}}
\cabecera{6}{tab:tec-ad63}{Matriz de decisión ponderada de AD-63}{\textbf{Criterio (driver)} & \textbf{Peso} & \hspace{0pt}\textbf{WorkManager} & \hspace{0pt}\textbf{SyncAdapter} & \hspace{0pt}\textbf{Servicio en primer plano} & \hspace{0pt}\textbf{Firebase}}
Ejecución garantizada con red disponible y reintento tras reinicio (QA-14, AD-44) & 30\,\% & 5 & 4 & 2 & 4 \\
Protocolo propio idempotente por identificador del cliente (QA-14, AD-48) & 25\,\% & 5 & 5 & 5 & 1 \\
Inmediatez del envío al recuperar la red (QA-14) & 15\,\% & 4 & 3 & 5 & 4 \\
Envío por el gateway con token OIDC (QA-11, AD-04) & 10\,\% & 5 & 5 & 5 & 1 \\
Experiencia del equipo & 10\,\% & 5 & 2 & 3 & 3 \\
Licencia y dependencia de un servicio externo & 10\,\% & 5 & 5 & 5 & 2 \\
\midrule
\textbf{Total ponderado} & 100\,\% & \textbf{4,85} & 4,10 & 3,90 & 2,65 \\
\end{longtable}
\endgroup

\paragraph{Alternativas descartadas.}
\begin{itemize}
  \item SyncAdapter pierde en ejecución garantizada (QA-14) y en experiencia del equipo, porque exige una cuenta del sistema y un proveedor de contenido.
  \item El servicio en primer plano gana en inmediatez del envío, pero pierde en ejecución garantizada tras un reinicio (QA-14).
  \item Firebase pierde en protocolo idempotente propio (QA-14) y en envío por el gateway (QA-11), porque sincroniza contra su propio servicio.
\end{itemize}

\paragraph{Compromiso aceptado.} El sistema operativo decide el instante de la ejecución de WorkManager.
El envío puede demorar unos segundos después de recuperar la red.

\paragraph{Riesgo.} Riesgo R-12 del ATAM: la variante \texttt{debug} admite tráfico sin TLS (L-07).
Riesgo R-13 del ATAM: desde un teléfono físico el emisor del token exige cambiar la configuración de Keycloak (L-08).

\paragraph{Dónde se ve.} Vista de implementación móvil: paquete \texttt{infraestructura/sync}.
Vista física de la implementación: \texttt{DanosCampo.apk} en el Teléfono Android.

\paragraph{Evidencia en el código.} QA-14 se verificó con 115 evaluaciones aplicadas, 5 en conflicto, 0 pérdidas, 0 sobrescrituras y 0 duplicados en 27,8\,s.
Archivos de evidencia:
{\raggedright
\begin{itemize}
  \item \texttt{danos/\br{}app/\br{}src/\br{}main/\br{}java/\br{}com/\br{}example/\br{}danos/\br{}infraestructura/\br{}sync/\br{}SincronizacionWorker.kt}
  \item \texttt{danos/\br{}app/\br{}src/\br{}main/\br{}java/\br{}com/\br{}example/\br{}danos/\br{}infraestructura/\br{}sync/\br{}EnviadorDesastres.kt}
\end{itemize}
\par}

%---------------------------------------------------------------------- AD-64
\subsection{AD-64. Biblioteca de resiliencia}
\label{sec:tec-ad64}

\paragraph{Contexto.} Esta decisión materializa el tiempo de espera, el reintento con espera creciente y el cortacircuitos de las fronteras (AD-05) y el cortacircuitos por adaptador (AD-37).
La biblioteca atiende QA-04, cuya respuesta exige cortar la propagación del fallo en $\le$ 60\,s.
También atiende QA-01, porque la caída de un sistema externo no debe degradar las demás funciones.

\paragraph{Alternativas.} Se consideran Microsoft.Extensions.Http.Resilience, una implementación propia y Steeltoe.
El \cref{tab:tec-ad64} presenta la matriz de decisión ponderada.
La alternativa elegida obtiene 4,90 y la segunda mejor, Implementación propia, obtiene 3,15.

\begingroup\small
\begin{longtable}{L{6.74cm} C{1.10cm} C{2.50cm} C{2.50cm} C{2.50cm}}
\cabecera{5}{tab:tec-ad64}{Matriz de decisión ponderada de AD-64}{\textbf{Criterio (driver)} & \textbf{Peso} & \hspace{0pt}\textbf{MS Resilience (Polly v8)} & \hspace{0pt}\textbf{Implementación propia} & \hspace{0pt}\textbf{Steeltoe}}
Tiempo de espera, reintento con variación aleatoria y cortacircuitos configurables (QA-01, AD-05) & 30\,\% & 5 & 3 & 2 \\
Corte de la propagación en $\le$ 60\,s verificable (QA-04) & 20\,\% & 5 & 3 & 3 \\
Integración con la fábrica de clientes HTTP de AD-50 & 15\,\% & 5 & 3 & 4 \\
Estado del circuito seguro ante concurrencia (QA-04) & 15\,\% & 5 & 2 & 4 \\
Experiencia del equipo & 10\,\% & 4 & 4 & 1 \\
Licencia y mantenimiento & 10\,\% & 5 & 5 & 4 \\
\midrule
\textbf{Total ponderado} & 100\,\% & \textbf{4,90} & 3,15 & 2,90 \\
\end{longtable}
\endgroup

\paragraph{Alternativas descartadas.}
\begin{itemize}
  \item La implementación propia iguala en experiencia y licencia, pero pierde en estado del circuito seguro ante concurrencia (QA-04).
  \item Steeltoe pierde en políticas configurables (QA-01) y en experiencia del equipo.
\end{itemize}

\paragraph{Compromiso aceptado.} Polly mide una proporción de fallos en una ventana, no un número de fallos consecutivos.
La regla de AD-05 se expresa como 3 llamadas mínimas con 100\,\% de fallos en 30\,s.

\paragraph{Riesgo.} Un reintento sobre una operación no idempotente duplicaría el efecto.
\texttt{Resiliencia.cs} reintenta solo operaciones idempotentes.
No corresponde a un riesgo del ATAM.

\paragraph{Dónde se ve.} Vista de implementación: proyecto \texttt{Plataforma.Comun.Adaptadores} (\texttt{Resiliencia.cs}).
Modelo de integración: cortacircuitos de cada adaptador.

\paragraph{Evidencia en el código.} La configuración usa 3\,s por intento, 2 reintentos, ventana de 30\,s y apertura de 15\,s. Las 12 pruebas de \texttt{AdaptadorSe6.Integracion} y las 14 de \texttt{AdaptadorSe7.Integracion} verifican el cortacircuitos y el reintento sin pérdida (QA-01).
Archivos de evidencia:
{\raggedright
\begin{itemize}
  \item \texttt{proyectoparte2/\br{}comun/\br{}Plataforma.Comun.Adaptadores/\br{}Resiliencia.cs}
  \item \texttt{proyectoparte2/\br{}Directory.Packages.props}
\end{itemize}
\par}

%---------------------------------------------------------------------- AD-65
\subsection{AD-65. Empaquetado, ejecución y escalado}
\label{sec:tec-ad65}

\paragraph{Contexto.} Esta decisión materializa las réplicas adicionales del camino de lectura (AD-18), la réplica activa con sondeo de salud (AD-24) y la ejecución híbrida con adaptadores efímeros (AD-25).
El empaquetado atiende QA-05, que exige escalar solo los componentes bajo presión sin modificar código.
También atiende QA-02 por el reemplazo de réplicas y QA-01 por el escalado de los adaptadores según su cola.

\paragraph{Alternativas.} Se consideran contenedores con Kubernetes, máquinas virtuales, Docker Swarm, Nomad y Knative.
El \cref{tab:tec-ad65} presenta la matriz de decisión ponderada.
La alternativa elegida obtiene 4,60 y la segunda mejor, Knative, obtiene 3,85.

\begingroup\small
\begin{longtable}{L{3.93cm} C{1.10cm} C{1.95cm} C{1.95cm} C{1.95cm} C{1.95cm} C{1.95cm}}
\cabecera{7}{tab:tec-ad65}{Matriz de decisión ponderada de AD-65}{\textbf{Criterio (driver)} & \textbf{Peso} & \hspace{0pt}\textbf{Docker + K8s (HPA, KEDA)} & \hspace{0pt}\textbf{Máquinas virtuales} & \hspace{0pt}\textbf{Docker Swarm} & \hspace{0pt}\textbf{Nomad} & \hspace{0pt}\textbf{Knative}}
Escalado automático selectivo sin modificar código (QA-05, AD-18) & 25\,\% & 5 & 2 & 2 & 4 & 5 \\
Escalado de 0 a N por largo de cola (QA-01, AD-25) & 15\,\% & 5 & 1 & 1 & 3 & 3 \\
Sondeo de salud y reemplazo de réplicas (QA-02, AD-24) & 20\,\% & 5 & 3 & 4 & 4 & 5 \\
Misma imagen en el computador y en el clúster (QA-09) & 15\,\% & 5 & 2 & 5 & 4 & 4 \\
Experiencia del equipo & 10\,\% & 4 & 3 & 3 & 1 & 1 \\
Costo de operación y licencia & 15\,\% & 3 & 2 & 5 & 3 & 3 \\
\midrule
\textbf{Total ponderado} & 100\,\% & \textbf{4,60} & 2,15 & 3,25 & 3,40 & 3,85 \\
\end{longtable}
\endgroup

\paragraph{Alternativas descartadas.}
\begin{itemize}
  \item Las máquinas virtuales pierden en escalado selectivo (QA-05) y en escalado por largo de cola (QA-01).
  \item Docker Swarm gana en costo de operación, pero pierde en escalado automático (QA-05), que no ofrece.
  \item Nomad pierde en escalado por cola y en licencia, porque cambió a BSL en 2023.
  \item Knative iguala en escalado selectivo y sondeo de salud, pero pierde en escalado por largo de cola (QA-01), porque está orientado a HTTP.
\end{itemize}

\paragraph{Compromiso aceptado.} En Docker Compose los adaptadores quedan siempre encendidos y no hay escalado automático.
El escalado de 0 a N solo existe en los manifiestos de Kubernetes (AD-25).

\paragraph{Riesgo.} El despliegue en Kubernetes no se verificó de punta a punta (L-05).
El escalado de 0 a N con KEDA no se midió (QA-01) y QA-05 no se midió con 10\,000 usuarios (L-03, L-04).

\paragraph{Dónde se ve.} Vista física de la implementación: nodo Docker Engine con los servicios del archivo de composición y anotación de Kubernetes con HPA y KEDA.

\paragraph{Evidencia en el código.} Los manifiestos declaran HPA de 2 a 6 réplicas para \texttt{consulta} y \texttt{consulta-\br{}danos}, de 2 a 5 para \texttt{consulta-\br{}programa}, y \texttt{ScaledObject} de KEDA de 0 a 4 réplicas para la salida de SE-6.
Archivos de evidencia:
{\raggedright
\begin{itemize}
  \item \texttt{proyectoparte2/\br{}despliegue/\br{}compose/\br{}docker-compose.yml}
  \item \texttt{proyectoparte2/\br{}despliegue/\br{}k8s/\br{}base/\br{}20-servicios.yaml}
  \item \texttt{proyectoparte2/\br{}despliegue/\br{}k8s/\br{}base/\br{}adaptador-se6.yaml}
\end{itemize}
\par}

%---------------------------------------------------------------------- AD-66
\subsection{AD-66. Plataforma del adaptador de SE-8}
\label{sec:tec-ad66}

\paragraph{Contexto.} CON-9, agregada el 19 de noviembre de 2026, exige que la solución final incluya Jakarta EE y .NET (E-44).
AD-03 ubica la heterogeneidad tecnológica (CON-2) en los adaptadores, no en los microservicios: un adaptador es el único punto de la arquitectura donde una pila distinta no toca el núcleo.
El adaptador de SE-8 estaba diseñado y no implementado (L-06).
QA-08 le exige una instancia encendida (AD-25), estado de entrega por alerta y \enquote{entregada} solo con acuse.
QA-09 pide que incorporar un sistema externo modifique 0 componentes existentes.
Esta decisión elige la plataforma de ese adaptador.

\paragraph{Alternativas.} Se consideran Jakarta EE 11 sobre WildFly 42.0.0.Final (imagen \texttt{quay.io/\br{}wildfly/\br{}wildfly:\br{}42.0.0.Final-\br{}jdk21}, Java 21), Spring Boot 3 sobre Java 21 y .NET 10 con \texttt{Plataforma.\br{}Comun.\br{}Adaptadores}, como los adaptadores de SE-6 y SE-7.
El \cref{tab:tec-ad66} presenta la matriz de decisión ponderada.
La alternativa elegida obtiene 4,35 y la segunda mejor, .NET 10, obtiene 4,00.

\begingroup\small
\begin{longtable}{L{6.0cm} C{1.10cm} C{2.6cm} C{2.6cm} C{2.6cm}}
\cabecera{5}{tab:tec-ad66}{Matriz de decisión ponderada de AD-66}{\textbf{Criterio (driver)} & \textbf{Peso} & \hspace{0pt}\textbf{Jakarta EE 11 (WildFly 42)} & \hspace{0pt}\textbf{Spring Boot 3} & \hspace{0pt}\textbf{.NET 10}}
Cumplimiento de CON-9 & 25\,\% & 5 & 2 & 1 \\
Adaptador nuevo sin modificar componentes existentes (QA-09, AD-03) & 20\,\% & 5 & 5 & 5 \\
Cortacircuitos del Cuadro 10 con aviso de apertura para EV-23 (QA-01, QA-08, AD-05, AD-37) & 15\,\% & 3 & 4 & 5 \\
Consumo transaccional: ack después del commit y DLQ a la quinta entrega (AD-02, AD-45) & 15\,\% & 4 & 4 & 5 \\
Instancia siempre encendida: memoria y arranque (QA-08, AD-25) & 10\,\% & 4 & 4 & 5 \\
Experiencia del equipo (especificaciones vistas en el curso) & 10\,\% & 4 & 3 & 5 \\
Licencia y costo & 5\,\% & 5 & 5 & 5 \\
\midrule
\textbf{Total ponderado} & 100\,\% & \textbf{4,35} & 3,65 & 4,00 \\
\end{longtable}
\endgroup

\paragraph{Alternativas descartadas.}
\begin{itemize}
  \item .NET 10 gana en resiliencia, mensajería y arranque, porque reutiliza \texttt{Plataforma.\br{}Comun.\br{}Adaptadores} y \texttt{ConsumidorIdempotente}. Pierde en el criterio que origina la decisión: no cumple CON-9.
  \item Spring Boot 3 corre sobre Java y usa algunas API de Jakarta, pero no es la plataforma Jakarta EE: cumple CON-9 solo en parte. Tiene un cortacircuitos con ventana de tiempo (Resilience4j) y es una tercera pila de marco que el equipo no vio en el curso.
\end{itemize}

\paragraph{Compromiso aceptado.} El adaptador usa las especificaciones estándar: Jakarta Messaging, EJB (\texttt{@Stateless}, \texttt{@Singleton}, Timer Service), Jakarta Persistence, Jakarta REST, CDI, Jakarta Concurrency y MicroProfile Config.
Dos piezas son propias y no vienen del estándar:
\begin{itemize}
  \item \textbf{Cortacircuitos propio en lugar de MicroProfile Fault Tolerance.} El \texttt{@CircuitBreaker} de MicroProfile mide una ventana de N llamadas, no la ventana de \qty{30}{\second} del Cuadro 10, y no avisa cuando el circuito se abre o se cierra. Ese aviso hace falta para publicar EV-23 por la bandeja (E-15). La clase \texttt{Cortacircuitos} aplica la regla de Polly (al menos 3 llamadas en \qty{30}{\second} con \qty{100}{\percent} de fallo, \qty{15}{\second} abierto y una sola llamada de prueba) y es segura ante concurrencia. Tiene 9 pruebas unitarias propias, una de ellas con 200 hilos.
  \item \textbf{EJB \texttt{@Singleton @Startup} en lugar de MDB.} Un MDB necesita un resource adapter JCA para el proveedor de mensajería, y RabbitMQ no tiene uno mantenido. El consumo usa igual la API estándar de Jakarta Messaging 3.1 (\texttt{Connection}, \texttt{Session} transaccional, \texttt{MessageConsumer}) con el cliente JMS de RabbitMQ en modo AMQP sobre la cola existente. El hilo lo da un \texttt{ManagedExecutorService}. Con \texttt{nackOnRollback}, \texttt{rollback()} manda el mensaje a la DLQ y cerrar la sesión lo reencola.
\end{itemize}
Se acepta además una segunda pila que mantener en la frontera. Esa pila duplica en Java la lectura del sobre y la parte del catálogo que usa el adaptador. Los esquemas JSON siguen siendo una sola fuente: el WAR los copia de \texttt{Plataforma.\br{}Comun.\br{}Eventos} al compilar.

\paragraph{Riesgo.} Un cambio del sobre canónico o de las claves de ruteo obliga a cambiar dos implementaciones. Lo detecta la validación contra los mismos esquemas, en las pruebas y al consumir. Corresponde al riesgo R-15 del ATAM.
Mientras Notificación no exista, el adaptador recibe EV-03 y EV-04 como sustituto temporal (L-11).

\paragraph{Dónde se ve.} Vista lógica: componente Adaptador de SE-8 [\texttt{adaptador-se8}] «implementado».
Vista de procesos: proceso \texttt{adaptador-se8} ×1, siempre encendido.
Vista física de la implementación: entorno de ejecución WildFly 42 dentro del contenedor \texttt{adaptador-se8} (\cref{fig:vista-despliegue-impl}, \cref{tab:nodos-impl}).
Vista de implementación: paquete \texttt{adaptador-\br{}se8-\br{}jee} (\cref{fig:vista-implementacion}).
Modelo de integración: fila de SE-8 del \cref{tab:integracion}.

\paragraph{Evidencia en el código.} 52 pruebas unitarias y 7 de integración en verde. Las de integración usan Testcontainers con la topología real del bus (cola quorum con \texttt{delivery-limit} 5) y la base del migrador. Demuestran el ack después del commit JPA, la DLQ a la quinta entrega fallida, el duplicado sin segundo envío y el circuito con EV-23.
La prueba \texttt{prueba-se8.sh} mide contra el compose: \qty{1,8}{\second} del registro a la alerta entregada con acuse, circuito abierto a los \qty{25,6}{\second} de la caída, y cola en 0 a los \qty{12,9}{\second} de la vuelta, con las 20 alertas entregadas una sola vez.
El \cref{tab:que-se-modifica} registra los 0 componentes existentes modificados.
Archivos de evidencia:
{\raggedright
\begin{itemize}
  \item \texttt{proyectoparte2/\br{}adaptadores/\br{}adaptador-se8-jee/\br{}adaptador/\br{}src/\br{}main/\br{}java/\br{}plataforma/\br{}adaptador/\br{}se8/\br{}jee/\br{}AdaptadorSe8.java}
  \item \texttt{proyectoparte2/\br{}adaptadores/\br{}adaptador-se8-jee/\br{}adaptador/\br{}src/\br{}main/\br{}java/\br{}plataforma/\br{}adaptador/\br{}se8/\br{}consumo/\br{}BucleConsumo.java}
  \item \texttt{proyectoparte2/\br{}adaptadores/\br{}adaptador-se8-jee/\br{}adaptador/\br{}src/\br{}main/\br{}java/\br{}plataforma/\br{}adaptador/\br{}se8/\br{}resiliencia/\br{}Cortacircuitos.java}
  \item \texttt{proyectoparte2/\br{}adaptadores/\br{}adaptador-se8-jee/\br{}Dockerfile}
  \item \texttt{proyectoparte2/\br{}pruebas/\br{}humo/\br{}prueba-se8.sh}
\end{itemize}
\par}

%------------------------------------------------------------------------------
\subsection{Restricciones de las tecnologías elegidas}
\label{sec:tec-restricciones}

Las tecnologías elegidas imponen límites al entorno implementado.
El \cref{tab:tec-restricciones} los relaciona con la decisión que los origina y con el escenario afectado.
Ninguno modifica una decisión estructural.

\begingroup\small
\begin{longtable}{L{2.6cm} L{1.3cm} L{6.4cm} L{1.7cm} L{3.3cm}}
\cabecera{5}{tab:tec-restricciones}{Restricciones que imponen las tecnologías elegidas}{\textbf{Restricción} &
\textbf{Origen} & \textbf{Efecto} & \textbf{QA afectado} & \textbf{Qué haría falta}}
Variante \texttt{debug} y Keycloak sin TLS (L-07) & AD-58, AD-61, AD-63 &
La app de entrega admite \texttt{http://}. Keycloak corre sin TLS. El tráfico entre el teléfono, el gateway y el proveedor de identidad viaja sin cifrar. Corresponde al riesgo R-12 del ATAM. &
QA-11, QA-12 & TLS en el borde y URL HTTPS en la variante \texttt{release} \\
Emisor del token en \texttt{localhost} (L-08) & AD-57, AD-58 &
El \texttt{iss} del token es \texttt{http://localhost:8180}. El emulador exige \texttt{adb reverse}. Un teléfono físico exige cambiar \texttt{KC\_HOSTNAME} e \texttt{Identidad\_\_Emisor} a la vez. Corresponde al riesgo R-13 del ATAM. &
QA-11, QA-14 & Un nombre DNS común para Keycloak \\
Imagen comunitaria de MinIO (DT-11) & AD-56 &
El proyecto original dejó de publicar imágenes. El archivo de composición usa el fork \texttt{pgsty/minio}, cuya continuidad no está garantizada. &
QA-02, QA-14 & Sustituir por otro almacén con API S3 sin cambiar el código \\
Política de uso de las teselas de OpenStreetMap & AD-60 &
Las teselas públicas no tienen acuerdo de servicio, exigen atribución y prohíben el uso intensivo. Durante una emergencia con 10\,000 usuarios la carga del mapa dependería de un tercero. &
QA-05, QA-06 & Servidor de teselas propio o proveedor con acuerdo \\
Keycloak en \texttt{start-dev} & AD-58 &
El modo de desarrollo usa una base embebida, no exige TLS y no está pensado para producción. El realm se importa desde \texttt{realm-plataforma.json} en cada arranque. Es una sola instancia. &
QA-02, QA-11 & Modo \texttt{start} con base PostgreSQL, TLS y dos instancias \\
Sustituto temporal de Notificación (L-11) & AD-66 &
EV-03 y EV-04 llegan a \texttt{q.salida.se8} por dos enlaces temporales, y el adaptador alerta a un único destino de difusión. No hay suscripciones ni resolución por ámbito, y el estado de entrega se consulta en el adaptador. &
QA-08 & Implementar Notificación (AD-13), quitar los enlaces y poner \texttt{Sustituto\_\_Activo=false} \\
Dos pilas en la frontera & AD-66 &
El adaptador de SE-8 duplica en Java la lectura del sobre y la parte del catálogo que usa. Un cambio del contrato se hace en .NET y en Java. Corresponde al riesgo R-15 del ATAM. &
QA-09, QA-10 & Mantener los esquemas JSON como única fuente y validarlos en las pruebas de las dos pilas \\
\end{longtable}
\endgroup
.
%  Generado a partir de la especificación de la versión 1.2. No carga paquetes.
%==============================================================================
\section{Decisiones de tecnología}
\label{sec:tecnologia}

El enunciado del proyecto (§8) pide seleccionar las tecnologías en la segunda entrega.
CON-8 exigía independencia tecnológica en la primera entrega.
Esta sección sustituye esa condición para la segunda entrega.
El Anexo~\ref{sec:anexo-a} deja de ser una excepción a CON-8 y describe la implementación de FP-06 con las tecnologías elegidas aquí.

Las decisiones se numeran de AD-50 a AD-66.
AD-66 responde a CON-9, que se agregó después de AD-50 a AD-65.
Ninguna modifica una decisión estructural.
Cada una materializa uno o varios mecanismos de AD-01 a AD-49, y su justificación parte de los escenarios QA-nn que ese mecanismo debe cumplir (CON-7).

\subsection{Método de decisión}
\label{sec:tec-metodo}

Cada decisión sigue cinco pasos:
\begin{enumerate}
  \item Se identifica el mecanismo del SAD que la tecnología materializa y los escenarios QA-nn que ese mecanismo atiende.
  \item Se eligen entre tres y cinco alternativas con uso real para ese mecanismo.
  \item Se definen los criterios. Cada criterio cita el QA-nn, AD-nn o CON-n del que proviene. La experiencia del equipo y la licencia o el costo se incluyen con un peso menor que los criterios de calidad.
  \item Se asigna un peso a cada criterio. Los pesos de una decisión suman 100\,\%.
  \item Se califica cada alternativa en cada criterio y se calcula el total ponderado.
\end{enumerate}

La escala de calificación tiene cinco niveles:
\begin{itemize}
  \item 1: la alternativa no satisface el criterio o lo impide.
  \item 2: lo satisface con desarrollo propio considerable.
  \item 3: lo satisface con configuración adicional o con complementos de terceros.
  \item 4: lo satisface de forma nativa con limitaciones menores.
  \item 5: lo satisface de forma nativa y la implementación lo verifica.
\end{itemize}

El total ponderado de una alternativa es $\sum_i p_i \, c_i / 100$, donde $p_i$ es el peso en porcentaje y $c_i$ la calificación del criterio $i$.
El máximo es 5.
Las cifras de verificación citadas provienen del perfil local de un computador, con 20 clientes, y no reemplazan la prueba con 10\,000 usuarios (L-04).

\subsection{Resumen de las decisiones}
\label{sec:tec-resumen}

El \cref{tab:tecnologia-resumen} relaciona cada decisión con su tecnología, el mecanismo que materializa y los escenarios que favorece.
Las versiones son las de los archivos de proyecto.

\begingroup\small
\begin{longtable}{L{1.15cm} L{2.2cm} L{6.0cm} L{2.6cm} L{2.6cm}}
\cabecera{5}{tab:tecnologia-resumen}{Resumen de las decisiones de tecnología}{\textbf{ID} & \textbf{Categoría} &
\textbf{Tecnología y versión} & \textbf{Mecanismo} & \textbf{QA que favorece}}
AD-50 & Lenguaje y plataforma & C\# sobre .NET 10 (SDK 10.0.100, \texttt{net10.0}) & AD-01, AD-08, AD-09, AD-45 & QA-06, QA-07, QA-14 \\
AD-51 & Marco web & ASP.NET Core 10 (controladores y \emph{minimal APIs}) & AD-08, AD-22, AD-24 & QA-06, QA-07, QA-11 \\
AD-52 & Persistencia & PostgreSQL 18 con PostGIS 3.6 (\texttt{postgis/postgis:18-3.6}), una base por microservicio y réplica \texttt{db-replica} & AD-01, AD-17, AD-20, AD-24, AD-41, AD-45 & QA-03, QA-04, QA-06, QA-13, QA-15 \\
AD-53 & Acceso a datos & Entity Framework Core 10.0.14, Npgsql 10.0.3 y NetTopologySuite 10.0.3, con contenedor migrador & AD-08, AD-17, AD-41, AD-45 & QA-06, QA-10, QA-13 \\
AD-54 & Mensajería & RabbitMQ 4.4 (3 nodos, colas quorum acotadas, DLQ) con RabbitMQ.Client 7.2.2 & AD-02, AD-05, AD-36, AD-45 & QA-01, QA-04, QA-08 \\
AD-55 & Caché & Redis 8.8 (\texttt{redis:8.8.4-alpine}) con StackExchange.Redis 3.4.29 & AD-16, AD-17 & QA-05, QA-06 \\
AD-56 & Almacén de objetos & MinIO con API S3 (paquete Minio 7.0.0) & AD-24, AD-31 & QA-02, QA-03, QA-14 \\
AD-57 & Puerta de entrada & YARP 2.3.0 con catálogo \texttt{rutas.json} (56 rutas) & AD-04, AD-19, AD-22 & QA-05, QA-06, QA-11 \\
AD-58 & Identidad & Keycloak 26.5.2 (realm \texttt{plataforma}), AppAuth-Android 0.11.1, angular-auth-oidc-client 22.0.2, cliente OpenID Connect de ASP.NET Core para el tablero y validador JWT propio en \texttt{Plataforma.Comun.Seguridad} & AD-04, AD-22, AD-49 & QA-11, QA-14 \\
AD-59 & Cliente web & Angular 22.2 (TypeScript 6.0) servido por nginx 1.29 detrás del gateway & AD-04, AD-08 & QA-06, QA-10, QA-11 \\
AD-60 & Mapas & Leaflet 1.9.4 con teselas de OpenStreetMap & AD-16, AD-17 & QA-06 \\
AD-61 & Cliente móvil & Android nativo: Kotlin 2.0.21, Jetpack Compose (BOM 2024.09.00), Hilt 2.60.1, minSdk 26, targetSdk 36 & AD-31, AD-44, AD-48 & QA-07, QA-14 \\
AD-62 & Persistencia móvil & Room 2.8.5 sobre SQLite, migraciones explícitas 1$\to$2$\to$3 & AD-31, AD-44, AD-48 & QA-14 \\
AD-63 & Sincronización móvil & WorkManager 2.12.0, Retrofit 3.0.0, OkHttp 5.4.0 y kotlinx.serialization 1.9.0 & AD-31, AD-44, AD-48 & QA-11, QA-14 \\
AD-64 & Resiliencia & Microsoft.Extensions.Http.Resilience 10.11.0 (Polly v8) & AD-05, AD-37 & QA-01, QA-04 \\
AD-65 & Despliegue & Contenedores Docker (\texttt{mcr.microsoft.com/\br{}dotnet/\br{}aspnet:10.0}), Docker Compose para un computador y Kubernetes con HPA y KEDA para escalar & AD-18, AD-24, AD-25 & QA-01, QA-02, QA-05 \\
AD-66 & Adaptador de SE-8 & Jakarta EE 11 sobre WildFly 42.0.0.Final (Java 21): Jakarta Messaging con el cliente JMS de RabbitMQ 3.9.1, EJB, Jakarta Persistence, Jakarta REST y MicroProfile Config & AD-03, AD-05, AD-09, AD-25, AD-45 & QA-08, QA-09, CON-9 \\
\end{longtable}
\endgroup

%---------------------------------------------------------------------- AD-50
\subsection{AD-50. Lenguaje y plataforma de los microservicios}
\label{sec:tec-ad50}

\paragraph{Contexto.} Esta decisión materializa los microservicios con almacén propio (AD-01), las cuatro capas de cada microservicio (AD-08), el contrato común de los adaptadores (AD-09) y la bandeja de salida transaccional (AD-45).
El lenguaje determina la latencia del camino síncrono, que exige QA-06 (p95 $<$ 2\,s) y QA-07 (p95 $<$ 3\,s).
También determina si la base común \texttt{comun} y el prototipo de FP-06 del Anexo~\ref{sec:anexo-a} se reutilizan, lo que afecta la verificación de QA-14.
CON-2 exige integrarse con sistemas de tecnologías diferentes, y esa heterogeneidad la absorben los adaptadores (AD-03), no el lenguaje de los microservicios.

\paragraph{Alternativas.} Se consideran C\# sobre .NET 10, Java 21, Node.js 22 con TypeScript y Go.
Las cuatro tienen licencia libre y ejecución en contenedores.
El \cref{tab:tec-ad50} presenta la matriz de decisión ponderada.
La alternativa elegida obtiene 4,90 y la segunda mejor, Java 21, obtiene 3,85.

\begingroup\small
\begin{longtable}{L{4.96cm} C{1.10cm} C{2.25cm} C{2.25cm} C{2.25cm} C{2.25cm}}
\cabecera{6}{tab:tec-ad50}{Matriz de decisión ponderada de AD-50}{\textbf{Criterio (driver)} & \textbf{Peso} & \hspace{0pt}\textbf{.NET 10} & \hspace{0pt}\textbf{Java 21} & \hspace{0pt}\textbf{Node.js 22} & \hspace{0pt}\textbf{Go}}
Rendimiento del camino síncrono (QA-06, QA-07) & 25\,\% & 5 & 4 & 3 & 5 \\
Bibliotecas maduras para bus, resiliencia, datos espaciales y JWT (AD-02, AD-05, AD-17, AD-22) & 25\,\% & 5 & 5 & 3 & 3 \\
Reutilización de \texttt{comun} y del prototipo de FP-06 (QA-14, AD-09) & 20\,\% & 5 & 2 & 2 & 2 \\
Memoria y arranque de cada réplica (QA-05, AD-18) & 10\,\% & 4 & 3 & 4 & 5 \\
Experiencia del equipo & 10\,\% & 5 & 4 & 3 & 1 \\
Licencia y costo & 10\,\% & 5 & 5 & 5 & 5 \\
\midrule
\textbf{Total ponderado} & 100\,\% & \textbf{4,90} & 3,85 & 3,10 & 3,50 \\
\end{longtable}
\endgroup

\paragraph{Alternativas descartadas.}
\begin{itemize}
  \item Java 21 pierde en la reutilización de \texttt{comun} y del prototipo de FP-06 (QA-14), que habría que reescribir.
  \item Node.js 22 pierde en el rendimiento del camino síncrono (QA-06, QA-07) y en la madurez de las bibliotecas para datos espaciales y bus.
  \item Go gana en memoria y arranque de cada réplica (QA-05), pero pierde en bibliotecas para los mecanismos del SAD y en experiencia del equipo.
\end{itemize}

\paragraph{Compromiso aceptado.} Una sola pila reduce el costo de mantener \texttt{comun} (AD-09) y acepta que todos los microservicios dependen del ciclo de soporte de .NET.
La heterogeneidad que exige CON-2 queda confinada en los adaptadores (AD-03).
El núcleo sigue en .NET aun con CON-9: Jakarta EE se confina a la frontera según AD-03 y AD-66.

\paragraph{Riesgo.} El fin del soporte de .NET 10 obligaría a migrar todos los microservicios a la vez.
El archivo \texttt{global.json} fija el SDK con avance solo de la versión de características, lo que acota este riesgo.
No corresponde a un riesgo del ATAM.

\paragraph{Dónde se ve.} Vista de implementación: proyectos \texttt{*.Api}, \texttt{*.Aplicacion}, \texttt{*.Dominio} e \texttt{*.Infraestructura} de cada microservicio.
Vista física de la implementación: artefactos \texttt{plataforma/<servicio>:local} sobre \texttt{mcr.microsoft.com/\br{}dotnet/\br{}aspnet:10.0}.

\paragraph{Evidencia en el código.} La suite de .NET suma 1\,008 pruebas en 19 proyectos en verde. QA-07 se midió con p95 de 1\,334\,ms en el perfil local (L-04).
Archivos de evidencia:
{\raggedright
\begin{itemize}
  \item \texttt{proyectoparte2/\br{}global.json}
  \item \texttt{proyectoparte2/\br{}Directory.Build.props}
\end{itemize}
\par}

%---------------------------------------------------------------------- AD-51
\subsection{AD-51. Marco de las API de los microservicios}
\label{sec:tec-ad51}

\paragraph{Contexto.} Esta decisión materializa la capa de acceso de cada microservicio (AD-08), el nivel 2 de la autorización (AD-22) y el sondeo de salud de la réplica activa (AD-24).
El marco atiende la latencia de QA-06 y QA-07 y el rechazo del 100\,\% de las modificaciones no autorizadas de QA-11.
La elección está condicionada por AD-50, porque el marco debe pertenecer a la pila elegida.

\paragraph{Alternativas.} Se consideran ASP.NET Core 10, Spring Boot 3 y NestJS.
Spring Boot 3 y NestJS se evalúan como el marco que se habría usado con Java y con Node.js.
El \cref{tab:tec-ad51} presenta la matriz de decisión ponderada.
La alternativa elegida obtiene 4,85 y la segunda mejor, Spring Boot 3, obtiene 3,75.

\begingroup\small
\begin{longtable}{L{6.74cm} C{1.10cm} C{2.50cm} C{2.50cm} C{2.50cm}}
\cabecera{5}{tab:tec-ad51}{Matriz de decisión ponderada de AD-51}{\textbf{Criterio (driver)} & \textbf{Peso} & \hspace{0pt}\textbf{ASP.NET Core 10} & \hspace{0pt}\textbf{Spring Boot 3} & \hspace{0pt}\textbf{NestJS}}
Latencia por petición (QA-06, QA-07) & 25\,\% & 5 & 4 & 3 \\
Validación JWT y autorización de nivel 2 en la tubería (QA-11, AD-22) & 20\,\% & 5 & 5 & 4 \\
Sondeos de salud vivo y listo (QA-02, QA-04, AD-24) & 15\,\% & 4 & 5 & 4 \\
Coherencia con AD-50 y con el gateway (QA-10, AD-57) & 20\,\% & 5 & 1 & 1 \\
Experiencia del equipo & 10\,\% & 5 & 3 & 3 \\
Licencia y costo & 10\,\% & 5 & 5 & 5 \\
\midrule
\textbf{Total ponderado} & 100\,\% & \textbf{4,85} & 3,75 & 3,15 \\
\end{longtable}
\endgroup

\paragraph{Alternativas descartadas.}
\begin{itemize}
  \item Spring Boot 3 gana en sondeos de salud con Actuator (QA-02), pero pierde en coherencia con AD-50, porque introduce una segunda pila.
  \item NestJS pierde en latencia por petición (QA-06, QA-07) y en coherencia con AD-50.
\end{itemize}

\paragraph{Compromiso aceptado.} Los sondeos de salud se componen con \texttt{Plataforma.Comun.Salud} en lugar de un módulo integrado como Actuator (AD-24).

\paragraph{Riesgo.} Una actualización mayor del marco afecta a todos los microservicios a la vez.
No corresponde a un riesgo del ATAM.

\paragraph{Dónde se ve.} Vista de implementación: proyectos \texttt{*.Api} con \texttt{Sdk.Web}.
Vista física de la implementación: artefactos \texttt{plataforma/<servicio>:local}.

\paragraph{Evidencia en el código.} Los proyectos \texttt{*.Api.csproj} usan \texttt{Microsoft.NET.Sdk.Web}. QA-07 se cumple con 100 de 100 registros respondidos con 201.
Archivos de evidencia:
{\raggedright
\begin{itemize}
  \item \texttt{proyectoparte2/\br{}servicios/\br{}gestion-desastres/\br{}src/\br{}GestionDesastres.Api}
  \item \texttt{proyectoparte2/\br{}comun/\br{}Plataforma.Comun.Salud}
\end{itemize}
\par}

%---------------------------------------------------------------------- AD-52
\subsection{AD-52. Motor de base de datos}
\label{sec:tec-ad52}

\paragraph{Contexto.} Esta decisión materializa el almacén propio de cada microservicio (AD-01), la geometría precalculada de las zonas (AD-17), los historiales versionados y de solo inserción (AD-20, AD-41), la réplica activa (AD-24) y la bandeja de salida (AD-45).
El motor atiende la consulta espacial de QA-06, la publicación sin pérdida de QA-04, la inmutabilidad de QA-13 y QA-15 y el RPO $\le$ 60\,s de QA-03.

\paragraph{Alternativas.} Se consideran PostgreSQL con PostGIS, MySQL 8, SQL Server y MongoDB.
El \cref{tab:tec-ad52} presenta la matriz de decisión ponderada.
La alternativa elegida obtiene 4,85 y la segunda mejor, SQL Server, obtiene 4,10.

\begingroup\small
\begin{longtable}{L{4.96cm} C{1.10cm} C{2.25cm} C{2.25cm} C{2.25cm} C{2.25cm}}
\cabecera{6}{tab:tec-ad52}{Matriz de decisión ponderada de AD-52}{\textbf{Criterio (driver)} & \textbf{Peso} & \hspace{0pt}\textbf{PostgreSQL + PostGIS} & \hspace{0pt}\textbf{MySQL 8} & \hspace{0pt}\textbf{SQL Server} & \hspace{0pt}\textbf{MongoDB}}
Consulta espacial con índice (QA-06, AD-17) & 25\,\% & 5 & 3 & 4 & 3 \\
Transacción ACID para la bandeja y el historial (QA-04, AD-20, AD-45) & 20\,\% & 5 & 4 & 5 & 3 \\
Inmutabilidad garantizada por permisos del almacén (QA-13, QA-15, AD-41) & 15\,\% & 5 & 4 & 5 & 3 \\
Réplica continua para la réplica activa y el RPO (QA-03, AD-24, AD-27) & 15\,\% & 4 & 4 & 5 & 5 \\
Experiencia del equipo & 10\,\% & 5 & 4 & 3 & 3 \\
Licencia y costo & 15\,\% & 5 & 4 & 2 & 3 \\
\midrule
\textbf{Total ponderado} & 100\,\% & \textbf{4,85} & 3,75 & 4,10 & 3,30 \\
\end{longtable}
\endgroup

\paragraph{Alternativas descartadas.}
\begin{itemize}
  \item MySQL 8 pierde en consulta espacial con índice (QA-06), porque sus funciones espaciales son menos completas que PostGIS.
  \item SQL Server gana en réplica continua (QA-03), pero pierde en licencia y costo, porque su licencia comercial se cobra por núcleo.
  \item MongoDB gana en réplica continua (QA-03), pero pierde en transacción ACID para la bandeja (QA-04) y en inmutabilidad por permisos (QA-13).
\end{itemize}

\paragraph{Compromiso aceptado.} En el entorno de un computador las nueve bases comparten la instancia \texttt{db}.
La separación por microservicio se conserva por base y por rol con mínimo privilegio (AD-01), no por servidor.

\paragraph{Riesgo.} La instancia \texttt{db} es un punto único de fallo en el entorno de un computador.
La réplica \texttt{db-replica} reduce la pérdida de datos, pero la promoción no es automática y no hay segunda región (L-01).
QA-02 y QA-03 no se midieron.

\paragraph{Dónde se ve.} Vista física de la implementación: nodos \texttt{db} y \texttt{db-replica}.
Vista de implementación: paquete \texttt{Infraestructura/Persistencia} de cada microservicio y migradores.

\paragraph{Evidencia en el código.} Los servicios \texttt{db} y \texttt{db-replica} usan la imagen \texttt{postgis/postgis:18-3.6}. QA-06 se midió con p95 de 34\,ms en el consolidado del desastre en el perfil local (L-04).
Archivos de evidencia:
{\raggedright
\begin{itemize}
  \item \texttt{proyectoparte2/\br{}despliegue/\br{}compose/\br{}docker-compose.yml}
\end{itemize}
\par}

%---------------------------------------------------------------------- AD-53
\subsection{AD-53. Acceso a datos y migraciones}
\label{sec:tec-ad53}

\paragraph{Contexto.} Esta decisión materializa la capa de infraestructura de cada microservicio (AD-08), el mapeo de las geometrías (AD-17), los permisos del historial de solo inserción (AD-41) y la escritura atómica con la bandeja (AD-45).
El acceso a datos afecta la latencia de QA-06, la evolución del esquema que pide QA-10 y la inmutabilidad de QA-13.

\paragraph{Alternativas.} Se consideran Entity Framework Core, Dapper, NHibernate y SQL directo con Npgsql sin mapeador.
El \cref{tab:tec-ad53} presenta la matriz de decisión ponderada.
La alternativa elegida obtiene 4,80 y la segunda mejor, Dapper, obtiene 3,60.

\begingroup\small
\begin{longtable}{L{4.96cm} C{1.10cm} C{2.25cm} C{2.25cm} C{2.25cm} C{2.25cm}}
\cabecera{6}{tab:tec-ad53}{Matriz de decisión ponderada de AD-53}{\textbf{Criterio (driver)} & \textbf{Peso} & \hspace{0pt}\textbf{EF Core 10} & \hspace{0pt}\textbf{Dapper} & \hspace{0pt}\textbf{NHibernate} & \hspace{0pt}\textbf{SQL directo}}
Migraciones versionadas con los permisos del historial (QA-10, QA-13, AD-41) & 25\,\% & 5 & 2 & 3 & 1 \\
Tipos espaciales mapeados (QA-06, AD-17) & 20\,\% & 5 & 3 & 2 & 3 \\
Control del SQL en las lecturas críticas (QA-06) & 20\,\% & 4 & 5 & 3 & 5 \\
Aislamiento en la capa de infraestructura y pruebas (QA-10, AD-08) & 15\,\% & 5 & 4 & 4 & 3 \\
Experiencia del equipo & 10\,\% & 5 & 4 & 2 & 4 \\
Licencia y costo & 10\,\% & 5 & 5 & 5 & 5 \\
\midrule
\textbf{Total ponderado} & 100\,\% & \textbf{4,80} & 3,60 & 3,05 & 3,20 \\
\end{longtable}
\endgroup

\paragraph{Alternativas descartadas.}
\begin{itemize}
  \item Dapper gana en control del SQL (QA-06), pero pierde en migraciones versionadas (QA-10, QA-13), que exigirían otra herramienta.
  \item NHibernate pierde en tipos espaciales mapeados (QA-06) y en experiencia del equipo.
  \item SQL directo gana en control del SQL (QA-06), pero pierde en migraciones versionadas y en aislamiento de la capa de infraestructura (QA-10).
\end{itemize}

\paragraph{Compromiso aceptado.} El mapeador agrega una capa sobre el SQL.
Las lecturas críticas del modelo de lectura (AD-15) pueden usar SQL explícito dentro del mismo contexto cuando QA-06 lo exija.
El esquema lo aplica un contenedor migrador con el rol dueño, y el microservicio corre con un rol sin permisos de DDL (AD-41).

\paragraph{Riesgo.} Una consulta generada por el mapeador puede degradar la latencia del mapa de daños, que se midió con p95 de 1\,628\,ms frente al límite de 2\,s de QA-06.
No corresponde a un riesgo del ATAM.

\paragraph{Dónde se ve.} Vista de implementación: paquete \texttt{Infraestructura/Persistencia} y migradores (\texttt{*-migrador}).
Vista física de la implementación: nodos \texttt{db} y \texttt{db-replica}.

\paragraph{Evidencia en el código.} El servicio \texttt{reconstruccion-migrador} del archivo de composición crea los roles y aplica las migraciones antes de arrancar el microservicio.
Archivos de evidencia:
{\raggedright
\begin{itemize}
  \item \texttt{proyectoparte2/\br{}Directory.Packages.props}
  \item \texttt{proyectoparte2/\br{}comun/\br{}Plataforma.Comun.Migraciones}
  \item \texttt{proyectoparte2/\br{}servicios/\br{}evaluacion-danos/\br{}src/\br{}EvaluacionDanos.Infraestructura/\br{}Persistencia}
\end{itemize}
\par}

%---------------------------------------------------------------------- AD-54
\subsection{AD-54. Bus de eventos}
\label{sec:tec-ad54}

\paragraph{Contexto.} Esta decisión materializa el bus de eventos con colas por consumidor (AD-02), las colas acotadas de la frontera (AD-05), la cola de salida con reintento (AD-36) y la entrega desde la bandeja de salida (AD-45).
El bus atiende QA-04, cuya respuesta exige 0 eventos perdidos ante el fallo de un componente interno, y QA-01, que exige procesar el 100\,\% de lo encolado cuando el sistema externo vuelve.

\paragraph{Alternativas.} Se consideran RabbitMQ, Apache Kafka, NATS JetStream y MassTransit sobre RabbitMQ.
MassTransit se evalúa como capa de abstracción sobre el mismo intermediario.
El \cref{tab:tec-ad54} presenta la matriz de decisión ponderada.
La alternativa elegida obtiene 4,90 y la segunda mejor, NATS JetStream, obtiene 3,85.

\begingroup\small
\begin{longtable}{L{4.96cm} C{1.10cm} C{2.25cm} C{2.25cm} C{2.25cm} C{2.25cm}}
\cabecera{6}{tab:tec-ad54}{Matriz de decisión ponderada de AD-54}{\textbf{Criterio (driver)} & \textbf{Peso} & \hspace{0pt}\textbf{RabbitMQ 4.4} & \hspace{0pt}\textbf{Kafka} & \hspace{0pt}\textbf{NATS JetStream} & \hspace{0pt}\textbf{MassTransit}}
Cola acotada por consumidor con rechazo al productor (QA-04, AD-05) & 25\,\% & 5 & 2 & 4 & 4 \\
Confirmación por mensaje, límite de entregas y DLQ (QA-01, QA-04, AD-36) & 20\,\% & 5 & 2 & 4 & 5 \\
Replicación de la cola ante la caída de un nodo (QA-02, QA-04) & 15\,\% & 5 & 5 & 5 & 5 \\
Topología explícita y declarada al desplegar (QA-09, QA-10) & 15\,\% & 5 & 3 & 3 & 2 \\
Experiencia del equipo & 10\,\% & 4 & 2 & 1 & 3 \\
Licencia y costo & 15\,\% & 5 & 5 & 5 & 2 \\
\midrule
\textbf{Total ponderado} & 100\,\% & \textbf{4,90} & 3,05 & 3,85 & 3,65 \\
\end{longtable}
\endgroup

\paragraph{Alternativas descartadas.}
\begin{itemize}
  \item Kafka iguala en replicación (QA-04), pero pierde en cola acotada por consumidor y en confirmación por mensaje con DLQ, que no ofrece de forma nativa.
  \item NATS JetStream pierde en topología declarada (QA-09) y en experiencia del equipo.
  \item MassTransit gana en confirmación y DLQ (QA-04), pero pierde en topología explícita, porque la oculta, y en licencia, porque sus autores anunciaron licencia comercial para la versión 9.
\end{itemize}

\paragraph{Compromiso aceptado.} La retención de un registro de eventos, que Kafka ofrece, no existe en RabbitMQ.
La recuperación de los mensajes pendientes descansa en la bandeja de salida de cada productor (AD-45).
La política \texttt{plataforma-colas} acota cada cola a 100\,000 mensajes y rechaza la publicación al llenarse.

\paragraph{Riesgo.} Una cola llena rechaza la publicación y el relevo de la bandeja reintenta.
Si la saturación persiste, la bandeja crece en la base del productor.
La alarma de la bandeja no se registra en Auditoría (L-02).

\paragraph{Dónde se ve.} Vista de procesos: exchange \texttt{plataforma.eventos} y colas \texttt{q.*} con su \texttt{dlq.*}.
Vista física de la implementación: nodos \texttt{bus1}, \texttt{bus2} y \texttt{bus3}.

\paragraph{Evidencia en el código.} Las 71 pruebas de \texttt{Plataforma.Comun.Pruebas} con RabbitMQ real verifican consumidor detenido sin pérdida, mensaje envenenado apartado y cola llena sin descarte (QA-04).
Archivos de evidencia:
{\raggedright
\begin{itemize}
  \item \texttt{proyectoparte2/\br{}despliegue/\br{}bus/\br{}definiciones.json}
  \item \texttt{proyectoparte2/\br{}comun/\br{}Plataforma.Comun.Mensajeria}
\end{itemize}
\par}

%---------------------------------------------------------------------- AD-55
\subsection{AD-55. Motor de caché}
\label{sec:tec-ad55}

\paragraph{Contexto.} Esta decisión materializa la caché de geometrías y de respuestas del mapa (AD-17) y la marca de tiempo del dato consultado (AD-16).
El motor atiende QA-06, que exige p95 $<$ 2\,s y antigüedad del dato $\le$ 2\,s en el p95.
También atiende QA-05, porque la caché separa la carga de lectura de la base de datos.

\paragraph{Alternativas.} Se consideran Redis, Valkey, Memcached y una vista materializada en PostgreSQL.
El \cref{tab:tec-ad55} presenta la matriz de decisión ponderada.
La alternativa elegida obtiene 4,70 y la segunda mejor, Valkey, obtiene 4,45.

\begingroup\small
\begin{longtable}{L{4.96cm} C{1.10cm} C{2.25cm} C{2.25cm} C{2.25cm} C{2.25cm}}
\cabecera{6}{tab:tec-ad55}{Matriz de decisión ponderada de AD-55}{\textbf{Criterio (driver)} & \textbf{Peso} & \hspace{0pt}\textbf{Redis 8.8} & \hspace{0pt}\textbf{Valkey} & \hspace{0pt}\textbf{Memcached} & \hspace{0pt}\textbf{Vista materializada}}
Lectura con TTL e invalidación por clave (QA-06, AD-17) & 30\,\% & 5 & 5 & 4 & 3 \\
Antigüedad del dato controlada (QA-06, AD-16) & 20\,\% & 4 & 4 & 4 & 2 \\
Componente separado de la base que escala aparte (QA-05, AD-18) & 15\,\% & 5 & 5 & 5 & 1 \\
Cliente maduro en la pila de AD-50 & 15\,\% & 5 & 4 & 3 & 5 \\
Experiencia del equipo & 10\,\% & 5 & 3 & 2 & 4 \\
Licencia y costo & 10\,\% & 4 & 5 & 5 & 5 \\
\midrule
\textbf{Total ponderado} & 100\,\% & \textbf{4,70} & 4,45 & 3,90 & 3,10 \\
\end{longtable}
\endgroup

\paragraph{Alternativas descartadas.}
\begin{itemize}
  \item Valkey gana en licencia, porque conserva la licencia BSD, pero pierde en experiencia del equipo y en madurez del cliente probado con la pila.
  \item Memcached pierde en invalidación por clave (QA-06), porque no ofrece estructuras ni invalidación por patrón.
  \item La vista materializada gana en cliente y licencia, pero pierde en componente separado (QA-05), porque concentra la lectura en la misma base.
\end{itemize}

\paragraph{Compromiso aceptado.} Redis 8 se distribuye con licencias RSALv2, SSPL y AGPLv3.
La AGPLv3 permite su uso en este proyecto académico.
Valkey es un sustituto compatible con el mismo protocolo si la licencia cambia.

\paragraph{Riesgo.} Riesgo R-14 del ATAM: la antigüedad del dato presentado depende de la carga.
En el perfil local se midió p95 de 1\,485\,ms con máximo de 11,8\,s.
Una corrida anterior dio p95 de 4,2\,s con 10 muestras.

\paragraph{Dónde se ve.} Vista física de la implementación: nodo \texttt{cache}.

\paragraph{Evidencia en el código.} QA-06 se midió con p95 de 142\,ms en el mapa de desastres y con la marca \texttt{datoAl} en el 100\,\% de las respuestas.
Archivos de evidencia:
{\raggedright
\begin{itemize}
  \item \texttt{proyectoparte2/\br{}servicios/\br{}consulta-danos/\br{}src/\br{}ConsultaDanos.Infraestructura/\br{}Cache.cs}
  \item \texttt{proyectoparte2/\br{}servicios/\br{}consulta/\br{}src/\br{}Consulta.Infraestructura/\br{}Adaptadores.cs}
\end{itemize}
\par}

%---------------------------------------------------------------------- AD-56
\subsection{AD-56. Almacén de archivos}
\label{sec:tec-ad56}

\paragraph{Contexto.} Esta decisión materializa el almacén de archivos de la vista física, donde se guardan las evidencias de Evaluación de Daños y los documentos de Reconstrucción.
El almacén debe mantener la evidencia disponible aunque falle una réplica del microservicio (AD-24, QA-02).
También debe conservar el 100\,\% de las fotos sincronizadas desde el cliente móvil (AD-31, QA-14).

\paragraph{Alternativas.} Se consideran MinIO, el disco local del contenedor, una columna binaria en PostgreSQL y SeaweedFS.
El \cref{tab:tec-ad56} presenta la matriz de decisión ponderada.
La alternativa elegida obtiene 4,20 y la segunda mejor, SeaweedFS, obtiene 4,05.

\begingroup\small
\begin{longtable}{L{4.96cm} C{1.10cm} C{2.25cm} C{2.25cm} C{2.25cm} C{2.25cm}}
\cabecera{6}{tab:tec-ad56}{Matriz de decisión ponderada de AD-56}{\textbf{Criterio (driver)} & \textbf{Peso} & \hspace{0pt}\textbf{MinIO} & \hspace{0pt}\textbf{Disco local} & \hspace{0pt}\textbf{Columna binaria} & \hspace{0pt}\textbf{SeaweedFS}}
Evidencia disponible ante la pérdida de una réplica (QA-02, AD-24) & 25\,\% & 4 & 1 & 4 & 4 \\
Integridad de la evidencia sincronizada (QA-14) & 20\,\% & 5 & 3 & 5 & 5 \\
Tamaño de la base y de su respaldo (QA-03, QA-06) & 15\,\% & 5 & 4 & 1 & 5 \\
Interfaz S3 estándar que permite sustituir el producto (QA-10) & 15\,\% & 5 & 1 & 1 & 4 \\
Experiencia del equipo & 10\,\% & 4 & 5 & 4 & 1 \\
Licencia y continuidad del producto & 15\,\% & 2 & 5 & 5 & 4 \\
\midrule
\textbf{Total ponderado} & 100\,\% & \textbf{4,20} & 2,85 & 3,45 & 4,05 \\
\end{longtable}
\endgroup

\paragraph{Alternativas descartadas.}
\begin{itemize}
  \item El disco local gana en experiencia y licencia, pero pierde en disponibilidad ante la pérdida de una réplica (QA-02), que AD-24 descarta.
  \item La columna binaria iguala en integridad (QA-14), pero pierde en tamaño de la base y de su respaldo (QA-03).
  \item SeaweedFS gana en continuidad del producto, pero pierde en experiencia del equipo y en compatibilidad S3 verificada.
\end{itemize}

\paragraph{Compromiso aceptado.} El proyecto original de MinIO dejó de publicar imágenes y el archivo de composición usa el fork comunitario \texttt{pgsty/minio} (DT-11).
El código depende solo de la API S3, lo que permite sustituir el producto sin modificar los microservicios.

\paragraph{Riesgo.} El fork comunitario puede dejar de mantenerse.
En el entorno de un computador \texttt{archivos} es un solo nodo, y su caída deja sin evidencias a las réplicas de \texttt{evaluacion-danos}.
No corresponde a un riesgo del ATAM.

\paragraph{Dónde se ve.} Vista física de la implementación: nodo \texttt{archivos}.

\paragraph{Evidencia en el código.} La corrida de QA-14 conservó 120 de 120 fotos a través del gateway.
Archivos de evidencia:
{\raggedright
\begin{itemize}
  \item \texttt{proyectoparte2/\br{}servicios/\br{}evaluacion-danos/\br{}src/\br{}EvaluacionDanos.Infraestructura/\br{}Archivos/\br{}AlmacenEvidenciasMinio.cs}
  \item \texttt{proyectoparte2/\br{}servicios/\br{}reconstruccion/\br{}src/\br{}Reconstruccion.Infraestructura/\br{}Archivos/\br{}AlmacenDocumentosMinio.cs}
\end{itemize}
\par}

%---------------------------------------------------------------------- AD-57
\subsection{AD-57. API Gateway}
\label{sec:tec-ad57}

\paragraph{Contexto.} Esta decisión materializa el punto único de entrada (AD-04), la priorización y degradación por tipo de cliente (AD-19) y el nivel 1 de la autorización (AD-22).
El gateway atiende QA-11, que exige el rechazo del 100\,\% de las modificaciones de un organismo no responsable.
También atiende QA-05 y QA-06, porque prioriza el tráfico crítico cuando la carga crece.

\paragraph{Alternativas.} Se consideran YARP, Kong, Traefik, Ocelot y nginx.
El \cref{tab:tec-ad57} presenta la matriz de decisión ponderada.
La alternativa elegida obtiene 4,85 y la segunda mejor, Kong, obtiene 3,45.

\begingroup\small
\begin{longtable}{L{3.93cm} C{1.10cm} C{1.95cm} C{1.95cm} C{1.95cm} C{1.95cm} C{1.95cm}}
\cabecera{7}{tab:tec-ad57}{Matriz de decisión ponderada de AD-57}{\textbf{Criterio (driver)} & \textbf{Peso} & \hspace{0pt}\textbf{YARP 2.3} & \hspace{0pt}\textbf{Kong} & \hspace{0pt}\textbf{Traefik} & \hspace{0pt}\textbf{Ocelot} & \hspace{0pt}\textbf{nginx}}
Priorización por clase de tráfico y degradación (QA-05, QA-06, AD-19) & 25\,\% & 5 & 3 & 2 & 3 & 2 \\
Validación JWT y regla de dominio por ruta (QA-11, AD-22) & 25\,\% & 5 & 4 & 3 & 4 & 2 \\
Catálogo de rutas validado al arrancar (QA-11) & 15\,\% & 5 & 4 & 3 & 3 & 2 \\
Latencia agregada por el intermediario (QA-06) & 15\,\% & 4 & 4 & 4 & 3 & 5 \\
Experiencia del equipo en la pila de AD-50 & 10\,\% & 5 & 2 & 3 & 4 & 3 \\
Licencia y mantenimiento & 10\,\% & 5 & 3 & 5 & 3 & 5 \\
\midrule
\textbf{Total ponderado} & 100\,\% & \textbf{4,85} & 3,45 & 3,10 & 3,35 & 2,85 \\
\end{longtable}
\endgroup

\paragraph{Alternativas descartadas.}
\begin{itemize}
  \item Kong pierde en priorización con degradación (QA-05), que exigiría complementos propios en otro lenguaje.
  \item Traefik pierde en priorización y en validación de la regla de dominio por ruta (QA-11).
  \item Ocelot pierde en catálogo validado al arrancar (QA-11) y en mantenimiento del proyecto.
  \item nginx gana en latencia agregada (QA-06), pero pierde en validación JWT con regla de dominio (QA-11), que no ofrece sin módulos adicionales.
\end{itemize}

\paragraph{Compromiso aceptado.} La priorización y la regla de dominio son código propio en el gateway.
Ese código se prueba con la misma suite de .NET, pero lo mantiene el equipo.

\paragraph{Riesgo.} Riesgo R-12 del ATAM: el gateway recibe tráfico sin TLS en el entorno de laboratorio (L-07).
Riesgo R-13 del ATAM: la validación depende de las llaves y del emisor de Keycloak (L-08).
Los rechazos del gateway no se registran en Auditoría (L-02).

\paragraph{Dónde se ve.} Vista lógica: componente API Gateway.
Vista física de la implementación: nodo \texttt{gateway} con dos réplicas y el artefacto \texttt{rutas.json}.

\paragraph{Evidencia en el código.} QA-11 se verificó con 218 de 218 intentos rechazados (100\,\%), 0 falsos negativos y 0 datos modificados sobre 12 rutas con regla de dominio.
Archivos de evidencia:
{\raggedright
\begin{itemize}
  \item \texttt{proyectoparte2/\br{}gateway/\br{}src/\br{}Plataforma.Gateway/\br{}rutas.json}
  \item \texttt{proyectoparte2/\br{}gateway/\br{}src/\br{}Plataforma.Gateway/\br{}Prioridad/\br{}Priorizacion.cs}
  \item \texttt{proyectoparte2/\br{}gateway/\br{}src/\br{}Plataforma.Gateway/\br{}Seguridad/\br{}AutorizacionNivel1.cs}
\end{itemize}
\par}

%---------------------------------------------------------------------- AD-58
\subsection{AD-58. Proveedor de identidad y clientes OIDC}
\label{sec:tec-ad58}

\paragraph{Contexto.} Esta decisión materializa el proveedor de identidad federado con OIDC y PKCE (AD-49) y la validación del token en el gateway y en cada microservicio (AD-04, AD-22).
Sus clientes son la consola, el cliente móvil y el tablero de daños. El tablero usa el cliente confidencial \texttt{consulta-\br{}danos-\br{}tablero} y una sesión con cookie que valida la propia página, porque su ruta es pública en el gateway (E-41).
El proveedor atiende QA-11, porque emite los roles y el atributo de organismo que usa la regla de dominio.
También atiende QA-14, porque el refresco del token sin conexión permite trabajar días en campo y sincronizar al volver.

\paragraph{Alternativas.} Se consideran Keycloak, Auth0, Duende IdentityServer y una identidad propia.
El \cref{tab:tec-ad58} presenta la matriz de decisión ponderada.
La alternativa elegida obtiene 4,90 y la segunda mejor, Duende IdentityServer, obtiene 3,75.

\begingroup\small
\begin{longtable}{L{4.96cm} C{1.10cm} C{2.25cm} C{2.25cm} C{2.25cm} C{2.25cm}}
\cabecera{6}{tab:tec-ad58}{Matriz de decisión ponderada de AD-58}{\textbf{Criterio (driver)} & \textbf{Peso} & \hspace{0pt}\textbf{Keycloak 26} & \hspace{0pt}\textbf{Auth0} & \hspace{0pt}\textbf{Duende IdentityServer} & \hspace{0pt}\textbf{Identidad propia}}
OIDC con PKCE para la consola y el móvil, con refresco sin conexión (QA-11, QA-14, AD-49) & 25\,\% & 5 & 5 & 4 & 2 \\
Roles y organismo como \emph{claims} para la regla de dominio (QA-11, AD-22) & 20\,\% & 5 & 4 & 4 & 3 \\
Despliegue dentro de la plataforma, sin servicio externo en la emergencia (QA-02, CON-3) & 15\,\% & 5 & 2 & 5 & 5 \\
Configuración declarada en un archivo versionado (QA-10) & 10\,\% & 5 & 3 & 4 & 3 \\
Experiencia del equipo & 10\,\% & 4 & 2 & 2 & 3 \\
Licencia y costo & 20\,\% & 5 & 2 & 3 & 5 \\
\midrule
\textbf{Total ponderado} & 100\,\% & \textbf{4,90} & 3,25 & 3,75 & 3,45 \\
\end{longtable}
\endgroup

\paragraph{Alternativas descartadas.}
\begin{itemize}
  \item Auth0 iguala en OIDC con PKCE (QA-11), pero pierde en despliegue dentro de la plataforma (CON-3) y en costo por usuario activo.
  \item Duende IdentityServer pierde en licencia, porque exige licencia comercial por encima de un umbral de ingresos, y en experiencia del equipo.
  \item La identidad propia gana en licencia y despliegue, pero pierde en OIDC con PKCE (QA-11), porque obligaría a construir la seguridad desde cero.
\end{itemize}

\paragraph{Compromiso aceptado.} El token se valida dos veces, en el gateway y en el microservicio, con un validador propio en \texttt{Plataforma.Comun.Seguridad}.
La doble validación agrega trabajo por petición a cambio de la defensa en profundidad de AD-22.

\paragraph{Riesgo.} Riesgo R-13 del ATAM: el emisor del token es \texttt{http://localhost:8180} y debe coincidir en el gateway y en los microservicios (L-08).
La recarga de las llaves JWT está pendiente (L-09).
Riesgo R-12 del ATAM: Keycloak corre en \texttt{start-dev} sin TLS (L-07).

\paragraph{Dónde se ve.} Vista lógica: componente externo Proveedor de identidad.
Vista física de la implementación: nodo \texttt{keycloak}.
Vista de implementación móvil: paquete \texttt{infraestructura/identidad}.

\paragraph{Evidencia en el código.} Siete pruebas instrumentadas de la app corren contra el gateway y Keycloak, incluido el rechazo por rol.
Archivos de evidencia:
{\raggedright
\begin{itemize}
  \item \texttt{proyectoparte2/\br{}despliegue/\br{}keycloak/\br{}realm-plataforma.json}
  \item \texttt{proyectoparte2/\br{}comun/\br{}Plataforma.Comun.Seguridad/\br{}IdentidadPorJwt.cs}
  \item \texttt{danos/\br{}app/\br{}src/\br{}main/\br{}java/\br{}com/\br{}example/\br{}danos/\br{}infraestructura/\br{}identidad/\br{}ConexionOidcBuilder.kt}
\end{itemize}
\par}

%---------------------------------------------------------------------- AD-59
\subsection{AD-59. Consola de coordinación}
\label{sec:tec-ad59}

\paragraph{Contexto.} Esta decisión materializa la capa de presentación web (AD-08) de la Consola de coordinación, con los módulos \texttt{mapa}, \texttt{evaluaciones} y \texttt{reconstruccion}.
La consola entra por el gateway (AD-04) y atiende QA-06 en la carga del mapa, QA-11 en el inicio de sesión y QA-10 en la incorporación de módulos.

\paragraph{Alternativas.} Se consideran Angular, React 19, Vue 3 y Blazor.
El \cref{tab:tec-ad59} presenta la matriz de decisión ponderada.
La alternativa elegida obtiene 4,80 y la segunda mejor, Vue 3, obtiene 4,10.

\begingroup\small
\begin{longtable}{L{4.96cm} C{1.10cm} C{2.25cm} C{2.25cm} C{2.25cm} C{2.25cm}}
\cabecera{6}{tab:tec-ad59}{Matriz de decisión ponderada de AD-59}{\textbf{Criterio (driver)} & \textbf{Peso} & \hspace{0pt}\textbf{Angular 22} & \hspace{0pt}\textbf{React 19} & \hspace{0pt}\textbf{Vue 3} & \hspace{0pt}\textbf{Blazor}}
Módulos con rutas, inyección y guardas (QA-10, AD-08) & 25\,\% & 5 & 3 & 4 & 4 \\
Cliente OIDC con PKCE certificado (QA-11, AD-49) & 20\,\% & 5 & 4 & 4 & 4 \\
Carga inicial y respuesta de la interfaz (QA-06) & 20\,\% & 4 & 4 & 5 & 2 \\
Archivos estáticos servidos detrás del gateway (QA-02, AD-04) & 10\,\% & 5 & 5 & 5 & 4 \\
Experiencia del equipo & 15\,\% & 5 & 3 & 2 & 3 \\
Licencia y costo & 10\,\% & 5 & 5 & 5 & 5 \\
\midrule
\textbf{Total ponderado} & 100\,\% & \textbf{4,80} & 3,80 & 4,10 & 3,55 \\
\end{longtable}
\endgroup

\paragraph{Alternativas descartadas.}
\begin{itemize}
  \item React 19 pierde en estructura por módulos con inyección y guardas (QA-10), que exige elegir y combinar bibliotecas.
  \item Vue 3 gana en carga inicial (QA-06), pero pierde en experiencia del equipo y en estructura por módulos.
  \item Blazor pierde en carga inicial (QA-06), por el tamaño del entorno WebAssembly que descarga el navegador.
\end{itemize}

\paragraph{Compromiso aceptado.} Angular produce un paquete inicial mayor que Vue.
La carga se mitiga con módulos cargados bajo demanda por ruta.

\paragraph{Riesgo.} Angular publica una versión mayor cada seis meses y obliga a actualizaciones frecuentes.
No corresponde a un riesgo del ATAM.

\paragraph{Dónde se ve.} Vista física de la implementación: nodo \texttt{web}.
Vista lógica: Consola de coordinación con los módulos \texttt{mapa}, \texttt{evaluaciones} y \texttt{reconstruccion}.

\paragraph{Evidencia en el código.} El \texttt{Dockerfile} compila con \texttt{node:22-alpine} y sirve los archivos con \texttt{nginx:1.29-alpine}.
Archivos de evidencia:
{\raggedright
\begin{itemize}
  \item \texttt{proyectoparte2/\br{}consola/\br{}package.json}
  \item \texttt{proyectoparte2/\br{}consola/\br{}Dockerfile}
  \item \texttt{proyectoparte2/\br{}consola/\br{}nginx.conf}
  \item \texttt{proyectoparte2/\br{}consola/\br{}src/\br{}app/\br{}modulos}
\end{itemize}
\par}

%---------------------------------------------------------------------- AD-60
\subsection{AD-60. Biblioteca de mapas}
\label{sec:tec-ad60}

\paragraph{Contexto.} Esta decisión materializa la presentación de las geometrías precalculadas (AD-17) y de su marca de tiempo (AD-16) en el módulo \texttt{mapa} de la consola.
La biblioteca atiende QA-06, porque el mapa es la consulta frecuente del coordinador.

\paragraph{Alternativas.} Se consideran Leaflet, OpenLayers, MapLibre GL y Google Maps.
El \cref{tab:tec-ad60} presenta la matriz de decisión ponderada.
La alternativa elegida obtiene 4,80 y la segunda mejor, OpenLayers, obtiene 3,90.

\begingroup\small
\begin{longtable}{L{4.96cm} C{1.10cm} C{2.25cm} C{2.25cm} C{2.25cm} C{2.25cm}}
\cabecera{6}{tab:tec-ad60}{Matriz de decisión ponderada de AD-60}{\textbf{Criterio (driver)} & \textbf{Peso} & \hspace{0pt}\textbf{Leaflet 1.9} & \hspace{0pt}\textbf{OpenLayers} & \hspace{0pt}\textbf{MapLibre GL} & \hspace{0pt}\textbf{Google Maps}}
Capas GeoJSON agregadas por zona (QA-06, AD-17) & 25\,\% & 5 & 5 & 5 & 4 \\
Peso de la biblioteca y tiempo de carga (QA-06) & 20\,\% & 5 & 3 & 2 & 3 \\
Independencia de un proveedor y de una llave (QA-02, CON-3) & 20\,\% & 4 & 4 & 4 & 1 \\
Integración con Angular y con el mapa del prototipo (QA-10) & 10\,\% & 5 & 3 & 3 & 3 \\
Experiencia del equipo & 10\,\% & 5 & 2 & 2 & 3 \\
Licencia y costo & 15\,\% & 5 & 5 & 5 & 2 \\
\midrule
\textbf{Total ponderado} & 100\,\% & \textbf{4,80} & 3,90 & 3,70 & 2,70 \\
\end{longtable}
\endgroup

\paragraph{Alternativas descartadas.}
\begin{itemize}
  \item OpenLayers iguala en capas GeoJSON (QA-06), pero pierde en peso de la biblioteca y en experiencia del equipo.
  \item MapLibre GL iguala en capas GeoJSON, pero pierde en peso de la biblioteca (QA-06), porque su ventaja con millones de puntos no hace falta con zonas agregadas.
  \item Google Maps pierde en independencia de un proveedor (CON-3) y en costo por carga de mapa.
\end{itemize}

\paragraph{Compromiso aceptado.} Leaflet dibuja en el navegador sin aceleración WebGL.
El volumen se controla con la agregación por zona del servidor (AD-17).

\paragraph{Riesgo.} Las teselas de OpenStreetMap no tienen acuerdo de servicio y su política de uso prohíbe el uso intensivo.
No corresponde a un riesgo del ATAM.

\paragraph{Dónde se ve.} Vista física de la implementación: nodo \texttt{web}.
Vista lógica: módulo \texttt{mapa} de la Consola de coordinación.

\paragraph{Evidencia en el código.} QA-06 se midió con p95 de 1\,628\,ms y p99 de 1\,935\,ms en el mapa de daños en el perfil local (L-04).
Archivos de evidencia:
{\raggedright
\begin{itemize}
  \item \texttt{proyectoparte2/\br{}consola/\br{}src/\br{}app/\br{}modulos/\br{}mapa/\br{}lienzo-mapa.ts}
  \item \texttt{proyectoparte2/\br{}consola/\br{}package.json}
\end{itemize}
\par}

%---------------------------------------------------------------------- AD-61
\subsection{AD-61. Plataforma del cliente móvil}
\label{sec:tec-ad61}

\paragraph{Contexto.} Esta decisión materializa la réplica local del cliente móvil (AD-31), la cola local de operaciones (AD-44) y el registro del desastre en campo (AD-48).
La plataforma atiende QA-14, cuya medida es 0 evaluaciones perdidas o sobrescritas en silencio al sincronizar 120 evaluaciones, y cuya respuesta exige 0 duplicados.

\paragraph{Alternativas.} Se consideran Android nativo con Kotlin, Flutter, React Native y Kotlin Multiplatform.
El \cref{tab:tec-ad61} presenta la matriz de decisión ponderada.
La alternativa elegida obtiene 4,60 y la segunda mejor, Kotlin Multiplatform, obtiene 4,00.

\begingroup\small
\begin{longtable}{L{4.96cm} C{1.10cm} C{2.25cm} C{2.25cm} C{2.25cm} C{2.25cm}}
\cabecera{6}{tab:tec-ad61}{Matriz de decisión ponderada de AD-61}{\textbf{Criterio (driver)} & \textbf{Peso} & \hspace{0pt}\textbf{Kotlin nativo} & \hspace{0pt}\textbf{Flutter} & \hspace{0pt}\textbf{React Native} & \hspace{0pt}\textbf{Kotlin Multiplatform}}
Trabajo en segundo plano y almacenamiento del sistema (QA-14, AD-44) & 25\,\% & 5 & 3 & 3 & 4 \\
Base local con transacciones y migraciones (QA-14, AD-31) & 20\,\% & 5 & 4 & 3 & 4 \\
Cámara y ubicación para las evidencias (QA-14) & 10\,\% & 5 & 4 & 4 & 4 \\
Reutilización del prototipo de FP-06 y de sus pruebas (QA-14) & 15\,\% & 5 & 1 & 1 & 4 \\
Portabilidad a otro sistema operativo móvil (QA-10) & 10\,\% & 1 & 5 & 5 & 4 \\
Experiencia del equipo & 10\,\% & 5 & 2 & 3 & 3 \\
Licencia y costo & 10\,\% & 5 & 5 & 5 & 5 \\
\midrule
\textbf{Total ponderado} & 100\,\% & \textbf{4,60} & 3,30 & 3,20 & 4,00 \\
\end{longtable}
\endgroup

\paragraph{Alternativas descartadas.}
\begin{itemize}
  \item Flutter gana en portabilidad (QA-10), pero pierde en reutilización del prototipo de FP-06 (QA-14), que habría que reescribir.
  \item React Native gana en portabilidad, pero pierde en base local transaccional (QA-14) y en reutilización del prototipo.
  \item Kotlin Multiplatform pierde en trabajo en segundo plano (QA-14), porque esa parte sigue siendo específica de cada plataforma.
\end{itemize}

\paragraph{Compromiso aceptado.} El cliente funciona solo en Android.
Un cliente para iOS exigiría otra implementación de la capa de presentación y de la sincronización.

\paragraph{Riesgo.} La app de entrega es la variante \texttt{debug}, que admite \texttt{http://} (L-07).
La evidencia solo admite fotos y no guarda un autor propio (L-10).

\paragraph{Dónde se ve.} Vista de implementación móvil: paquete \texttt{presentacion}.
Vista física de la implementación: \texttt{<<device>>} Teléfono Android con \texttt{DanosCampo.apk}.

\paragraph{Evidencia en el código.} La app \texttt{co.plataforma.danos} 2.1.0 tiene 81 pruebas unitarias y 16 instrumentadas en verde.
Archivos de evidencia:
{\raggedright
\begin{itemize}
  \item \texttt{danos/\br{}gradle/\br{}libs.versions.toml}
  \item \texttt{danos/\br{}app/\br{}build.gradle.kts}
\end{itemize}
\par}

%---------------------------------------------------------------------- AD-62
\subsection{AD-62. Base de datos local del cliente móvil}
\label{sec:tec-ad62}

\paragraph{Contexto.} Esta decisión materializa la réplica local (AD-31), la cola local de operaciones inmutables con copia completa del contenido (AD-44) y la tabla \texttt{registro\_desastre} del registro en campo (AD-48).
La base atiende QA-14, porque la escritura y su operación en la cola deben confirmarse en la misma transacción.

\paragraph{Alternativas.} Se consideran Room, Realm, SQLDelight y DataStore.
El \cref{tab:tec-ad62} presenta la matriz de decisión ponderada.
La alternativa elegida obtiene 4,90 y la segunda mejor, SQLDelight, obtiene 4,35.

\begingroup\small
\begin{longtable}{L{4.96cm} C{1.10cm} C{2.25cm} C{2.25cm} C{2.25cm} C{2.25cm}}
\cabecera{6}{tab:tec-ad62}{Matriz de decisión ponderada de AD-62}{\textbf{Criterio (driver)} & \textbf{Peso} & \hspace{0pt}\textbf{Room 2.8} & \hspace{0pt}\textbf{Realm} & \hspace{0pt}\textbf{SQLDelight} & \hspace{0pt}\textbf{DataStore}}
Transacción atómica entre la escritura y la cola local (QA-14, AD-44) & 30\,\% & 5 & 4 & 5 & 1 \\
Migraciones explícitas y verificables (QA-10, QA-14) & 20\,\% & 5 & 3 & 4 & 2 \\
Integración con Compose y corrutinas (AD-61) & 15\,\% & 5 & 4 & 4 & 4 \\
Continuidad del producto & 15\,\% & 5 & 1 & 5 & 5 \\
Experiencia del equipo & 10\,\% & 5 & 2 & 2 & 4 \\
Control del SQL verificado en compilación (QA-14) & 10\,\% & 4 & 2 & 5 & 1 \\
\midrule
\textbf{Total ponderado} & 100\,\% & \textbf{4,90} & 2,95 & 4,35 & 2,55 \\
\end{longtable}
\endgroup

\paragraph{Alternativas descartadas.}
\begin{itemize}
  \item Realm pierde en continuidad del producto, porque su fabricante retiró el SDK de dispositivos.
  \item SQLDelight gana en control del SQL verificado en compilación, pero pierde en migraciones con esquema exportado y en experiencia del equipo.
  \item DataStore pierde en transacción atómica con la cola (QA-14), porque guarda pares clave y valor sin relaciones.
\end{itemize}

\paragraph{Compromiso aceptado.} Cada cambio de esquema exige una migración escrita a mano y su esquema exportado.
Ese costo es el precio de no perder datos locales al actualizar la app (QA-14).

\paragraph{Riesgo.} Una migración defectuosa puede impedir abrir la base local con operaciones pendientes.
No corresponde a un riesgo del ATAM.

\paragraph{Dónde se ve.} Vista de implementación móvil: paquete \texttt{infraestructura/local}.
Vista física de la implementación: artefacto \texttt{danos.db} en el Teléfono Android.

\paragraph{Evidencia en el código.} \texttt{AppDatabase} declara la versión 3 con \texttt{MIGRACION\_1\_2} y \texttt{MIGRACION\_2\_3}, y la carpeta de esquemas guarda las versiones 1, 2 y 3.
Archivos de evidencia:
{\raggedright
\begin{itemize}
  \item \texttt{danos/\br{}app/\br{}src/\br{}main/\br{}java/\br{}com/\br{}example/\br{}danos/\br{}infraestructura/\br{}local/\br{}AppDatabase.kt}
  \item \texttt{danos/\br{}app/\br{}schemas}
\end{itemize}
\par}

%---------------------------------------------------------------------- AD-63
\subsection{AD-63. Sincronización del cliente móvil}
\label{sec:tec-ad63}

\paragraph{Contexto.} Esta decisión materializa la sincronización diferida (AD-31), el reenvío de la cola local (AD-44) y el envío idempotente del registro del desastre por el identificador del cliente (AD-48).
La sincronización atiende QA-14, cuya respuesta exige reenvío sin duplicados y 120 evaluaciones en $\le$ 10\,min.
También atiende QA-11, porque todo envío pasa por el gateway con el token OIDC.

\paragraph{Alternativas.} Se consideran WorkManager, SyncAdapter, un servicio en primer plano y Firebase.
El \cref{tab:tec-ad63} presenta la matriz de decisión ponderada.
La alternativa elegida obtiene 4,85 y la segunda mejor, SyncAdapter, obtiene 4,10.

\begingroup\small
\begin{longtable}{L{4.96cm} C{1.10cm} C{2.25cm} C{2.25cm} C{2.25cm} C{2.25cm}}
\cabecera{6}{tab:tec-ad63}{Matriz de decisión ponderada de AD-63}{\textbf{Criterio (driver)} & \textbf{Peso} & \hspace{0pt}\textbf{WorkManager} & \hspace{0pt}\textbf{SyncAdapter} & \hspace{0pt}\textbf{Servicio en primer plano} & \hspace{0pt}\textbf{Firebase}}
Ejecución garantizada con red disponible y reintento tras reinicio (QA-14, AD-44) & 30\,\% & 5 & 4 & 2 & 4 \\
Protocolo propio idempotente por identificador del cliente (QA-14, AD-48) & 25\,\% & 5 & 5 & 5 & 1 \\
Inmediatez del envío al recuperar la red (QA-14) & 15\,\% & 4 & 3 & 5 & 4 \\
Envío por el gateway con token OIDC (QA-11, AD-04) & 10\,\% & 5 & 5 & 5 & 1 \\
Experiencia del equipo & 10\,\% & 5 & 2 & 3 & 3 \\
Licencia y dependencia de un servicio externo & 10\,\% & 5 & 5 & 5 & 2 \\
\midrule
\textbf{Total ponderado} & 100\,\% & \textbf{4,85} & 4,10 & 3,90 & 2,65 \\
\end{longtable}
\endgroup

\paragraph{Alternativas descartadas.}
\begin{itemize}
  \item SyncAdapter pierde en ejecución garantizada (QA-14) y en experiencia del equipo, porque exige una cuenta del sistema y un proveedor de contenido.
  \item El servicio en primer plano gana en inmediatez del envío, pero pierde en ejecución garantizada tras un reinicio (QA-14).
  \item Firebase pierde en protocolo idempotente propio (QA-14) y en envío por el gateway (QA-11), porque sincroniza contra su propio servicio.
\end{itemize}

\paragraph{Compromiso aceptado.} El sistema operativo decide el instante de la ejecución de WorkManager.
El envío puede demorar unos segundos después de recuperar la red.

\paragraph{Riesgo.} Riesgo R-12 del ATAM: la variante \texttt{debug} admite tráfico sin TLS (L-07).
Riesgo R-13 del ATAM: desde un teléfono físico el emisor del token exige cambiar la configuración de Keycloak (L-08).

\paragraph{Dónde se ve.} Vista de implementación móvil: paquete \texttt{infraestructura/sync}.
Vista física de la implementación: \texttt{DanosCampo.apk} en el Teléfono Android.

\paragraph{Evidencia en el código.} QA-14 se verificó con 115 evaluaciones aplicadas, 5 en conflicto, 0 pérdidas, 0 sobrescrituras y 0 duplicados en 27,8\,s.
Archivos de evidencia:
{\raggedright
\begin{itemize}
  \item \texttt{danos/\br{}app/\br{}src/\br{}main/\br{}java/\br{}com/\br{}example/\br{}danos/\br{}infraestructura/\br{}sync/\br{}SincronizacionWorker.kt}
  \item \texttt{danos/\br{}app/\br{}src/\br{}main/\br{}java/\br{}com/\br{}example/\br{}danos/\br{}infraestructura/\br{}sync/\br{}EnviadorDesastres.kt}
\end{itemize}
\par}

%---------------------------------------------------------------------- AD-64
\subsection{AD-64. Biblioteca de resiliencia}
\label{sec:tec-ad64}

\paragraph{Contexto.} Esta decisión materializa el tiempo de espera, el reintento con espera creciente y el cortacircuitos de las fronteras (AD-05) y el cortacircuitos por adaptador (AD-37).
La biblioteca atiende QA-04, cuya respuesta exige cortar la propagación del fallo en $\le$ 60\,s.
También atiende QA-01, porque la caída de un sistema externo no debe degradar las demás funciones.

\paragraph{Alternativas.} Se consideran Microsoft.Extensions.Http.Resilience, una implementación propia y Steeltoe.
El \cref{tab:tec-ad64} presenta la matriz de decisión ponderada.
La alternativa elegida obtiene 4,90 y la segunda mejor, Implementación propia, obtiene 3,15.

\begingroup\small
\begin{longtable}{L{6.74cm} C{1.10cm} C{2.50cm} C{2.50cm} C{2.50cm}}
\cabecera{5}{tab:tec-ad64}{Matriz de decisión ponderada de AD-64}{\textbf{Criterio (driver)} & \textbf{Peso} & \hspace{0pt}\textbf{MS Resilience (Polly v8)} & \hspace{0pt}\textbf{Implementación propia} & \hspace{0pt}\textbf{Steeltoe}}
Tiempo de espera, reintento con variación aleatoria y cortacircuitos configurables (QA-01, AD-05) & 30\,\% & 5 & 3 & 2 \\
Corte de la propagación en $\le$ 60\,s verificable (QA-04) & 20\,\% & 5 & 3 & 3 \\
Integración con la fábrica de clientes HTTP de AD-50 & 15\,\% & 5 & 3 & 4 \\
Estado del circuito seguro ante concurrencia (QA-04) & 15\,\% & 5 & 2 & 4 \\
Experiencia del equipo & 10\,\% & 4 & 4 & 1 \\
Licencia y mantenimiento & 10\,\% & 5 & 5 & 4 \\
\midrule
\textbf{Total ponderado} & 100\,\% & \textbf{4,90} & 3,15 & 2,90 \\
\end{longtable}
\endgroup

\paragraph{Alternativas descartadas.}
\begin{itemize}
  \item La implementación propia iguala en experiencia y licencia, pero pierde en estado del circuito seguro ante concurrencia (QA-04).
  \item Steeltoe pierde en políticas configurables (QA-01) y en experiencia del equipo.
\end{itemize}

\paragraph{Compromiso aceptado.} Polly mide una proporción de fallos en una ventana, no un número de fallos consecutivos.
La regla de AD-05 se expresa como 3 llamadas mínimas con 100\,\% de fallos en 30\,s.

\paragraph{Riesgo.} Un reintento sobre una operación no idempotente duplicaría el efecto.
\texttt{Resiliencia.cs} reintenta solo operaciones idempotentes.
No corresponde a un riesgo del ATAM.

\paragraph{Dónde se ve.} Vista de implementación: proyecto \texttt{Plataforma.Comun.Adaptadores} (\texttt{Resiliencia.cs}).
Modelo de integración: cortacircuitos de cada adaptador.

\paragraph{Evidencia en el código.} La configuración usa 3\,s por intento, 2 reintentos, ventana de 30\,s y apertura de 15\,s. Las 12 pruebas de \texttt{AdaptadorSe6.Integracion} y las 14 de \texttt{AdaptadorSe7.Integracion} verifican el cortacircuitos y el reintento sin pérdida (QA-01).
Archivos de evidencia:
{\raggedright
\begin{itemize}
  \item \texttt{proyectoparte2/\br{}comun/\br{}Plataforma.Comun.Adaptadores/\br{}Resiliencia.cs}
  \item \texttt{proyectoparte2/\br{}Directory.Packages.props}
\end{itemize}
\par}

%---------------------------------------------------------------------- AD-65
\subsection{AD-65. Empaquetado, ejecución y escalado}
\label{sec:tec-ad65}

\paragraph{Contexto.} Esta decisión materializa las réplicas adicionales del camino de lectura (AD-18), la réplica activa con sondeo de salud (AD-24) y la ejecución híbrida con adaptadores efímeros (AD-25).
El empaquetado atiende QA-05, que exige escalar solo los componentes bajo presión sin modificar código.
También atiende QA-02 por el reemplazo de réplicas y QA-01 por el escalado de los adaptadores según su cola.

\paragraph{Alternativas.} Se consideran contenedores con Kubernetes, máquinas virtuales, Docker Swarm, Nomad y Knative.
El \cref{tab:tec-ad65} presenta la matriz de decisión ponderada.
La alternativa elegida obtiene 4,60 y la segunda mejor, Knative, obtiene 3,85.

\begingroup\small
\begin{longtable}{L{3.93cm} C{1.10cm} C{1.95cm} C{1.95cm} C{1.95cm} C{1.95cm} C{1.95cm}}
\cabecera{7}{tab:tec-ad65}{Matriz de decisión ponderada de AD-65}{\textbf{Criterio (driver)} & \textbf{Peso} & \hspace{0pt}\textbf{Docker + K8s (HPA, KEDA)} & \hspace{0pt}\textbf{Máquinas virtuales} & \hspace{0pt}\textbf{Docker Swarm} & \hspace{0pt}\textbf{Nomad} & \hspace{0pt}\textbf{Knative}}
Escalado automático selectivo sin modificar código (QA-05, AD-18) & 25\,\% & 5 & 2 & 2 & 4 & 5 \\
Escalado de 0 a N por largo de cola (QA-01, AD-25) & 15\,\% & 5 & 1 & 1 & 3 & 3 \\
Sondeo de salud y reemplazo de réplicas (QA-02, AD-24) & 20\,\% & 5 & 3 & 4 & 4 & 5 \\
Misma imagen en el computador y en el clúster (QA-09) & 15\,\% & 5 & 2 & 5 & 4 & 4 \\
Experiencia del equipo & 10\,\% & 4 & 3 & 3 & 1 & 1 \\
Costo de operación y licencia & 15\,\% & 3 & 2 & 5 & 3 & 3 \\
\midrule
\textbf{Total ponderado} & 100\,\% & \textbf{4,60} & 2,15 & 3,25 & 3,40 & 3,85 \\
\end{longtable}
\endgroup

\paragraph{Alternativas descartadas.}
\begin{itemize}
  \item Las máquinas virtuales pierden en escalado selectivo (QA-05) y en escalado por largo de cola (QA-01).
  \item Docker Swarm gana en costo de operación, pero pierde en escalado automático (QA-05), que no ofrece.
  \item Nomad pierde en escalado por cola y en licencia, porque cambió a BSL en 2023.
  \item Knative iguala en escalado selectivo y sondeo de salud, pero pierde en escalado por largo de cola (QA-01), porque está orientado a HTTP.
\end{itemize}

\paragraph{Compromiso aceptado.} En Docker Compose los adaptadores quedan siempre encendidos y no hay escalado automático.
El escalado de 0 a N solo existe en los manifiestos de Kubernetes (AD-25).

\paragraph{Riesgo.} El despliegue en Kubernetes no se verificó de punta a punta (L-05).
El escalado de 0 a N con KEDA no se midió (QA-01) y QA-05 no se midió con 10\,000 usuarios (L-03, L-04).

\paragraph{Dónde se ve.} Vista física de la implementación: nodo Docker Engine con los servicios del archivo de composición y anotación de Kubernetes con HPA y KEDA.

\paragraph{Evidencia en el código.} Los manifiestos declaran HPA de 2 a 6 réplicas para \texttt{consulta} y \texttt{consulta-\br{}danos}, de 2 a 5 para \texttt{consulta-\br{}programa}, y \texttt{ScaledObject} de KEDA de 0 a 4 réplicas para la salida de SE-6.
Archivos de evidencia:
{\raggedright
\begin{itemize}
  \item \texttt{proyectoparte2/\br{}despliegue/\br{}compose/\br{}docker-compose.yml}
  \item \texttt{proyectoparte2/\br{}despliegue/\br{}k8s/\br{}base/\br{}20-servicios.yaml}
  \item \texttt{proyectoparte2/\br{}despliegue/\br{}k8s/\br{}base/\br{}adaptador-se6.yaml}
\end{itemize}
\par}

%---------------------------------------------------------------------- AD-66
\subsection{AD-66. Plataforma del adaptador de SE-8}
\label{sec:tec-ad66}

\paragraph{Contexto.} CON-9, agregada el 19 de noviembre de 2026, exige que la solución final incluya Jakarta EE y .NET (E-44).
AD-03 ubica la heterogeneidad tecnológica (CON-2) en los adaptadores, no en los microservicios: un adaptador es el único punto de la arquitectura donde una pila distinta no toca el núcleo.
El adaptador de SE-8 estaba diseñado y no implementado (L-06).
QA-08 le exige una instancia encendida (AD-25), estado de entrega por alerta y \enquote{entregada} solo con acuse.
QA-09 pide que incorporar un sistema externo modifique 0 componentes existentes.
Esta decisión elige la plataforma de ese adaptador.

\paragraph{Alternativas.} Se consideran Jakarta EE 11 sobre WildFly 42.0.0.Final (imagen \texttt{quay.io/\br{}wildfly/\br{}wildfly:\br{}42.0.0.Final-\br{}jdk21}, Java 21), Spring Boot 3 sobre Java 21 y .NET 10 con \texttt{Plataforma.\br{}Comun.\br{}Adaptadores}, como los adaptadores de SE-6 y SE-7.
El \cref{tab:tec-ad66} presenta la matriz de decisión ponderada.
La alternativa elegida obtiene 4,35 y la segunda mejor, .NET 10, obtiene 4,00.

\begingroup\small
\begin{longtable}{L{6.0cm} C{1.10cm} C{2.6cm} C{2.6cm} C{2.6cm}}
\cabecera{5}{tab:tec-ad66}{Matriz de decisión ponderada de AD-66}{\textbf{Criterio (driver)} & \textbf{Peso} & \hspace{0pt}\textbf{Jakarta EE 11 (WildFly 42)} & \hspace{0pt}\textbf{Spring Boot 3} & \hspace{0pt}\textbf{.NET 10}}
Cumplimiento de CON-9 & 25\,\% & 5 & 2 & 1 \\
Adaptador nuevo sin modificar componentes existentes (QA-09, AD-03) & 20\,\% & 5 & 5 & 5 \\
Cortacircuitos del Cuadro 10 con aviso de apertura para EV-23 (QA-01, QA-08, AD-05, AD-37) & 15\,\% & 3 & 4 & 5 \\
Consumo transaccional: ack después del commit y DLQ a la quinta entrega (AD-02, AD-45) & 15\,\% & 4 & 4 & 5 \\
Instancia siempre encendida: memoria y arranque (QA-08, AD-25) & 10\,\% & 4 & 4 & 5 \\
Experiencia del equipo (especificaciones vistas en el curso) & 10\,\% & 4 & 3 & 5 \\
Licencia y costo & 5\,\% & 5 & 5 & 5 \\
\midrule
\textbf{Total ponderado} & 100\,\% & \textbf{4,35} & 3,65 & 4,00 \\
\end{longtable}
\endgroup

\paragraph{Alternativas descartadas.}
\begin{itemize}
  \item .NET 10 gana en resiliencia, mensajería y arranque, porque reutiliza \texttt{Plataforma.\br{}Comun.\br{}Adaptadores} y \texttt{ConsumidorIdempotente}. Pierde en el criterio que origina la decisión: no cumple CON-9.
  \item Spring Boot 3 corre sobre Java y usa algunas API de Jakarta, pero no es la plataforma Jakarta EE: cumple CON-9 solo en parte. Tiene un cortacircuitos con ventana de tiempo (Resilience4j) y es una tercera pila de marco que el equipo no vio en el curso.
\end{itemize}

\paragraph{Compromiso aceptado.} El adaptador usa las especificaciones estándar: Jakarta Messaging, EJB (\texttt{@Stateless}, \texttt{@Singleton}, Timer Service), Jakarta Persistence, Jakarta REST, CDI, Jakarta Concurrency y MicroProfile Config.
Dos piezas son propias y no vienen del estándar:
\begin{itemize}
  \item \textbf{Cortacircuitos propio en lugar de MicroProfile Fault Tolerance.} El \texttt{@CircuitBreaker} de MicroProfile mide una ventana de N llamadas, no la ventana de \qty{30}{\second} del Cuadro 10, y no avisa cuando el circuito se abre o se cierra. Ese aviso hace falta para publicar EV-23 por la bandeja (E-15). La clase \texttt{Cortacircuitos} aplica la regla de Polly (al menos 3 llamadas en \qty{30}{\second} con \qty{100}{\percent} de fallo, \qty{15}{\second} abierto y una sola llamada de prueba) y es segura ante concurrencia. Tiene 9 pruebas unitarias propias, una de ellas con 200 hilos.
  \item \textbf{EJB \texttt{@Singleton @Startup} en lugar de MDB.} Un MDB necesita un resource adapter JCA para el proveedor de mensajería, y RabbitMQ no tiene uno mantenido. El consumo usa igual la API estándar de Jakarta Messaging 3.1 (\texttt{Connection}, \texttt{Session} transaccional, \texttt{MessageConsumer}) con el cliente JMS de RabbitMQ en modo AMQP sobre la cola existente. El hilo lo da un \texttt{ManagedExecutorService}. Con \texttt{nackOnRollback}, \texttt{rollback()} manda el mensaje a la DLQ y cerrar la sesión lo reencola.
\end{itemize}
Se acepta además una segunda pila que mantener en la frontera. Esa pila duplica en Java la lectura del sobre y la parte del catálogo que usa el adaptador. Los esquemas JSON siguen siendo una sola fuente: el WAR los copia de \texttt{Plataforma.\br{}Comun.\br{}Eventos} al compilar.

\paragraph{Riesgo.} Un cambio del sobre canónico o de las claves de ruteo obliga a cambiar dos implementaciones. Lo detecta la validación contra los mismos esquemas, en las pruebas y al consumir. Corresponde al riesgo R-15 del ATAM.
Mientras Notificación no exista, el adaptador recibe EV-03 y EV-04 como sustituto temporal (L-11).

\paragraph{Dónde se ve.} Vista lógica: componente Adaptador de SE-8 [\texttt{adaptador-se8}] «implementado».
Vista de procesos: proceso \texttt{adaptador-se8} ×1, siempre encendido.
Vista física de la implementación: entorno de ejecución WildFly 42 dentro del contenedor \texttt{adaptador-se8} (\cref{fig:vista-despliegue-impl}, \cref{tab:nodos-impl}).
Vista de implementación: paquete \texttt{adaptador-\br{}se8-\br{}jee} (\cref{fig:vista-implementacion}).
Modelo de integración: fila de SE-8 del \cref{tab:integracion}.

\paragraph{Evidencia en el código.} 52 pruebas unitarias y 7 de integración en verde. Las de integración usan Testcontainers con la topología real del bus (cola quorum con \texttt{delivery-limit} 5) y la base del migrador. Demuestran el ack después del commit JPA, la DLQ a la quinta entrega fallida, el duplicado sin segundo envío y el circuito con EV-23.
La prueba \texttt{prueba-se8.sh} mide contra el compose: \qty{1,8}{\second} del registro a la alerta entregada con acuse, circuito abierto a los \qty{25,6}{\second} de la caída, y cola en 0 a los \qty{12,9}{\second} de la vuelta, con las 20 alertas entregadas una sola vez.
El \cref{tab:que-se-modifica} registra los 0 componentes existentes modificados.
Archivos de evidencia:
{\raggedright
\begin{itemize}
  \item \texttt{proyectoparte2/\br{}adaptadores/\br{}adaptador-se8-jee/\br{}adaptador/\br{}src/\br{}main/\br{}java/\br{}plataforma/\br{}adaptador/\br{}se8/\br{}jee/\br{}AdaptadorSe8.java}
  \item \texttt{proyectoparte2/\br{}adaptadores/\br{}adaptador-se8-jee/\br{}adaptador/\br{}src/\br{}main/\br{}java/\br{}plataforma/\br{}adaptador/\br{}se8/\br{}consumo/\br{}BucleConsumo.java}
  \item \texttt{proyectoparte2/\br{}adaptadores/\br{}adaptador-se8-jee/\br{}adaptador/\br{}src/\br{}main/\br{}java/\br{}plataforma/\br{}adaptador/\br{}se8/\br{}resiliencia/\br{}Cortacircuitos.java}
  \item \texttt{proyectoparte2/\br{}adaptadores/\br{}adaptador-se8-jee/\br{}Dockerfile}
  \item \texttt{proyectoparte2/\br{}pruebas/\br{}humo/\br{}prueba-se8.sh}
\end{itemize}
\par}

%------------------------------------------------------------------------------
\subsection{Restricciones de las tecnologías elegidas}
\label{sec:tec-restricciones}

Las tecnologías elegidas imponen límites al entorno implementado.
El \cref{tab:tec-restricciones} los relaciona con la decisión que los origina y con el escenario afectado.
Ninguno modifica una decisión estructural.

\begingroup\small
\begin{longtable}{L{2.6cm} L{1.3cm} L{6.4cm} L{1.7cm} L{3.3cm}}
\cabecera{5}{tab:tec-restricciones}{Restricciones que imponen las tecnologías elegidas}{\textbf{Restricción} &
\textbf{Origen} & \textbf{Efecto} & \textbf{QA afectado} & \textbf{Qué haría falta}}
Variante \texttt{debug} y Keycloak sin TLS (L-07) & AD-58, AD-61, AD-63 &
La app de entrega admite \texttt{http://}. Keycloak corre sin TLS. El tráfico entre el teléfono, el gateway y el proveedor de identidad viaja sin cifrar. Corresponde al riesgo R-12 del ATAM. &
QA-11, QA-12 & TLS en el borde y URL HTTPS en la variante \texttt{release} \\
Emisor del token en \texttt{localhost} (L-08) & AD-57, AD-58 &
El \texttt{iss} del token es \texttt{http://localhost:8180}. El emulador exige \texttt{adb reverse}. Un teléfono físico exige cambiar \texttt{KC\_HOSTNAME} e \texttt{Identidad\_\_Emisor} a la vez. Corresponde al riesgo R-13 del ATAM. &
QA-11, QA-14 & Un nombre DNS común para Keycloak \\
Imagen comunitaria de MinIO (DT-11) & AD-56 &
El proyecto original dejó de publicar imágenes. El archivo de composición usa el fork \texttt{pgsty/minio}, cuya continuidad no está garantizada. &
QA-02, QA-14 & Sustituir por otro almacén con API S3 sin cambiar el código \\
Política de uso de las teselas de OpenStreetMap & AD-60 &
Las teselas públicas no tienen acuerdo de servicio, exigen atribución y prohíben el uso intensivo. Durante una emergencia con 10\,000 usuarios la carga del mapa dependería de un tercero. &
QA-05, QA-06 & Servidor de teselas propio o proveedor con acuerdo \\
Keycloak en \texttt{start-dev} & AD-58 &
El modo de desarrollo usa una base embebida, no exige TLS y no está pensado para producción. El realm se importa desde \texttt{realm-plataforma.json} en cada arranque. Es una sola instancia. &
QA-02, QA-11 & Modo \texttt{start} con base PostgreSQL, TLS y dos instancias \\
Sustituto temporal de Notificación (L-11) & AD-66 &
EV-03 y EV-04 llegan a \texttt{q.salida.se8} por dos enlaces temporales, y el adaptador alerta a un único destino de difusión. No hay suscripciones ni resolución por ámbito, y el estado de entrega se consulta en el adaptador. &
QA-08 & Implementar Notificación (AD-13), quitar los enlaces y poner \texttt{Sustituto\_\_Activo=false} \\
Dos pilas en la frontera & AD-66 &
El adaptador de SE-8 duplica en Java la lectura del sobre y la parte del catálogo que usa. Un cambio del contrato se hace en .NET y en Java. Corresponde al riesgo R-15 del ATAM. &
QA-09, QA-10 & Mantener los esquemas JSON como única fuente y validarlos en las pruebas de las dos pilas \\
\end{longtable}
\endgroup
.
%  Generado a partir de la especificación de la versión 1.2. No carga paquetes.
%==============================================================================
\section{Decisiones de tecnología}
\label{sec:tecnologia}

El enunciado del proyecto (§8) pide seleccionar las tecnologías en la segunda entrega.
CON-8 exigía independencia tecnológica en la primera entrega.
Esta sección sustituye esa condición para la segunda entrega.
El Anexo~\ref{sec:anexo-a} deja de ser una excepción a CON-8 y describe la implementación de FP-06 con las tecnologías elegidas aquí.

Las decisiones se numeran de AD-50 a AD-66.
AD-66 responde a CON-9, que se agregó después de AD-50 a AD-65.
Ninguna modifica una decisión estructural.
Cada una materializa uno o varios mecanismos de AD-01 a AD-49, y su justificación parte de los escenarios QA-nn que ese mecanismo debe cumplir (CON-7).

\subsection{Método de decisión}
\label{sec:tec-metodo}

Cada decisión sigue cinco pasos:
\begin{enumerate}
  \item Se identifica el mecanismo del SAD que la tecnología materializa y los escenarios QA-nn que ese mecanismo atiende.
  \item Se eligen entre tres y cinco alternativas con uso real para ese mecanismo.
  \item Se definen los criterios. Cada criterio cita el QA-nn, AD-nn o CON-n del que proviene. La experiencia del equipo y la licencia o el costo se incluyen con un peso menor que los criterios de calidad.
  \item Se asigna un peso a cada criterio. Los pesos de una decisión suman 100\,\%.
  \item Se califica cada alternativa en cada criterio y se calcula el total ponderado.
\end{enumerate}

La escala de calificación tiene cinco niveles:
\begin{itemize}
  \item 1: la alternativa no satisface el criterio o lo impide.
  \item 2: lo satisface con desarrollo propio considerable.
  \item 3: lo satisface con configuración adicional o con complementos de terceros.
  \item 4: lo satisface de forma nativa con limitaciones menores.
  \item 5: lo satisface de forma nativa y la implementación lo verifica.
\end{itemize}

El total ponderado de una alternativa es $\sum_i p_i \, c_i / 100$, donde $p_i$ es el peso en porcentaje y $c_i$ la calificación del criterio $i$.
El máximo es 5.
Las cifras de verificación citadas provienen del perfil local de un computador, con 20 clientes, y no reemplazan la prueba con 10\,000 usuarios (L-04).

\subsection{Resumen de las decisiones}
\label{sec:tec-resumen}

El \cref{tab:tecnologia-resumen} relaciona cada decisión con su tecnología, el mecanismo que materializa y los escenarios que favorece.
Las versiones son las de los archivos de proyecto.

\begingroup\small
\begin{longtable}{L{1.15cm} L{2.2cm} L{6.0cm} L{2.6cm} L{2.6cm}}
\cabecera{5}{tab:tecnologia-resumen}{Resumen de las decisiones de tecnología}{\textbf{ID} & \textbf{Categoría} &
\textbf{Tecnología y versión} & \textbf{Mecanismo} & \textbf{QA que favorece}}
AD-50 & Lenguaje y plataforma & C\# sobre .NET 10 (SDK 10.0.100, \texttt{net10.0}) & AD-01, AD-08, AD-09, AD-45 & QA-06, QA-07, QA-14 \\
AD-51 & Marco web & ASP.NET Core 10 (controladores y \emph{minimal APIs}) & AD-08, AD-22, AD-24 & QA-06, QA-07, QA-11 \\
AD-52 & Persistencia & PostgreSQL 18 con PostGIS 3.6 (\texttt{postgis/postgis:18-3.6}), una base por microservicio y réplica \texttt{db-replica} & AD-01, AD-17, AD-20, AD-24, AD-41, AD-45 & QA-03, QA-04, QA-06, QA-13, QA-15 \\
AD-53 & Acceso a datos & Entity Framework Core 10.0.14, Npgsql 10.0.3 y NetTopologySuite 10.0.3, con contenedor migrador & AD-08, AD-17, AD-41, AD-45 & QA-06, QA-10, QA-13 \\
AD-54 & Mensajería & RabbitMQ 4.4 (3 nodos, colas quorum acotadas, DLQ) con RabbitMQ.Client 7.2.2 & AD-02, AD-05, AD-36, AD-45 & QA-01, QA-04, QA-08 \\
AD-55 & Caché & Redis 8.8 (\texttt{redis:8.8.4-alpine}) con StackExchange.Redis 3.4.29 & AD-16, AD-17 & QA-05, QA-06 \\
AD-56 & Almacén de objetos & MinIO con API S3 (paquete Minio 7.0.0) & AD-24, AD-31 & QA-02, QA-03, QA-14 \\
AD-57 & Puerta de entrada & YARP 2.3.0 con catálogo \texttt{rutas.json} (56 rutas) & AD-04, AD-19, AD-22 & QA-05, QA-06, QA-11 \\
AD-58 & Identidad & Keycloak 26.5.2 (realm \texttt{plataforma}), AppAuth-Android 0.11.1, angular-auth-oidc-client 22.0.2, cliente OpenID Connect de ASP.NET Core para el tablero y validador JWT propio en \texttt{Plataforma.Comun.Seguridad} & AD-04, AD-22, AD-49 & QA-11, QA-14 \\
AD-59 & Cliente web & Angular 22.2 (TypeScript 6.0) servido por nginx 1.29 detrás del gateway & AD-04, AD-08 & QA-06, QA-10, QA-11 \\
AD-60 & Mapas & Leaflet 1.9.4 con teselas de OpenStreetMap & AD-16, AD-17 & QA-06 \\
AD-61 & Cliente móvil & Android nativo: Kotlin 2.0.21, Jetpack Compose (BOM 2024.09.00), Hilt 2.60.1, minSdk 26, targetSdk 36 & AD-31, AD-44, AD-48 & QA-07, QA-14 \\
AD-62 & Persistencia móvil & Room 2.8.5 sobre SQLite, migraciones explícitas 1$\to$2$\to$3 & AD-31, AD-44, AD-48 & QA-14 \\
AD-63 & Sincronización móvil & WorkManager 2.12.0, Retrofit 3.0.0, OkHttp 5.4.0 y kotlinx.serialization 1.9.0 & AD-31, AD-44, AD-48 & QA-11, QA-14 \\
AD-64 & Resiliencia & Microsoft.Extensions.Http.Resilience 10.11.0 (Polly v8) & AD-05, AD-37 & QA-01, QA-04 \\
AD-65 & Despliegue & Contenedores Docker (\texttt{mcr.microsoft.com/\br{}dotnet/\br{}aspnet:10.0}), Docker Compose para un computador y Kubernetes con HPA y KEDA para escalar & AD-18, AD-24, AD-25 & QA-01, QA-02, QA-05 \\
AD-66 & Adaptador de SE-8 & Jakarta EE 11 sobre WildFly 42.0.0.Final (Java 21): Jakarta Messaging con el cliente JMS de RabbitMQ 3.9.1, EJB, Jakarta Persistence, Jakarta REST y MicroProfile Config & AD-03, AD-05, AD-09, AD-25, AD-45 & QA-08, QA-09, CON-9 \\
\end{longtable}
\endgroup

%---------------------------------------------------------------------- AD-50
\subsection{AD-50. Lenguaje y plataforma de los microservicios}
\label{sec:tec-ad50}

\paragraph{Contexto.} Esta decisión materializa los microservicios con almacén propio (AD-01), las cuatro capas de cada microservicio (AD-08), el contrato común de los adaptadores (AD-09) y la bandeja de salida transaccional (AD-45).
El lenguaje determina la latencia del camino síncrono, que exige QA-06 (p95 $<$ 2\,s) y QA-07 (p95 $<$ 3\,s).
También determina si la base común \texttt{comun} y el prototipo de FP-06 del Anexo~\ref{sec:anexo-a} se reutilizan, lo que afecta la verificación de QA-14.
CON-2 exige integrarse con sistemas de tecnologías diferentes, y esa heterogeneidad la absorben los adaptadores (AD-03), no el lenguaje de los microservicios.

\paragraph{Alternativas.} Se consideran C\# sobre .NET 10, Java 21, Node.js 22 con TypeScript y Go.
Las cuatro tienen licencia libre y ejecución en contenedores.
El \cref{tab:tec-ad50} presenta la matriz de decisión ponderada.
La alternativa elegida obtiene 4,90 y la segunda mejor, Java 21, obtiene 3,85.

\begingroup\small
\begin{longtable}{L{4.96cm} C{1.10cm} C{2.25cm} C{2.25cm} C{2.25cm} C{2.25cm}}
\cabecera{6}{tab:tec-ad50}{Matriz de decisión ponderada de AD-50}{\textbf{Criterio (driver)} & \textbf{Peso} & \hspace{0pt}\textbf{.NET 10} & \hspace{0pt}\textbf{Java 21} & \hspace{0pt}\textbf{Node.js 22} & \hspace{0pt}\textbf{Go}}
Rendimiento del camino síncrono (QA-06, QA-07) & 25\,\% & 5 & 4 & 3 & 5 \\
Bibliotecas maduras para bus, resiliencia, datos espaciales y JWT (AD-02, AD-05, AD-17, AD-22) & 25\,\% & 5 & 5 & 3 & 3 \\
Reutilización de \texttt{comun} y del prototipo de FP-06 (QA-14, AD-09) & 20\,\% & 5 & 2 & 2 & 2 \\
Memoria y arranque de cada réplica (QA-05, AD-18) & 10\,\% & 4 & 3 & 4 & 5 \\
Experiencia del equipo & 10\,\% & 5 & 4 & 3 & 1 \\
Licencia y costo & 10\,\% & 5 & 5 & 5 & 5 \\
\midrule
\textbf{Total ponderado} & 100\,\% & \textbf{4,90} & 3,85 & 3,10 & 3,50 \\
\end{longtable}
\endgroup

\paragraph{Alternativas descartadas.}
\begin{itemize}
  \item Java 21 pierde en la reutilización de \texttt{comun} y del prototipo de FP-06 (QA-14), que habría que reescribir.
  \item Node.js 22 pierde en el rendimiento del camino síncrono (QA-06, QA-07) y en la madurez de las bibliotecas para datos espaciales y bus.
  \item Go gana en memoria y arranque de cada réplica (QA-05), pero pierde en bibliotecas para los mecanismos del SAD y en experiencia del equipo.
\end{itemize}

\paragraph{Compromiso aceptado.} Una sola pila reduce el costo de mantener \texttt{comun} (AD-09) y acepta que todos los microservicios dependen del ciclo de soporte de .NET.
La heterogeneidad que exige CON-2 queda confinada en los adaptadores (AD-03).
El núcleo sigue en .NET aun con CON-9: Jakarta EE se confina a la frontera según AD-03 y AD-66.

\paragraph{Riesgo.} El fin del soporte de .NET 10 obligaría a migrar todos los microservicios a la vez.
El archivo \texttt{global.json} fija el SDK con avance solo de la versión de características, lo que acota este riesgo.
No corresponde a un riesgo del ATAM.

\paragraph{Dónde se ve.} Vista de implementación: proyectos \texttt{*.Api}, \texttt{*.Aplicacion}, \texttt{*.Dominio} e \texttt{*.Infraestructura} de cada microservicio.
Vista física de la implementación: artefactos \texttt{plataforma/<servicio>:local} sobre \texttt{mcr.microsoft.com/\br{}dotnet/\br{}aspnet:10.0}.

\paragraph{Evidencia en el código.} La suite de .NET suma 1\,008 pruebas en 19 proyectos en verde. QA-07 se midió con p95 de 1\,334\,ms en el perfil local (L-04).
Archivos de evidencia:
{\raggedright
\begin{itemize}
  \item \texttt{proyectoparte2/\br{}global.json}
  \item \texttt{proyectoparte2/\br{}Directory.Build.props}
\end{itemize}
\par}

%---------------------------------------------------------------------- AD-51
\subsection{AD-51. Marco de las API de los microservicios}
\label{sec:tec-ad51}

\paragraph{Contexto.} Esta decisión materializa la capa de acceso de cada microservicio (AD-08), el nivel 2 de la autorización (AD-22) y el sondeo de salud de la réplica activa (AD-24).
El marco atiende la latencia de QA-06 y QA-07 y el rechazo del 100\,\% de las modificaciones no autorizadas de QA-11.
La elección está condicionada por AD-50, porque el marco debe pertenecer a la pila elegida.

\paragraph{Alternativas.} Se consideran ASP.NET Core 10, Spring Boot 3 y NestJS.
Spring Boot 3 y NestJS se evalúan como el marco que se habría usado con Java y con Node.js.
El \cref{tab:tec-ad51} presenta la matriz de decisión ponderada.
La alternativa elegida obtiene 4,85 y la segunda mejor, Spring Boot 3, obtiene 3,75.

\begingroup\small
\begin{longtable}{L{6.74cm} C{1.10cm} C{2.50cm} C{2.50cm} C{2.50cm}}
\cabecera{5}{tab:tec-ad51}{Matriz de decisión ponderada de AD-51}{\textbf{Criterio (driver)} & \textbf{Peso} & \hspace{0pt}\textbf{ASP.NET Core 10} & \hspace{0pt}\textbf{Spring Boot 3} & \hspace{0pt}\textbf{NestJS}}
Latencia por petición (QA-06, QA-07) & 25\,\% & 5 & 4 & 3 \\
Validación JWT y autorización de nivel 2 en la tubería (QA-11, AD-22) & 20\,\% & 5 & 5 & 4 \\
Sondeos de salud vivo y listo (QA-02, QA-04, AD-24) & 15\,\% & 4 & 5 & 4 \\
Coherencia con AD-50 y con el gateway (QA-10, AD-57) & 20\,\% & 5 & 1 & 1 \\
Experiencia del equipo & 10\,\% & 5 & 3 & 3 \\
Licencia y costo & 10\,\% & 5 & 5 & 5 \\
\midrule
\textbf{Total ponderado} & 100\,\% & \textbf{4,85} & 3,75 & 3,15 \\
\end{longtable}
\endgroup

\paragraph{Alternativas descartadas.}
\begin{itemize}
  \item Spring Boot 3 gana en sondeos de salud con Actuator (QA-02), pero pierde en coherencia con AD-50, porque introduce una segunda pila.
  \item NestJS pierde en latencia por petición (QA-06, QA-07) y en coherencia con AD-50.
\end{itemize}

\paragraph{Compromiso aceptado.} Los sondeos de salud se componen con \texttt{Plataforma.Comun.Salud} en lugar de un módulo integrado como Actuator (AD-24).

\paragraph{Riesgo.} Una actualización mayor del marco afecta a todos los microservicios a la vez.
No corresponde a un riesgo del ATAM.

\paragraph{Dónde se ve.} Vista de implementación: proyectos \texttt{*.Api} con \texttt{Sdk.Web}.
Vista física de la implementación: artefactos \texttt{plataforma/<servicio>:local}.

\paragraph{Evidencia en el código.} Los proyectos \texttt{*.Api.csproj} usan \texttt{Microsoft.NET.Sdk.Web}. QA-07 se cumple con 100 de 100 registros respondidos con 201.
Archivos de evidencia:
{\raggedright
\begin{itemize}
  \item \texttt{proyectoparte2/\br{}servicios/\br{}gestion-desastres/\br{}src/\br{}GestionDesastres.Api}
  \item \texttt{proyectoparte2/\br{}comun/\br{}Plataforma.Comun.Salud}
\end{itemize}
\par}

%---------------------------------------------------------------------- AD-52
\subsection{AD-52. Motor de base de datos}
\label{sec:tec-ad52}

\paragraph{Contexto.} Esta decisión materializa el almacén propio de cada microservicio (AD-01), la geometría precalculada de las zonas (AD-17), los historiales versionados y de solo inserción (AD-20, AD-41), la réplica activa (AD-24) y la bandeja de salida (AD-45).
El motor atiende la consulta espacial de QA-06, la publicación sin pérdida de QA-04, la inmutabilidad de QA-13 y QA-15 y el RPO $\le$ 60\,s de QA-03.

\paragraph{Alternativas.} Se consideran PostgreSQL con PostGIS, MySQL 8, SQL Server y MongoDB.
El \cref{tab:tec-ad52} presenta la matriz de decisión ponderada.
La alternativa elegida obtiene 4,85 y la segunda mejor, SQL Server, obtiene 4,10.

\begingroup\small
\begin{longtable}{L{4.96cm} C{1.10cm} C{2.25cm} C{2.25cm} C{2.25cm} C{2.25cm}}
\cabecera{6}{tab:tec-ad52}{Matriz de decisión ponderada de AD-52}{\textbf{Criterio (driver)} & \textbf{Peso} & \hspace{0pt}\textbf{PostgreSQL + PostGIS} & \hspace{0pt}\textbf{MySQL 8} & \hspace{0pt}\textbf{SQL Server} & \hspace{0pt}\textbf{MongoDB}}
Consulta espacial con índice (QA-06, AD-17) & 25\,\% & 5 & 3 & 4 & 3 \\
Transacción ACID para la bandeja y el historial (QA-04, AD-20, AD-45) & 20\,\% & 5 & 4 & 5 & 3 \\
Inmutabilidad garantizada por permisos del almacén (QA-13, QA-15, AD-41) & 15\,\% & 5 & 4 & 5 & 3 \\
Réplica continua para la réplica activa y el RPO (QA-03, AD-24, AD-27) & 15\,\% & 4 & 4 & 5 & 5 \\
Experiencia del equipo & 10\,\% & 5 & 4 & 3 & 3 \\
Licencia y costo & 15\,\% & 5 & 4 & 2 & 3 \\
\midrule
\textbf{Total ponderado} & 100\,\% & \textbf{4,85} & 3,75 & 4,10 & 3,30 \\
\end{longtable}
\endgroup

\paragraph{Alternativas descartadas.}
\begin{itemize}
  \item MySQL 8 pierde en consulta espacial con índice (QA-06), porque sus funciones espaciales son menos completas que PostGIS.
  \item SQL Server gana en réplica continua (QA-03), pero pierde en licencia y costo, porque su licencia comercial se cobra por núcleo.
  \item MongoDB gana en réplica continua (QA-03), pero pierde en transacción ACID para la bandeja (QA-04) y en inmutabilidad por permisos (QA-13).
\end{itemize}

\paragraph{Compromiso aceptado.} En el entorno de un computador las nueve bases comparten la instancia \texttt{db}.
La separación por microservicio se conserva por base y por rol con mínimo privilegio (AD-01), no por servidor.

\paragraph{Riesgo.} La instancia \texttt{db} es un punto único de fallo en el entorno de un computador.
La réplica \texttt{db-replica} reduce la pérdida de datos, pero la promoción no es automática y no hay segunda región (L-01).
QA-02 y QA-03 no se midieron.

\paragraph{Dónde se ve.} Vista física de la implementación: nodos \texttt{db} y \texttt{db-replica}.
Vista de implementación: paquete \texttt{Infraestructura/Persistencia} de cada microservicio y migradores.

\paragraph{Evidencia en el código.} Los servicios \texttt{db} y \texttt{db-replica} usan la imagen \texttt{postgis/postgis:18-3.6}. QA-06 se midió con p95 de 34\,ms en el consolidado del desastre en el perfil local (L-04).
Archivos de evidencia:
{\raggedright
\begin{itemize}
  \item \texttt{proyectoparte2/\br{}despliegue/\br{}compose/\br{}docker-compose.yml}
\end{itemize}
\par}

%---------------------------------------------------------------------- AD-53
\subsection{AD-53. Acceso a datos y migraciones}
\label{sec:tec-ad53}

\paragraph{Contexto.} Esta decisión materializa la capa de infraestructura de cada microservicio (AD-08), el mapeo de las geometrías (AD-17), los permisos del historial de solo inserción (AD-41) y la escritura atómica con la bandeja (AD-45).
El acceso a datos afecta la latencia de QA-06, la evolución del esquema que pide QA-10 y la inmutabilidad de QA-13.

\paragraph{Alternativas.} Se consideran Entity Framework Core, Dapper, NHibernate y SQL directo con Npgsql sin mapeador.
El \cref{tab:tec-ad53} presenta la matriz de decisión ponderada.
La alternativa elegida obtiene 4,80 y la segunda mejor, Dapper, obtiene 3,60.

\begingroup\small
\begin{longtable}{L{4.96cm} C{1.10cm} C{2.25cm} C{2.25cm} C{2.25cm} C{2.25cm}}
\cabecera{6}{tab:tec-ad53}{Matriz de decisión ponderada de AD-53}{\textbf{Criterio (driver)} & \textbf{Peso} & \hspace{0pt}\textbf{EF Core 10} & \hspace{0pt}\textbf{Dapper} & \hspace{0pt}\textbf{NHibernate} & \hspace{0pt}\textbf{SQL directo}}
Migraciones versionadas con los permisos del historial (QA-10, QA-13, AD-41) & 25\,\% & 5 & 2 & 3 & 1 \\
Tipos espaciales mapeados (QA-06, AD-17) & 20\,\% & 5 & 3 & 2 & 3 \\
Control del SQL en las lecturas críticas (QA-06) & 20\,\% & 4 & 5 & 3 & 5 \\
Aislamiento en la capa de infraestructura y pruebas (QA-10, AD-08) & 15\,\% & 5 & 4 & 4 & 3 \\
Experiencia del equipo & 10\,\% & 5 & 4 & 2 & 4 \\
Licencia y costo & 10\,\% & 5 & 5 & 5 & 5 \\
\midrule
\textbf{Total ponderado} & 100\,\% & \textbf{4,80} & 3,60 & 3,05 & 3,20 \\
\end{longtable}
\endgroup

\paragraph{Alternativas descartadas.}
\begin{itemize}
  \item Dapper gana en control del SQL (QA-06), pero pierde en migraciones versionadas (QA-10, QA-13), que exigirían otra herramienta.
  \item NHibernate pierde en tipos espaciales mapeados (QA-06) y en experiencia del equipo.
  \item SQL directo gana en control del SQL (QA-06), pero pierde en migraciones versionadas y en aislamiento de la capa de infraestructura (QA-10).
\end{itemize}

\paragraph{Compromiso aceptado.} El mapeador agrega una capa sobre el SQL.
Las lecturas críticas del modelo de lectura (AD-15) pueden usar SQL explícito dentro del mismo contexto cuando QA-06 lo exija.
El esquema lo aplica un contenedor migrador con el rol dueño, y el microservicio corre con un rol sin permisos de DDL (AD-41).

\paragraph{Riesgo.} Una consulta generada por el mapeador puede degradar la latencia del mapa de daños, que se midió con p95 de 1\,628\,ms frente al límite de 2\,s de QA-06.
No corresponde a un riesgo del ATAM.

\paragraph{Dónde se ve.} Vista de implementación: paquete \texttt{Infraestructura/Persistencia} y migradores (\texttt{*-migrador}).
Vista física de la implementación: nodos \texttt{db} y \texttt{db-replica}.

\paragraph{Evidencia en el código.} El servicio \texttt{reconstruccion-migrador} del archivo de composición crea los roles y aplica las migraciones antes de arrancar el microservicio.
Archivos de evidencia:
{\raggedright
\begin{itemize}
  \item \texttt{proyectoparte2/\br{}Directory.Packages.props}
  \item \texttt{proyectoparte2/\br{}comun/\br{}Plataforma.Comun.Migraciones}
  \item \texttt{proyectoparte2/\br{}servicios/\br{}evaluacion-danos/\br{}src/\br{}EvaluacionDanos.Infraestructura/\br{}Persistencia}
\end{itemize}
\par}

%---------------------------------------------------------------------- AD-54
\subsection{AD-54. Bus de eventos}
\label{sec:tec-ad54}

\paragraph{Contexto.} Esta decisión materializa el bus de eventos con colas por consumidor (AD-02), las colas acotadas de la frontera (AD-05), la cola de salida con reintento (AD-36) y la entrega desde la bandeja de salida (AD-45).
El bus atiende QA-04, cuya respuesta exige 0 eventos perdidos ante el fallo de un componente interno, y QA-01, que exige procesar el 100\,\% de lo encolado cuando el sistema externo vuelve.

\paragraph{Alternativas.} Se consideran RabbitMQ, Apache Kafka, NATS JetStream y MassTransit sobre RabbitMQ.
MassTransit se evalúa como capa de abstracción sobre el mismo intermediario.
El \cref{tab:tec-ad54} presenta la matriz de decisión ponderada.
La alternativa elegida obtiene 4,90 y la segunda mejor, NATS JetStream, obtiene 3,85.

\begingroup\small
\begin{longtable}{L{4.96cm} C{1.10cm} C{2.25cm} C{2.25cm} C{2.25cm} C{2.25cm}}
\cabecera{6}{tab:tec-ad54}{Matriz de decisión ponderada de AD-54}{\textbf{Criterio (driver)} & \textbf{Peso} & \hspace{0pt}\textbf{RabbitMQ 4.4} & \hspace{0pt}\textbf{Kafka} & \hspace{0pt}\textbf{NATS JetStream} & \hspace{0pt}\textbf{MassTransit}}
Cola acotada por consumidor con rechazo al productor (QA-04, AD-05) & 25\,\% & 5 & 2 & 4 & 4 \\
Confirmación por mensaje, límite de entregas y DLQ (QA-01, QA-04, AD-36) & 20\,\% & 5 & 2 & 4 & 5 \\
Replicación de la cola ante la caída de un nodo (QA-02, QA-04) & 15\,\% & 5 & 5 & 5 & 5 \\
Topología explícita y declarada al desplegar (QA-09, QA-10) & 15\,\% & 5 & 3 & 3 & 2 \\
Experiencia del equipo & 10\,\% & 4 & 2 & 1 & 3 \\
Licencia y costo & 15\,\% & 5 & 5 & 5 & 2 \\
\midrule
\textbf{Total ponderado} & 100\,\% & \textbf{4,90} & 3,05 & 3,85 & 3,65 \\
\end{longtable}
\endgroup

\paragraph{Alternativas descartadas.}
\begin{itemize}
  \item Kafka iguala en replicación (QA-04), pero pierde en cola acotada por consumidor y en confirmación por mensaje con DLQ, que no ofrece de forma nativa.
  \item NATS JetStream pierde en topología declarada (QA-09) y en experiencia del equipo.
  \item MassTransit gana en confirmación y DLQ (QA-04), pero pierde en topología explícita, porque la oculta, y en licencia, porque sus autores anunciaron licencia comercial para la versión 9.
\end{itemize}

\paragraph{Compromiso aceptado.} La retención de un registro de eventos, que Kafka ofrece, no existe en RabbitMQ.
La recuperación de los mensajes pendientes descansa en la bandeja de salida de cada productor (AD-45).
La política \texttt{plataforma-colas} acota cada cola a 100\,000 mensajes y rechaza la publicación al llenarse.

\paragraph{Riesgo.} Una cola llena rechaza la publicación y el relevo de la bandeja reintenta.
Si la saturación persiste, la bandeja crece en la base del productor.
La alarma de la bandeja no se registra en Auditoría (L-02).

\paragraph{Dónde se ve.} Vista de procesos: exchange \texttt{plataforma.eventos} y colas \texttt{q.*} con su \texttt{dlq.*}.
Vista física de la implementación: nodos \texttt{bus1}, \texttt{bus2} y \texttt{bus3}.

\paragraph{Evidencia en el código.} Las 71 pruebas de \texttt{Plataforma.Comun.Pruebas} con RabbitMQ real verifican consumidor detenido sin pérdida, mensaje envenenado apartado y cola llena sin descarte (QA-04).
Archivos de evidencia:
{\raggedright
\begin{itemize}
  \item \texttt{proyectoparte2/\br{}despliegue/\br{}bus/\br{}definiciones.json}
  \item \texttt{proyectoparte2/\br{}comun/\br{}Plataforma.Comun.Mensajeria}
\end{itemize}
\par}

%---------------------------------------------------------------------- AD-55
\subsection{AD-55. Motor de caché}
\label{sec:tec-ad55}

\paragraph{Contexto.} Esta decisión materializa la caché de geometrías y de respuestas del mapa (AD-17) y la marca de tiempo del dato consultado (AD-16).
El motor atiende QA-06, que exige p95 $<$ 2\,s y antigüedad del dato $\le$ 2\,s en el p95.
También atiende QA-05, porque la caché separa la carga de lectura de la base de datos.

\paragraph{Alternativas.} Se consideran Redis, Valkey, Memcached y una vista materializada en PostgreSQL.
El \cref{tab:tec-ad55} presenta la matriz de decisión ponderada.
La alternativa elegida obtiene 4,70 y la segunda mejor, Valkey, obtiene 4,45.

\begingroup\small
\begin{longtable}{L{4.96cm} C{1.10cm} C{2.25cm} C{2.25cm} C{2.25cm} C{2.25cm}}
\cabecera{6}{tab:tec-ad55}{Matriz de decisión ponderada de AD-55}{\textbf{Criterio (driver)} & \textbf{Peso} & \hspace{0pt}\textbf{Redis 8.8} & \hspace{0pt}\textbf{Valkey} & \hspace{0pt}\textbf{Memcached} & \hspace{0pt}\textbf{Vista materializada}}
Lectura con TTL e invalidación por clave (QA-06, AD-17) & 30\,\% & 5 & 5 & 4 & 3 \\
Antigüedad del dato controlada (QA-06, AD-16) & 20\,\% & 4 & 4 & 4 & 2 \\
Componente separado de la base que escala aparte (QA-05, AD-18) & 15\,\% & 5 & 5 & 5 & 1 \\
Cliente maduro en la pila de AD-50 & 15\,\% & 5 & 4 & 3 & 5 \\
Experiencia del equipo & 10\,\% & 5 & 3 & 2 & 4 \\
Licencia y costo & 10\,\% & 4 & 5 & 5 & 5 \\
\midrule
\textbf{Total ponderado} & 100\,\% & \textbf{4,70} & 4,45 & 3,90 & 3,10 \\
\end{longtable}
\endgroup

\paragraph{Alternativas descartadas.}
\begin{itemize}
  \item Valkey gana en licencia, porque conserva la licencia BSD, pero pierde en experiencia del equipo y en madurez del cliente probado con la pila.
  \item Memcached pierde en invalidación por clave (QA-06), porque no ofrece estructuras ni invalidación por patrón.
  \item La vista materializada gana en cliente y licencia, pero pierde en componente separado (QA-05), porque concentra la lectura en la misma base.
\end{itemize}

\paragraph{Compromiso aceptado.} Redis 8 se distribuye con licencias RSALv2, SSPL y AGPLv3.
La AGPLv3 permite su uso en este proyecto académico.
Valkey es un sustituto compatible con el mismo protocolo si la licencia cambia.

\paragraph{Riesgo.} Riesgo R-14 del ATAM: la antigüedad del dato presentado depende de la carga.
En el perfil local se midió p95 de 1\,485\,ms con máximo de 11,8\,s.
Una corrida anterior dio p95 de 4,2\,s con 10 muestras.

\paragraph{Dónde se ve.} Vista física de la implementación: nodo \texttt{cache}.

\paragraph{Evidencia en el código.} QA-06 se midió con p95 de 142\,ms en el mapa de desastres y con la marca \texttt{datoAl} en el 100\,\% de las respuestas.
Archivos de evidencia:
{\raggedright
\begin{itemize}
  \item \texttt{proyectoparte2/\br{}servicios/\br{}consulta-danos/\br{}src/\br{}ConsultaDanos.Infraestructura/\br{}Cache.cs}
  \item \texttt{proyectoparte2/\br{}servicios/\br{}consulta/\br{}src/\br{}Consulta.Infraestructura/\br{}Adaptadores.cs}
\end{itemize}
\par}

%---------------------------------------------------------------------- AD-56
\subsection{AD-56. Almacén de archivos}
\label{sec:tec-ad56}

\paragraph{Contexto.} Esta decisión materializa el almacén de archivos de la vista física, donde se guardan las evidencias de Evaluación de Daños y los documentos de Reconstrucción.
El almacén debe mantener la evidencia disponible aunque falle una réplica del microservicio (AD-24, QA-02).
También debe conservar el 100\,\% de las fotos sincronizadas desde el cliente móvil (AD-31, QA-14).

\paragraph{Alternativas.} Se consideran MinIO, el disco local del contenedor, una columna binaria en PostgreSQL y SeaweedFS.
El \cref{tab:tec-ad56} presenta la matriz de decisión ponderada.
La alternativa elegida obtiene 4,20 y la segunda mejor, SeaweedFS, obtiene 4,05.

\begingroup\small
\begin{longtable}{L{4.96cm} C{1.10cm} C{2.25cm} C{2.25cm} C{2.25cm} C{2.25cm}}
\cabecera{6}{tab:tec-ad56}{Matriz de decisión ponderada de AD-56}{\textbf{Criterio (driver)} & \textbf{Peso} & \hspace{0pt}\textbf{MinIO} & \hspace{0pt}\textbf{Disco local} & \hspace{0pt}\textbf{Columna binaria} & \hspace{0pt}\textbf{SeaweedFS}}
Evidencia disponible ante la pérdida de una réplica (QA-02, AD-24) & 25\,\% & 4 & 1 & 4 & 4 \\
Integridad de la evidencia sincronizada (QA-14) & 20\,\% & 5 & 3 & 5 & 5 \\
Tamaño de la base y de su respaldo (QA-03, QA-06) & 15\,\% & 5 & 4 & 1 & 5 \\
Interfaz S3 estándar que permite sustituir el producto (QA-10) & 15\,\% & 5 & 1 & 1 & 4 \\
Experiencia del equipo & 10\,\% & 4 & 5 & 4 & 1 \\
Licencia y continuidad del producto & 15\,\% & 2 & 5 & 5 & 4 \\
\midrule
\textbf{Total ponderado} & 100\,\% & \textbf{4,20} & 2,85 & 3,45 & 4,05 \\
\end{longtable}
\endgroup

\paragraph{Alternativas descartadas.}
\begin{itemize}
  \item El disco local gana en experiencia y licencia, pero pierde en disponibilidad ante la pérdida de una réplica (QA-02), que AD-24 descarta.
  \item La columna binaria iguala en integridad (QA-14), pero pierde en tamaño de la base y de su respaldo (QA-03).
  \item SeaweedFS gana en continuidad del producto, pero pierde en experiencia del equipo y en compatibilidad S3 verificada.
\end{itemize}

\paragraph{Compromiso aceptado.} El proyecto original de MinIO dejó de publicar imágenes y el archivo de composición usa el fork comunitario \texttt{pgsty/minio} (DT-11).
El código depende solo de la API S3, lo que permite sustituir el producto sin modificar los microservicios.

\paragraph{Riesgo.} El fork comunitario puede dejar de mantenerse.
En el entorno de un computador \texttt{archivos} es un solo nodo, y su caída deja sin evidencias a las réplicas de \texttt{evaluacion-danos}.
No corresponde a un riesgo del ATAM.

\paragraph{Dónde se ve.} Vista física de la implementación: nodo \texttt{archivos}.

\paragraph{Evidencia en el código.} La corrida de QA-14 conservó 120 de 120 fotos a través del gateway.
Archivos de evidencia:
{\raggedright
\begin{itemize}
  \item \texttt{proyectoparte2/\br{}servicios/\br{}evaluacion-danos/\br{}src/\br{}EvaluacionDanos.Infraestructura/\br{}Archivos/\br{}AlmacenEvidenciasMinio.cs}
  \item \texttt{proyectoparte2/\br{}servicios/\br{}reconstruccion/\br{}src/\br{}Reconstruccion.Infraestructura/\br{}Archivos/\br{}AlmacenDocumentosMinio.cs}
\end{itemize}
\par}

%---------------------------------------------------------------------- AD-57
\subsection{AD-57. API Gateway}
\label{sec:tec-ad57}

\paragraph{Contexto.} Esta decisión materializa el punto único de entrada (AD-04), la priorización y degradación por tipo de cliente (AD-19) y el nivel 1 de la autorización (AD-22).
El gateway atiende QA-11, que exige el rechazo del 100\,\% de las modificaciones de un organismo no responsable.
También atiende QA-05 y QA-06, porque prioriza el tráfico crítico cuando la carga crece.

\paragraph{Alternativas.} Se consideran YARP, Kong, Traefik, Ocelot y nginx.
El \cref{tab:tec-ad57} presenta la matriz de decisión ponderada.
La alternativa elegida obtiene 4,85 y la segunda mejor, Kong, obtiene 3,45.

\begingroup\small
\begin{longtable}{L{3.93cm} C{1.10cm} C{1.95cm} C{1.95cm} C{1.95cm} C{1.95cm} C{1.95cm}}
\cabecera{7}{tab:tec-ad57}{Matriz de decisión ponderada de AD-57}{\textbf{Criterio (driver)} & \textbf{Peso} & \hspace{0pt}\textbf{YARP 2.3} & \hspace{0pt}\textbf{Kong} & \hspace{0pt}\textbf{Traefik} & \hspace{0pt}\textbf{Ocelot} & \hspace{0pt}\textbf{nginx}}
Priorización por clase de tráfico y degradación (QA-05, QA-06, AD-19) & 25\,\% & 5 & 3 & 2 & 3 & 2 \\
Validación JWT y regla de dominio por ruta (QA-11, AD-22) & 25\,\% & 5 & 4 & 3 & 4 & 2 \\
Catálogo de rutas validado al arrancar (QA-11) & 15\,\% & 5 & 4 & 3 & 3 & 2 \\
Latencia agregada por el intermediario (QA-06) & 15\,\% & 4 & 4 & 4 & 3 & 5 \\
Experiencia del equipo en la pila de AD-50 & 10\,\% & 5 & 2 & 3 & 4 & 3 \\
Licencia y mantenimiento & 10\,\% & 5 & 3 & 5 & 3 & 5 \\
\midrule
\textbf{Total ponderado} & 100\,\% & \textbf{4,85} & 3,45 & 3,10 & 3,35 & 2,85 \\
\end{longtable}
\endgroup

\paragraph{Alternativas descartadas.}
\begin{itemize}
  \item Kong pierde en priorización con degradación (QA-05), que exigiría complementos propios en otro lenguaje.
  \item Traefik pierde en priorización y en validación de la regla de dominio por ruta (QA-11).
  \item Ocelot pierde en catálogo validado al arrancar (QA-11) y en mantenimiento del proyecto.
  \item nginx gana en latencia agregada (QA-06), pero pierde en validación JWT con regla de dominio (QA-11), que no ofrece sin módulos adicionales.
\end{itemize}

\paragraph{Compromiso aceptado.} La priorización y la regla de dominio son código propio en el gateway.
Ese código se prueba con la misma suite de .NET, pero lo mantiene el equipo.

\paragraph{Riesgo.} Riesgo R-12 del ATAM: el gateway recibe tráfico sin TLS en el entorno de laboratorio (L-07).
Riesgo R-13 del ATAM: la validación depende de las llaves y del emisor de Keycloak (L-08).
Los rechazos del gateway no se registran en Auditoría (L-02).

\paragraph{Dónde se ve.} Vista lógica: componente API Gateway.
Vista física de la implementación: nodo \texttt{gateway} con dos réplicas y el artefacto \texttt{rutas.json}.

\paragraph{Evidencia en el código.} QA-11 se verificó con 218 de 218 intentos rechazados (100\,\%), 0 falsos negativos y 0 datos modificados sobre 12 rutas con regla de dominio.
Archivos de evidencia:
{\raggedright
\begin{itemize}
  \item \texttt{proyectoparte2/\br{}gateway/\br{}src/\br{}Plataforma.Gateway/\br{}rutas.json}
  \item \texttt{proyectoparte2/\br{}gateway/\br{}src/\br{}Plataforma.Gateway/\br{}Prioridad/\br{}Priorizacion.cs}
  \item \texttt{proyectoparte2/\br{}gateway/\br{}src/\br{}Plataforma.Gateway/\br{}Seguridad/\br{}AutorizacionNivel1.cs}
\end{itemize}
\par}

%---------------------------------------------------------------------- AD-58
\subsection{AD-58. Proveedor de identidad y clientes OIDC}
\label{sec:tec-ad58}

\paragraph{Contexto.} Esta decisión materializa el proveedor de identidad federado con OIDC y PKCE (AD-49) y la validación del token en el gateway y en cada microservicio (AD-04, AD-22).
Sus clientes son la consola, el cliente móvil y el tablero de daños. El tablero usa el cliente confidencial \texttt{consulta-\br{}danos-\br{}tablero} y una sesión con cookie que valida la propia página, porque su ruta es pública en el gateway (E-41).
El proveedor atiende QA-11, porque emite los roles y el atributo de organismo que usa la regla de dominio.
También atiende QA-14, porque el refresco del token sin conexión permite trabajar días en campo y sincronizar al volver.

\paragraph{Alternativas.} Se consideran Keycloak, Auth0, Duende IdentityServer y una identidad propia.
El \cref{tab:tec-ad58} presenta la matriz de decisión ponderada.
La alternativa elegida obtiene 4,90 y la segunda mejor, Duende IdentityServer, obtiene 3,75.

\begingroup\small
\begin{longtable}{L{4.96cm} C{1.10cm} C{2.25cm} C{2.25cm} C{2.25cm} C{2.25cm}}
\cabecera{6}{tab:tec-ad58}{Matriz de decisión ponderada de AD-58}{\textbf{Criterio (driver)} & \textbf{Peso} & \hspace{0pt}\textbf{Keycloak 26} & \hspace{0pt}\textbf{Auth0} & \hspace{0pt}\textbf{Duende IdentityServer} & \hspace{0pt}\textbf{Identidad propia}}
OIDC con PKCE para la consola y el móvil, con refresco sin conexión (QA-11, QA-14, AD-49) & 25\,\% & 5 & 5 & 4 & 2 \\
Roles y organismo como \emph{claims} para la regla de dominio (QA-11, AD-22) & 20\,\% & 5 & 4 & 4 & 3 \\
Despliegue dentro de la plataforma, sin servicio externo en la emergencia (QA-02, CON-3) & 15\,\% & 5 & 2 & 5 & 5 \\
Configuración declarada en un archivo versionado (QA-10) & 10\,\% & 5 & 3 & 4 & 3 \\
Experiencia del equipo & 10\,\% & 4 & 2 & 2 & 3 \\
Licencia y costo & 20\,\% & 5 & 2 & 3 & 5 \\
\midrule
\textbf{Total ponderado} & 100\,\% & \textbf{4,90} & 3,25 & 3,75 & 3,45 \\
\end{longtable}
\endgroup

\paragraph{Alternativas descartadas.}
\begin{itemize}
  \item Auth0 iguala en OIDC con PKCE (QA-11), pero pierde en despliegue dentro de la plataforma (CON-3) y en costo por usuario activo.
  \item Duende IdentityServer pierde en licencia, porque exige licencia comercial por encima de un umbral de ingresos, y en experiencia del equipo.
  \item La identidad propia gana en licencia y despliegue, pero pierde en OIDC con PKCE (QA-11), porque obligaría a construir la seguridad desde cero.
\end{itemize}

\paragraph{Compromiso aceptado.} El token se valida dos veces, en el gateway y en el microservicio, con un validador propio en \texttt{Plataforma.Comun.Seguridad}.
La doble validación agrega trabajo por petición a cambio de la defensa en profundidad de AD-22.

\paragraph{Riesgo.} Riesgo R-13 del ATAM: el emisor del token es \texttt{http://localhost:8180} y debe coincidir en el gateway y en los microservicios (L-08).
La recarga de las llaves JWT está pendiente (L-09).
Riesgo R-12 del ATAM: Keycloak corre en \texttt{start-dev} sin TLS (L-07).

\paragraph{Dónde se ve.} Vista lógica: componente externo Proveedor de identidad.
Vista física de la implementación: nodo \texttt{keycloak}.
Vista de implementación móvil: paquete \texttt{infraestructura/identidad}.

\paragraph{Evidencia en el código.} Siete pruebas instrumentadas de la app corren contra el gateway y Keycloak, incluido el rechazo por rol.
Archivos de evidencia:
{\raggedright
\begin{itemize}
  \item \texttt{proyectoparte2/\br{}despliegue/\br{}keycloak/\br{}realm-plataforma.json}
  \item \texttt{proyectoparte2/\br{}comun/\br{}Plataforma.Comun.Seguridad/\br{}IdentidadPorJwt.cs}
  \item \texttt{danos/\br{}app/\br{}src/\br{}main/\br{}java/\br{}com/\br{}example/\br{}danos/\br{}infraestructura/\br{}identidad/\br{}ConexionOidcBuilder.kt}
\end{itemize}
\par}

%---------------------------------------------------------------------- AD-59
\subsection{AD-59. Consola de coordinación}
\label{sec:tec-ad59}

\paragraph{Contexto.} Esta decisión materializa la capa de presentación web (AD-08) de la Consola de coordinación, con los módulos \texttt{mapa}, \texttt{evaluaciones} y \texttt{reconstruccion}.
La consola entra por el gateway (AD-04) y atiende QA-06 en la carga del mapa, QA-11 en el inicio de sesión y QA-10 en la incorporación de módulos.

\paragraph{Alternativas.} Se consideran Angular, React 19, Vue 3 y Blazor.
El \cref{tab:tec-ad59} presenta la matriz de decisión ponderada.
La alternativa elegida obtiene 4,80 y la segunda mejor, Vue 3, obtiene 4,10.

\begingroup\small
\begin{longtable}{L{4.96cm} C{1.10cm} C{2.25cm} C{2.25cm} C{2.25cm} C{2.25cm}}
\cabecera{6}{tab:tec-ad59}{Matriz de decisión ponderada de AD-59}{\textbf{Criterio (driver)} & \textbf{Peso} & \hspace{0pt}\textbf{Angular 22} & \hspace{0pt}\textbf{React 19} & \hspace{0pt}\textbf{Vue 3} & \hspace{0pt}\textbf{Blazor}}
Módulos con rutas, inyección y guardas (QA-10, AD-08) & 25\,\% & 5 & 3 & 4 & 4 \\
Cliente OIDC con PKCE certificado (QA-11, AD-49) & 20\,\% & 5 & 4 & 4 & 4 \\
Carga inicial y respuesta de la interfaz (QA-06) & 20\,\% & 4 & 4 & 5 & 2 \\
Archivos estáticos servidos detrás del gateway (QA-02, AD-04) & 10\,\% & 5 & 5 & 5 & 4 \\
Experiencia del equipo & 15\,\% & 5 & 3 & 2 & 3 \\
Licencia y costo & 10\,\% & 5 & 5 & 5 & 5 \\
\midrule
\textbf{Total ponderado} & 100\,\% & \textbf{4,80} & 3,80 & 4,10 & 3,55 \\
\end{longtable}
\endgroup

\paragraph{Alternativas descartadas.}
\begin{itemize}
  \item React 19 pierde en estructura por módulos con inyección y guardas (QA-10), que exige elegir y combinar bibliotecas.
  \item Vue 3 gana en carga inicial (QA-06), pero pierde en experiencia del equipo y en estructura por módulos.
  \item Blazor pierde en carga inicial (QA-06), por el tamaño del entorno WebAssembly que descarga el navegador.
\end{itemize}

\paragraph{Compromiso aceptado.} Angular produce un paquete inicial mayor que Vue.
La carga se mitiga con módulos cargados bajo demanda por ruta.

\paragraph{Riesgo.} Angular publica una versión mayor cada seis meses y obliga a actualizaciones frecuentes.
No corresponde a un riesgo del ATAM.

\paragraph{Dónde se ve.} Vista física de la implementación: nodo \texttt{web}.
Vista lógica: Consola de coordinación con los módulos \texttt{mapa}, \texttt{evaluaciones} y \texttt{reconstruccion}.

\paragraph{Evidencia en el código.} El \texttt{Dockerfile} compila con \texttt{node:22-alpine} y sirve los archivos con \texttt{nginx:1.29-alpine}.
Archivos de evidencia:
{\raggedright
\begin{itemize}
  \item \texttt{proyectoparte2/\br{}consola/\br{}package.json}
  \item \texttt{proyectoparte2/\br{}consola/\br{}Dockerfile}
  \item \texttt{proyectoparte2/\br{}consola/\br{}nginx.conf}
  \item \texttt{proyectoparte2/\br{}consola/\br{}src/\br{}app/\br{}modulos}
\end{itemize}
\par}

%---------------------------------------------------------------------- AD-60
\subsection{AD-60. Biblioteca de mapas}
\label{sec:tec-ad60}

\paragraph{Contexto.} Esta decisión materializa la presentación de las geometrías precalculadas (AD-17) y de su marca de tiempo (AD-16) en el módulo \texttt{mapa} de la consola.
La biblioteca atiende QA-06, porque el mapa es la consulta frecuente del coordinador.

\paragraph{Alternativas.} Se consideran Leaflet, OpenLayers, MapLibre GL y Google Maps.
El \cref{tab:tec-ad60} presenta la matriz de decisión ponderada.
La alternativa elegida obtiene 4,80 y la segunda mejor, OpenLayers, obtiene 3,90.

\begingroup\small
\begin{longtable}{L{4.96cm} C{1.10cm} C{2.25cm} C{2.25cm} C{2.25cm} C{2.25cm}}
\cabecera{6}{tab:tec-ad60}{Matriz de decisión ponderada de AD-60}{\textbf{Criterio (driver)} & \textbf{Peso} & \hspace{0pt}\textbf{Leaflet 1.9} & \hspace{0pt}\textbf{OpenLayers} & \hspace{0pt}\textbf{MapLibre GL} & \hspace{0pt}\textbf{Google Maps}}
Capas GeoJSON agregadas por zona (QA-06, AD-17) & 25\,\% & 5 & 5 & 5 & 4 \\
Peso de la biblioteca y tiempo de carga (QA-06) & 20\,\% & 5 & 3 & 2 & 3 \\
Independencia de un proveedor y de una llave (QA-02, CON-3) & 20\,\% & 4 & 4 & 4 & 1 \\
Integración con Angular y con el mapa del prototipo (QA-10) & 10\,\% & 5 & 3 & 3 & 3 \\
Experiencia del equipo & 10\,\% & 5 & 2 & 2 & 3 \\
Licencia y costo & 15\,\% & 5 & 5 & 5 & 2 \\
\midrule
\textbf{Total ponderado} & 100\,\% & \textbf{4,80} & 3,90 & 3,70 & 2,70 \\
\end{longtable}
\endgroup

\paragraph{Alternativas descartadas.}
\begin{itemize}
  \item OpenLayers iguala en capas GeoJSON (QA-06), pero pierde en peso de la biblioteca y en experiencia del equipo.
  \item MapLibre GL iguala en capas GeoJSON, pero pierde en peso de la biblioteca (QA-06), porque su ventaja con millones de puntos no hace falta con zonas agregadas.
  \item Google Maps pierde en independencia de un proveedor (CON-3) y en costo por carga de mapa.
\end{itemize}

\paragraph{Compromiso aceptado.} Leaflet dibuja en el navegador sin aceleración WebGL.
El volumen se controla con la agregación por zona del servidor (AD-17).

\paragraph{Riesgo.} Las teselas de OpenStreetMap no tienen acuerdo de servicio y su política de uso prohíbe el uso intensivo.
No corresponde a un riesgo del ATAM.

\paragraph{Dónde se ve.} Vista física de la implementación: nodo \texttt{web}.
Vista lógica: módulo \texttt{mapa} de la Consola de coordinación.

\paragraph{Evidencia en el código.} QA-06 se midió con p95 de 1\,628\,ms y p99 de 1\,935\,ms en el mapa de daños en el perfil local (L-04).
Archivos de evidencia:
{\raggedright
\begin{itemize}
  \item \texttt{proyectoparte2/\br{}consola/\br{}src/\br{}app/\br{}modulos/\br{}mapa/\br{}lienzo-mapa.ts}
  \item \texttt{proyectoparte2/\br{}consola/\br{}package.json}
\end{itemize}
\par}

%---------------------------------------------------------------------- AD-61
\subsection{AD-61. Plataforma del cliente móvil}
\label{sec:tec-ad61}

\paragraph{Contexto.} Esta decisión materializa la réplica local del cliente móvil (AD-31), la cola local de operaciones (AD-44) y el registro del desastre en campo (AD-48).
La plataforma atiende QA-14, cuya medida es 0 evaluaciones perdidas o sobrescritas en silencio al sincronizar 120 evaluaciones, y cuya respuesta exige 0 duplicados.

\paragraph{Alternativas.} Se consideran Android nativo con Kotlin, Flutter, React Native y Kotlin Multiplatform.
El \cref{tab:tec-ad61} presenta la matriz de decisión ponderada.
La alternativa elegida obtiene 4,60 y la segunda mejor, Kotlin Multiplatform, obtiene 4,00.

\begingroup\small
\begin{longtable}{L{4.96cm} C{1.10cm} C{2.25cm} C{2.25cm} C{2.25cm} C{2.25cm}}
\cabecera{6}{tab:tec-ad61}{Matriz de decisión ponderada de AD-61}{\textbf{Criterio (driver)} & \textbf{Peso} & \hspace{0pt}\textbf{Kotlin nativo} & \hspace{0pt}\textbf{Flutter} & \hspace{0pt}\textbf{React Native} & \hspace{0pt}\textbf{Kotlin Multiplatform}}
Trabajo en segundo plano y almacenamiento del sistema (QA-14, AD-44) & 25\,\% & 5 & 3 & 3 & 4 \\
Base local con transacciones y migraciones (QA-14, AD-31) & 20\,\% & 5 & 4 & 3 & 4 \\
Cámara y ubicación para las evidencias (QA-14) & 10\,\% & 5 & 4 & 4 & 4 \\
Reutilización del prototipo de FP-06 y de sus pruebas (QA-14) & 15\,\% & 5 & 1 & 1 & 4 \\
Portabilidad a otro sistema operativo móvil (QA-10) & 10\,\% & 1 & 5 & 5 & 4 \\
Experiencia del equipo & 10\,\% & 5 & 2 & 3 & 3 \\
Licencia y costo & 10\,\% & 5 & 5 & 5 & 5 \\
\midrule
\textbf{Total ponderado} & 100\,\% & \textbf{4,60} & 3,30 & 3,20 & 4,00 \\
\end{longtable}
\endgroup

\paragraph{Alternativas descartadas.}
\begin{itemize}
  \item Flutter gana en portabilidad (QA-10), pero pierde en reutilización del prototipo de FP-06 (QA-14), que habría que reescribir.
  \item React Native gana en portabilidad, pero pierde en base local transaccional (QA-14) y en reutilización del prototipo.
  \item Kotlin Multiplatform pierde en trabajo en segundo plano (QA-14), porque esa parte sigue siendo específica de cada plataforma.
\end{itemize}

\paragraph{Compromiso aceptado.} El cliente funciona solo en Android.
Un cliente para iOS exigiría otra implementación de la capa de presentación y de la sincronización.

\paragraph{Riesgo.} La app de entrega es la variante \texttt{debug}, que admite \texttt{http://} (L-07).
La evidencia solo admite fotos y no guarda un autor propio (L-10).

\paragraph{Dónde se ve.} Vista de implementación móvil: paquete \texttt{presentacion}.
Vista física de la implementación: \texttt{<<device>>} Teléfono Android con \texttt{DanosCampo.apk}.

\paragraph{Evidencia en el código.} La app \texttt{co.plataforma.danos} 2.1.0 tiene 81 pruebas unitarias y 16 instrumentadas en verde.
Archivos de evidencia:
{\raggedright
\begin{itemize}
  \item \texttt{danos/\br{}gradle/\br{}libs.versions.toml}
  \item \texttt{danos/\br{}app/\br{}build.gradle.kts}
\end{itemize}
\par}

%---------------------------------------------------------------------- AD-62
\subsection{AD-62. Base de datos local del cliente móvil}
\label{sec:tec-ad62}

\paragraph{Contexto.} Esta decisión materializa la réplica local (AD-31), la cola local de operaciones inmutables con copia completa del contenido (AD-44) y la tabla \texttt{registro\_desastre} del registro en campo (AD-48).
La base atiende QA-14, porque la escritura y su operación en la cola deben confirmarse en la misma transacción.

\paragraph{Alternativas.} Se consideran Room, Realm, SQLDelight y DataStore.
El \cref{tab:tec-ad62} presenta la matriz de decisión ponderada.
La alternativa elegida obtiene 4,90 y la segunda mejor, SQLDelight, obtiene 4,35.

\begingroup\small
\begin{longtable}{L{4.96cm} C{1.10cm} C{2.25cm} C{2.25cm} C{2.25cm} C{2.25cm}}
\cabecera{6}{tab:tec-ad62}{Matriz de decisión ponderada de AD-62}{\textbf{Criterio (driver)} & \textbf{Peso} & \hspace{0pt}\textbf{Room 2.8} & \hspace{0pt}\textbf{Realm} & \hspace{0pt}\textbf{SQLDelight} & \hspace{0pt}\textbf{DataStore}}
Transacción atómica entre la escritura y la cola local (QA-14, AD-44) & 30\,\% & 5 & 4 & 5 & 1 \\
Migraciones explícitas y verificables (QA-10, QA-14) & 20\,\% & 5 & 3 & 4 & 2 \\
Integración con Compose y corrutinas (AD-61) & 15\,\% & 5 & 4 & 4 & 4 \\
Continuidad del producto & 15\,\% & 5 & 1 & 5 & 5 \\
Experiencia del equipo & 10\,\% & 5 & 2 & 2 & 4 \\
Control del SQL verificado en compilación (QA-14) & 10\,\% & 4 & 2 & 5 & 1 \\
\midrule
\textbf{Total ponderado} & 100\,\% & \textbf{4,90} & 2,95 & 4,35 & 2,55 \\
\end{longtable}
\endgroup

\paragraph{Alternativas descartadas.}
\begin{itemize}
  \item Realm pierde en continuidad del producto, porque su fabricante retiró el SDK de dispositivos.
  \item SQLDelight gana en control del SQL verificado en compilación, pero pierde en migraciones con esquema exportado y en experiencia del equipo.
  \item DataStore pierde en transacción atómica con la cola (QA-14), porque guarda pares clave y valor sin relaciones.
\end{itemize}

\paragraph{Compromiso aceptado.} Cada cambio de esquema exige una migración escrita a mano y su esquema exportado.
Ese costo es el precio de no perder datos locales al actualizar la app (QA-14).

\paragraph{Riesgo.} Una migración defectuosa puede impedir abrir la base local con operaciones pendientes.
No corresponde a un riesgo del ATAM.

\paragraph{Dónde se ve.} Vista de implementación móvil: paquete \texttt{infraestructura/local}.
Vista física de la implementación: artefacto \texttt{danos.db} en el Teléfono Android.

\paragraph{Evidencia en el código.} \texttt{AppDatabase} declara la versión 3 con \texttt{MIGRACION\_1\_2} y \texttt{MIGRACION\_2\_3}, y la carpeta de esquemas guarda las versiones 1, 2 y 3.
Archivos de evidencia:
{\raggedright
\begin{itemize}
  \item \texttt{danos/\br{}app/\br{}src/\br{}main/\br{}java/\br{}com/\br{}example/\br{}danos/\br{}infraestructura/\br{}local/\br{}AppDatabase.kt}
  \item \texttt{danos/\br{}app/\br{}schemas}
\end{itemize}
\par}

%---------------------------------------------------------------------- AD-63
\subsection{AD-63. Sincronización del cliente móvil}
\label{sec:tec-ad63}

\paragraph{Contexto.} Esta decisión materializa la sincronización diferida (AD-31), el reenvío de la cola local (AD-44) y el envío idempotente del registro del desastre por el identificador del cliente (AD-48).
La sincronización atiende QA-14, cuya respuesta exige reenvío sin duplicados y 120 evaluaciones en $\le$ 10\,min.
También atiende QA-11, porque todo envío pasa por el gateway con el token OIDC.

\paragraph{Alternativas.} Se consideran WorkManager, SyncAdapter, un servicio en primer plano y Firebase.
El \cref{tab:tec-ad63} presenta la matriz de decisión ponderada.
La alternativa elegida obtiene 4,85 y la segunda mejor, SyncAdapter, obtiene 4,10.

\begingroup\small
\begin{longtable}{L{4.96cm} C{1.10cm} C{2.25cm} C{2.25cm} C{2.25cm} C{2.25cm}}
\cabecera{6}{tab:tec-ad63}{Matriz de decisión ponderada de AD-63}{\textbf{Criterio (driver)} & \textbf{Peso} & \hspace{0pt}\textbf{WorkManager} & \hspace{0pt}\textbf{SyncAdapter} & \hspace{0pt}\textbf{Servicio en primer plano} & \hspace{0pt}\textbf{Firebase}}
Ejecución garantizada con red disponible y reintento tras reinicio (QA-14, AD-44) & 30\,\% & 5 & 4 & 2 & 4 \\
Protocolo propio idempotente por identificador del cliente (QA-14, AD-48) & 25\,\% & 5 & 5 & 5 & 1 \\
Inmediatez del envío al recuperar la red (QA-14) & 15\,\% & 4 & 3 & 5 & 4 \\
Envío por el gateway con token OIDC (QA-11, AD-04) & 10\,\% & 5 & 5 & 5 & 1 \\
Experiencia del equipo & 10\,\% & 5 & 2 & 3 & 3 \\
Licencia y dependencia de un servicio externo & 10\,\% & 5 & 5 & 5 & 2 \\
\midrule
\textbf{Total ponderado} & 100\,\% & \textbf{4,85} & 4,10 & 3,90 & 2,65 \\
\end{longtable}
\endgroup

\paragraph{Alternativas descartadas.}
\begin{itemize}
  \item SyncAdapter pierde en ejecución garantizada (QA-14) y en experiencia del equipo, porque exige una cuenta del sistema y un proveedor de contenido.
  \item El servicio en primer plano gana en inmediatez del envío, pero pierde en ejecución garantizada tras un reinicio (QA-14).
  \item Firebase pierde en protocolo idempotente propio (QA-14) y en envío por el gateway (QA-11), porque sincroniza contra su propio servicio.
\end{itemize}

\paragraph{Compromiso aceptado.} El sistema operativo decide el instante de la ejecución de WorkManager.
El envío puede demorar unos segundos después de recuperar la red.

\paragraph{Riesgo.} Riesgo R-12 del ATAM: la variante \texttt{debug} admite tráfico sin TLS (L-07).
Riesgo R-13 del ATAM: desde un teléfono físico el emisor del token exige cambiar la configuración de Keycloak (L-08).

\paragraph{Dónde se ve.} Vista de implementación móvil: paquete \texttt{infraestructura/sync}.
Vista física de la implementación: \texttt{DanosCampo.apk} en el Teléfono Android.

\paragraph{Evidencia en el código.} QA-14 se verificó con 115 evaluaciones aplicadas, 5 en conflicto, 0 pérdidas, 0 sobrescrituras y 0 duplicados en 27,8\,s.
Archivos de evidencia:
{\raggedright
\begin{itemize}
  \item \texttt{danos/\br{}app/\br{}src/\br{}main/\br{}java/\br{}com/\br{}example/\br{}danos/\br{}infraestructura/\br{}sync/\br{}SincronizacionWorker.kt}
  \item \texttt{danos/\br{}app/\br{}src/\br{}main/\br{}java/\br{}com/\br{}example/\br{}danos/\br{}infraestructura/\br{}sync/\br{}EnviadorDesastres.kt}
\end{itemize}
\par}

%---------------------------------------------------------------------- AD-64
\subsection{AD-64. Biblioteca de resiliencia}
\label{sec:tec-ad64}

\paragraph{Contexto.} Esta decisión materializa el tiempo de espera, el reintento con espera creciente y el cortacircuitos de las fronteras (AD-05) y el cortacircuitos por adaptador (AD-37).
La biblioteca atiende QA-04, cuya respuesta exige cortar la propagación del fallo en $\le$ 60\,s.
También atiende QA-01, porque la caída de un sistema externo no debe degradar las demás funciones.

\paragraph{Alternativas.} Se consideran Microsoft.Extensions.Http.Resilience, una implementación propia y Steeltoe.
El \cref{tab:tec-ad64} presenta la matriz de decisión ponderada.
La alternativa elegida obtiene 4,90 y la segunda mejor, Implementación propia, obtiene 3,15.

\begingroup\small
\begin{longtable}{L{6.74cm} C{1.10cm} C{2.50cm} C{2.50cm} C{2.50cm}}
\cabecera{5}{tab:tec-ad64}{Matriz de decisión ponderada de AD-64}{\textbf{Criterio (driver)} & \textbf{Peso} & \hspace{0pt}\textbf{MS Resilience (Polly v8)} & \hspace{0pt}\textbf{Implementación propia} & \hspace{0pt}\textbf{Steeltoe}}
Tiempo de espera, reintento con variación aleatoria y cortacircuitos configurables (QA-01, AD-05) & 30\,\% & 5 & 3 & 2 \\
Corte de la propagación en $\le$ 60\,s verificable (QA-04) & 20\,\% & 5 & 3 & 3 \\
Integración con la fábrica de clientes HTTP de AD-50 & 15\,\% & 5 & 3 & 4 \\
Estado del circuito seguro ante concurrencia (QA-04) & 15\,\% & 5 & 2 & 4 \\
Experiencia del equipo & 10\,\% & 4 & 4 & 1 \\
Licencia y mantenimiento & 10\,\% & 5 & 5 & 4 \\
\midrule
\textbf{Total ponderado} & 100\,\% & \textbf{4,90} & 3,15 & 2,90 \\
\end{longtable}
\endgroup

\paragraph{Alternativas descartadas.}
\begin{itemize}
  \item La implementación propia iguala en experiencia y licencia, pero pierde en estado del circuito seguro ante concurrencia (QA-04).
  \item Steeltoe pierde en políticas configurables (QA-01) y en experiencia del equipo.
\end{itemize}

\paragraph{Compromiso aceptado.} Polly mide una proporción de fallos en una ventana, no un número de fallos consecutivos.
La regla de AD-05 se expresa como 3 llamadas mínimas con 100\,\% de fallos en 30\,s.

\paragraph{Riesgo.} Un reintento sobre una operación no idempotente duplicaría el efecto.
\texttt{Resiliencia.cs} reintenta solo operaciones idempotentes.
No corresponde a un riesgo del ATAM.

\paragraph{Dónde se ve.} Vista de implementación: proyecto \texttt{Plataforma.Comun.Adaptadores} (\texttt{Resiliencia.cs}).
Modelo de integración: cortacircuitos de cada adaptador.

\paragraph{Evidencia en el código.} La configuración usa 3\,s por intento, 2 reintentos, ventana de 30\,s y apertura de 15\,s. Las 12 pruebas de \texttt{AdaptadorSe6.Integracion} y las 14 de \texttt{AdaptadorSe7.Integracion} verifican el cortacircuitos y el reintento sin pérdida (QA-01).
Archivos de evidencia:
{\raggedright
\begin{itemize}
  \item \texttt{proyectoparte2/\br{}comun/\br{}Plataforma.Comun.Adaptadores/\br{}Resiliencia.cs}
  \item \texttt{proyectoparte2/\br{}Directory.Packages.props}
\end{itemize}
\par}

%---------------------------------------------------------------------- AD-65
\subsection{AD-65. Empaquetado, ejecución y escalado}
\label{sec:tec-ad65}

\paragraph{Contexto.} Esta decisión materializa las réplicas adicionales del camino de lectura (AD-18), la réplica activa con sondeo de salud (AD-24) y la ejecución híbrida con adaptadores efímeros (AD-25).
El empaquetado atiende QA-05, que exige escalar solo los componentes bajo presión sin modificar código.
También atiende QA-02 por el reemplazo de réplicas y QA-01 por el escalado de los adaptadores según su cola.

\paragraph{Alternativas.} Se consideran contenedores con Kubernetes, máquinas virtuales, Docker Swarm, Nomad y Knative.
El \cref{tab:tec-ad65} presenta la matriz de decisión ponderada.
La alternativa elegida obtiene 4,60 y la segunda mejor, Knative, obtiene 3,85.

\begingroup\small
\begin{longtable}{L{3.93cm} C{1.10cm} C{1.95cm} C{1.95cm} C{1.95cm} C{1.95cm} C{1.95cm}}
\cabecera{7}{tab:tec-ad65}{Matriz de decisión ponderada de AD-65}{\textbf{Criterio (driver)} & \textbf{Peso} & \hspace{0pt}\textbf{Docker + K8s (HPA, KEDA)} & \hspace{0pt}\textbf{Máquinas virtuales} & \hspace{0pt}\textbf{Docker Swarm} & \hspace{0pt}\textbf{Nomad} & \hspace{0pt}\textbf{Knative}}
Escalado automático selectivo sin modificar código (QA-05, AD-18) & 25\,\% & 5 & 2 & 2 & 4 & 5 \\
Escalado de 0 a N por largo de cola (QA-01, AD-25) & 15\,\% & 5 & 1 & 1 & 3 & 3 \\
Sondeo de salud y reemplazo de réplicas (QA-02, AD-24) & 20\,\% & 5 & 3 & 4 & 4 & 5 \\
Misma imagen en el computador y en el clúster (QA-09) & 15\,\% & 5 & 2 & 5 & 4 & 4 \\
Experiencia del equipo & 10\,\% & 4 & 3 & 3 & 1 & 1 \\
Costo de operación y licencia & 15\,\% & 3 & 2 & 5 & 3 & 3 \\
\midrule
\textbf{Total ponderado} & 100\,\% & \textbf{4,60} & 2,15 & 3,25 & 3,40 & 3,85 \\
\end{longtable}
\endgroup

\paragraph{Alternativas descartadas.}
\begin{itemize}
  \item Las máquinas virtuales pierden en escalado selectivo (QA-05) y en escalado por largo de cola (QA-01).
  \item Docker Swarm gana en costo de operación, pero pierde en escalado automático (QA-05), que no ofrece.
  \item Nomad pierde en escalado por cola y en licencia, porque cambió a BSL en 2023.
  \item Knative iguala en escalado selectivo y sondeo de salud, pero pierde en escalado por largo de cola (QA-01), porque está orientado a HTTP.
\end{itemize}

\paragraph{Compromiso aceptado.} En Docker Compose los adaptadores quedan siempre encendidos y no hay escalado automático.
El escalado de 0 a N solo existe en los manifiestos de Kubernetes (AD-25).

\paragraph{Riesgo.} El despliegue en Kubernetes no se verificó de punta a punta (L-05).
El escalado de 0 a N con KEDA no se midió (QA-01) y QA-05 no se midió con 10\,000 usuarios (L-03, L-04).

\paragraph{Dónde se ve.} Vista física de la implementación: nodo Docker Engine con los servicios del archivo de composición y anotación de Kubernetes con HPA y KEDA.

\paragraph{Evidencia en el código.} Los manifiestos declaran HPA de 2 a 6 réplicas para \texttt{consulta} y \texttt{consulta-\br{}danos}, de 2 a 5 para \texttt{consulta-\br{}programa}, y \texttt{ScaledObject} de KEDA de 0 a 4 réplicas para la salida de SE-6.
Archivos de evidencia:
{\raggedright
\begin{itemize}
  \item \texttt{proyectoparte2/\br{}despliegue/\br{}compose/\br{}docker-compose.yml}
  \item \texttt{proyectoparte2/\br{}despliegue/\br{}k8s/\br{}base/\br{}20-servicios.yaml}
  \item \texttt{proyectoparte2/\br{}despliegue/\br{}k8s/\br{}base/\br{}adaptador-se6.yaml}
\end{itemize}
\par}

%---------------------------------------------------------------------- AD-66
\subsection{AD-66. Plataforma del adaptador de SE-8}
\label{sec:tec-ad66}

\paragraph{Contexto.} CON-9, agregada el 19 de noviembre de 2026, exige que la solución final incluya Jakarta EE y .NET (E-44).
AD-03 ubica la heterogeneidad tecnológica (CON-2) en los adaptadores, no en los microservicios: un adaptador es el único punto de la arquitectura donde una pila distinta no toca el núcleo.
El adaptador de SE-8 estaba diseñado y no implementado (L-06).
QA-08 le exige una instancia encendida (AD-25), estado de entrega por alerta y \enquote{entregada} solo con acuse.
QA-09 pide que incorporar un sistema externo modifique 0 componentes existentes.
Esta decisión elige la plataforma de ese adaptador.

\paragraph{Alternativas.} Se consideran Jakarta EE 11 sobre WildFly 42.0.0.Final (imagen \texttt{quay.io/\br{}wildfly/\br{}wildfly:\br{}42.0.0.Final-\br{}jdk21}, Java 21), Spring Boot 3 sobre Java 21 y .NET 10 con \texttt{Plataforma.\br{}Comun.\br{}Adaptadores}, como los adaptadores de SE-6 y SE-7.
El \cref{tab:tec-ad66} presenta la matriz de decisión ponderada.
La alternativa elegida obtiene 4,35 y la segunda mejor, .NET 10, obtiene 4,00.

\begingroup\small
\begin{longtable}{L{6.0cm} C{1.10cm} C{2.6cm} C{2.6cm} C{2.6cm}}
\cabecera{5}{tab:tec-ad66}{Matriz de decisión ponderada de AD-66}{\textbf{Criterio (driver)} & \textbf{Peso} & \hspace{0pt}\textbf{Jakarta EE 11 (WildFly 42)} & \hspace{0pt}\textbf{Spring Boot 3} & \hspace{0pt}\textbf{.NET 10}}
Cumplimiento de CON-9 & 25\,\% & 5 & 2 & 1 \\
Adaptador nuevo sin modificar componentes existentes (QA-09, AD-03) & 20\,\% & 5 & 5 & 5 \\
Cortacircuitos del Cuadro 10 con aviso de apertura para EV-23 (QA-01, QA-08, AD-05, AD-37) & 15\,\% & 3 & 4 & 5 \\
Consumo transaccional: ack después del commit y DLQ a la quinta entrega (AD-02, AD-45) & 15\,\% & 4 & 4 & 5 \\
Instancia siempre encendida: memoria y arranque (QA-08, AD-25) & 10\,\% & 4 & 4 & 5 \\
Experiencia del equipo (especificaciones vistas en el curso) & 10\,\% & 4 & 3 & 5 \\
Licencia y costo & 5\,\% & 5 & 5 & 5 \\
\midrule
\textbf{Total ponderado} & 100\,\% & \textbf{4,35} & 3,65 & 4,00 \\
\end{longtable}
\endgroup

\paragraph{Alternativas descartadas.}
\begin{itemize}
  \item .NET 10 gana en resiliencia, mensajería y arranque, porque reutiliza \texttt{Plataforma.\br{}Comun.\br{}Adaptadores} y \texttt{ConsumidorIdempotente}. Pierde en el criterio que origina la decisión: no cumple CON-9.
  \item Spring Boot 3 corre sobre Java y usa algunas API de Jakarta, pero no es la plataforma Jakarta EE: cumple CON-9 solo en parte. Tiene un cortacircuitos con ventana de tiempo (Resilience4j) y es una tercera pila de marco que el equipo no vio en el curso.
\end{itemize}

\paragraph{Compromiso aceptado.} El adaptador usa las especificaciones estándar: Jakarta Messaging, EJB (\texttt{@Stateless}, \texttt{@Singleton}, Timer Service), Jakarta Persistence, Jakarta REST, CDI, Jakarta Concurrency y MicroProfile Config.
Dos piezas son propias y no vienen del estándar:
\begin{itemize}
  \item \textbf{Cortacircuitos propio en lugar de MicroProfile Fault Tolerance.} El \texttt{@CircuitBreaker} de MicroProfile mide una ventana de N llamadas, no la ventana de \qty{30}{\second} del Cuadro 10, y no avisa cuando el circuito se abre o se cierra. Ese aviso hace falta para publicar EV-23 por la bandeja (E-15). La clase \texttt{Cortacircuitos} aplica la regla de Polly (al menos 3 llamadas en \qty{30}{\second} con \qty{100}{\percent} de fallo, \qty{15}{\second} abierto y una sola llamada de prueba) y es segura ante concurrencia. Tiene 9 pruebas unitarias propias, una de ellas con 200 hilos.
  \item \textbf{EJB \texttt{@Singleton @Startup} en lugar de MDB.} Un MDB necesita un resource adapter JCA para el proveedor de mensajería, y RabbitMQ no tiene uno mantenido. El consumo usa igual la API estándar de Jakarta Messaging 3.1 (\texttt{Connection}, \texttt{Session} transaccional, \texttt{MessageConsumer}) con el cliente JMS de RabbitMQ en modo AMQP sobre la cola existente. El hilo lo da un \texttt{ManagedExecutorService}. Con \texttt{nackOnRollback}, \texttt{rollback()} manda el mensaje a la DLQ y cerrar la sesión lo reencola.
\end{itemize}
Se acepta además una segunda pila que mantener en la frontera. Esa pila duplica en Java la lectura del sobre y la parte del catálogo que usa el adaptador. Los esquemas JSON siguen siendo una sola fuente: el WAR los copia de \texttt{Plataforma.\br{}Comun.\br{}Eventos} al compilar.

\paragraph{Riesgo.} Un cambio del sobre canónico o de las claves de ruteo obliga a cambiar dos implementaciones. Lo detecta la validación contra los mismos esquemas, en las pruebas y al consumir. Corresponde al riesgo R-15 del ATAM.
Mientras Notificación no exista, el adaptador recibe EV-03 y EV-04 como sustituto temporal (L-11).

\paragraph{Dónde se ve.} Vista lógica: componente Adaptador de SE-8 [\texttt{adaptador-se8}] «implementado».
Vista de procesos: proceso \texttt{adaptador-se8} ×1, siempre encendido.
Vista física de la implementación: entorno de ejecución WildFly 42 dentro del contenedor \texttt{adaptador-se8} (\cref{fig:vista-despliegue-impl}, \cref{tab:nodos-impl}).
Vista de implementación: paquete \texttt{adaptador-\br{}se8-\br{}jee} (\cref{fig:vista-implementacion}).
Modelo de integración: fila de SE-8 del \cref{tab:integracion}.

\paragraph{Evidencia en el código.} 52 pruebas unitarias y 7 de integración en verde. Las de integración usan Testcontainers con la topología real del bus (cola quorum con \texttt{delivery-limit} 5) y la base del migrador. Demuestran el ack después del commit JPA, la DLQ a la quinta entrega fallida, el duplicado sin segundo envío y el circuito con EV-23.
La prueba \texttt{prueba-se8.sh} mide contra el compose: \qty{1,8}{\second} del registro a la alerta entregada con acuse, circuito abierto a los \qty{25,6}{\second} de la caída, y cola en 0 a los \qty{12,9}{\second} de la vuelta, con las 20 alertas entregadas una sola vez.
El \cref{tab:que-se-modifica} registra los 0 componentes existentes modificados.
Archivos de evidencia:
{\raggedright
\begin{itemize}
  \item \texttt{proyectoparte2/\br{}adaptadores/\br{}adaptador-se8-jee/\br{}adaptador/\br{}src/\br{}main/\br{}java/\br{}plataforma/\br{}adaptador/\br{}se8/\br{}jee/\br{}AdaptadorSe8.java}
  \item \texttt{proyectoparte2/\br{}adaptadores/\br{}adaptador-se8-jee/\br{}adaptador/\br{}src/\br{}main/\br{}java/\br{}plataforma/\br{}adaptador/\br{}se8/\br{}consumo/\br{}BucleConsumo.java}
  \item \texttt{proyectoparte2/\br{}adaptadores/\br{}adaptador-se8-jee/\br{}adaptador/\br{}src/\br{}main/\br{}java/\br{}plataforma/\br{}adaptador/\br{}se8/\br{}resiliencia/\br{}Cortacircuitos.java}
  \item \texttt{proyectoparte2/\br{}adaptadores/\br{}adaptador-se8-jee/\br{}Dockerfile}
  \item \texttt{proyectoparte2/\br{}pruebas/\br{}humo/\br{}prueba-se8.sh}
\end{itemize}
\par}

%------------------------------------------------------------------------------
\subsection{Restricciones de las tecnologías elegidas}
\label{sec:tec-restricciones}

Las tecnologías elegidas imponen límites al entorno implementado.
El \cref{tab:tec-restricciones} los relaciona con la decisión que los origina y con el escenario afectado.
Ninguno modifica una decisión estructural.

\begingroup\small
\begin{longtable}{L{2.6cm} L{1.3cm} L{6.4cm} L{1.7cm} L{3.3cm}}
\cabecera{5}{tab:tec-restricciones}{Restricciones que imponen las tecnologías elegidas}{\textbf{Restricción} &
\textbf{Origen} & \textbf{Efecto} & \textbf{QA afectado} & \textbf{Qué haría falta}}
Variante \texttt{debug} y Keycloak sin TLS (L-07) & AD-58, AD-61, AD-63 &
La app de entrega admite \texttt{http://}. Keycloak corre sin TLS. El tráfico entre el teléfono, el gateway y el proveedor de identidad viaja sin cifrar. Corresponde al riesgo R-12 del ATAM. &
QA-11, QA-12 & TLS en el borde y URL HTTPS en la variante \texttt{release} \\
Emisor del token en \texttt{localhost} (L-08) & AD-57, AD-58 &
El \texttt{iss} del token es \texttt{http://localhost:8180}. El emulador exige \texttt{adb reverse}. Un teléfono físico exige cambiar \texttt{KC\_HOSTNAME} e \texttt{Identidad\_\_Emisor} a la vez. Corresponde al riesgo R-13 del ATAM. &
QA-11, QA-14 & Un nombre DNS común para Keycloak \\
Imagen comunitaria de MinIO (DT-11) & AD-56 &
El proyecto original dejó de publicar imágenes. El archivo de composición usa el fork \texttt{pgsty/minio}, cuya continuidad no está garantizada. &
QA-02, QA-14 & Sustituir por otro almacén con API S3 sin cambiar el código \\
Política de uso de las teselas de OpenStreetMap & AD-60 &
Las teselas públicas no tienen acuerdo de servicio, exigen atribución y prohíben el uso intensivo. Durante una emergencia con 10\,000 usuarios la carga del mapa dependería de un tercero. &
QA-05, QA-06 & Servidor de teselas propio o proveedor con acuerdo \\
Keycloak en \texttt{start-dev} & AD-58 &
El modo de desarrollo usa una base embebida, no exige TLS y no está pensado para producción. El realm se importa desde \texttt{realm-plataforma.json} en cada arranque. Es una sola instancia. &
QA-02, QA-11 & Modo \texttt{start} con base PostgreSQL, TLS y dos instancias \\
Sustituto temporal de Notificación (L-11) & AD-66 &
EV-03 y EV-04 llegan a \texttt{q.salida.se8} por dos enlaces temporales, y el adaptador alerta a un único destino de difusión. No hay suscripciones ni resolución por ámbito, y el estado de entrega se consulta en el adaptador. &
QA-08 & Implementar Notificación (AD-13), quitar los enlaces y poner \texttt{Sustituto\_\_Activo=false} \\
Dos pilas en la frontera & AD-66 &
El adaptador de SE-8 duplica en Java la lectura del sobre y la parte del catálogo que usa. Un cambio del contrato se hace en .NET y en Java. Corresponde al riesgo R-15 del ATAM. &
QA-09, QA-10 & Mantener los esquemas JSON como única fuente y validarlos en las pruebas de las dos pilas \\
\end{longtable}
\endgroup
.
%  Generado a partir de la especificación de la versión 1.2. No carga paquetes.
%==============================================================================
\section{Decisiones de tecnología}
\label{sec:tecnologia}

El enunciado del proyecto (§8) pide seleccionar las tecnologías en la segunda entrega.
CON-8 exigía independencia tecnológica en la primera entrega.
Esta sección sustituye esa condición para la segunda entrega.
El Anexo~\ref{sec:anexo-a} deja de ser una excepción a CON-8 y describe la implementación de FP-06 con las tecnologías elegidas aquí.

Las decisiones se numeran de AD-50 a AD-66.
AD-66 responde a CON-9, que se agregó después de AD-50 a AD-65.
Ninguna modifica una decisión estructural.
Cada una materializa uno o varios mecanismos de AD-01 a AD-49, y su justificación parte de los escenarios QA-nn que ese mecanismo debe cumplir (CON-7).

\subsection{Método de decisión}
\label{sec:tec-metodo}

Cada decisión sigue cinco pasos:
\begin{enumerate}
  \item Se identifica el mecanismo del SAD que la tecnología materializa y los escenarios QA-nn que ese mecanismo atiende.
  \item Se eligen entre tres y cinco alternativas con uso real para ese mecanismo.
  \item Se definen los criterios. Cada criterio cita el QA-nn, AD-nn o CON-n del que proviene. La experiencia del equipo y la licencia o el costo se incluyen con un peso menor que los criterios de calidad.
  \item Se asigna un peso a cada criterio. Los pesos de una decisión suman 100\,\%.
  \item Se califica cada alternativa en cada criterio y se calcula el total ponderado.
\end{enumerate}

La escala de calificación tiene cinco niveles:
\begin{itemize}
  \item 1: la alternativa no satisface el criterio o lo impide.
  \item 2: lo satisface con desarrollo propio considerable.
  \item 3: lo satisface con configuración adicional o con complementos de terceros.
  \item 4: lo satisface de forma nativa con limitaciones menores.
  \item 5: lo satisface de forma nativa y la implementación lo verifica.
\end{itemize}

El total ponderado de una alternativa es $\sum_i p_i \, c_i / 100$, donde $p_i$ es el peso en porcentaje y $c_i$ la calificación del criterio $i$.
El máximo es 5.
Las cifras de verificación citadas provienen del perfil local de un computador, con 20 clientes, y no reemplazan la prueba con 10\,000 usuarios (L-04).

\subsection{Resumen de las decisiones}
\label{sec:tec-resumen}

El \cref{tab:tecnologia-resumen} relaciona cada decisión con su tecnología, el mecanismo que materializa y los escenarios que favorece.
Las versiones son las de los archivos de proyecto.

\begingroup\small
\begin{longtable}{L{1.15cm} L{2.2cm} L{6.0cm} L{2.6cm} L{2.6cm}}
\cabecera{5}{tab:tecnologia-resumen}{Resumen de las decisiones de tecnología}{\textbf{ID} & \textbf{Categoría} &
\textbf{Tecnología y versión} & \textbf{Mecanismo} & \textbf{QA que favorece}}
AD-50 & Lenguaje y plataforma & C\# sobre .NET 10 (SDK 10.0.100, \texttt{net10.0}) & AD-01, AD-08, AD-09, AD-45 & QA-06, QA-07, QA-14 \\
AD-51 & Marco web & ASP.NET Core 10 (controladores y \emph{minimal APIs}) & AD-08, AD-22, AD-24 & QA-06, QA-07, QA-11 \\
AD-52 & Persistencia & PostgreSQL 18 con PostGIS 3.6 (\texttt{postgis/postgis:18-3.6}), una base por microservicio y réplica \texttt{db-replica} & AD-01, AD-17, AD-20, AD-24, AD-41, AD-45 & QA-03, QA-04, QA-06, QA-13, QA-15 \\
AD-53 & Acceso a datos & Entity Framework Core 10.0.14, Npgsql 10.0.3 y NetTopologySuite 10.0.3, con contenedor migrador & AD-08, AD-17, AD-41, AD-45 & QA-06, QA-10, QA-13 \\
AD-54 & Mensajería & RabbitMQ 4.4 (3 nodos, colas quorum acotadas, DLQ) con RabbitMQ.Client 7.2.2 & AD-02, AD-05, AD-36, AD-45 & QA-01, QA-04, QA-08 \\
AD-55 & Caché & Redis 8.8 (\texttt{redis:8.8.4-alpine}) con StackExchange.Redis 3.4.29 & AD-16, AD-17 & QA-05, QA-06 \\
AD-56 & Almacén de objetos & MinIO con API S3 (paquete Minio 7.0.0) & AD-24, AD-31 & QA-02, QA-03, QA-14 \\
AD-57 & Puerta de entrada & YARP 2.3.0 con catálogo \texttt{rutas.json} (56 rutas) & AD-04, AD-19, AD-22 & QA-05, QA-06, QA-11 \\
AD-58 & Identidad & Keycloak 26.5.2 (realm \texttt{plataforma}), AppAuth-Android 0.11.1, angular-auth-oidc-client 22.0.2, cliente OpenID Connect de ASP.NET Core para el tablero y validador JWT propio en \texttt{Plataforma.Comun.Seguridad} & AD-04, AD-22, AD-49 & QA-11, QA-14 \\
AD-59 & Cliente web & Angular 22.2 (TypeScript 6.0) servido por nginx 1.29 detrás del gateway & AD-04, AD-08 & QA-06, QA-10, QA-11 \\
AD-60 & Mapas & Leaflet 1.9.4 con teselas de OpenStreetMap & AD-16, AD-17 & QA-06 \\
AD-61 & Cliente móvil & Android nativo: Kotlin 2.0.21, Jetpack Compose (BOM 2024.09.00), Hilt 2.60.1, minSdk 26, targetSdk 36 & AD-31, AD-44, AD-48 & QA-07, QA-14 \\
AD-62 & Persistencia móvil & Room 2.8.5 sobre SQLite, migraciones explícitas 1$\to$2$\to$3 & AD-31, AD-44, AD-48 & QA-14 \\
AD-63 & Sincronización móvil & WorkManager 2.12.0, Retrofit 3.0.0, OkHttp 5.4.0 y kotlinx.serialization 1.9.0 & AD-31, AD-44, AD-48 & QA-11, QA-14 \\
AD-64 & Resiliencia & Microsoft.Extensions.Http.Resilience 10.11.0 (Polly v8) & AD-05, AD-37 & QA-01, QA-04 \\
AD-65 & Despliegue & Contenedores Docker (\texttt{mcr.microsoft.com/\br{}dotnet/\br{}aspnet:10.0}), Docker Compose para un computador y Kubernetes con HPA y KEDA para escalar & AD-18, AD-24, AD-25 & QA-01, QA-02, QA-05 \\
AD-66 & Adaptador de SE-8 & Jakarta EE 11 sobre WildFly 42.0.0.Final (Java 21): Jakarta Messaging con el cliente JMS de RabbitMQ 3.9.1, EJB, Jakarta Persistence, Jakarta REST y MicroProfile Config & AD-03, AD-05, AD-09, AD-25, AD-45 & QA-08, QA-09, CON-9 \\
\end{longtable}
\endgroup

%---------------------------------------------------------------------- AD-50
\subsection{AD-50. Lenguaje y plataforma de los microservicios}
\label{sec:tec-ad50}

\paragraph{Contexto.} Esta decisión materializa los microservicios con almacén propio (AD-01), las cuatro capas de cada microservicio (AD-08), el contrato común de los adaptadores (AD-09) y la bandeja de salida transaccional (AD-45).
El lenguaje determina la latencia del camino síncrono, que exige QA-06 (p95 $<$ 2\,s) y QA-07 (p95 $<$ 3\,s).
También determina si la base común \texttt{comun} y el prototipo de FP-06 del Anexo~\ref{sec:anexo-a} se reutilizan, lo que afecta la verificación de QA-14.
CON-2 exige integrarse con sistemas de tecnologías diferentes, y esa heterogeneidad la absorben los adaptadores (AD-03), no el lenguaje de los microservicios.

\paragraph{Alternativas.} Se consideran C\# sobre .NET 10, Java 21, Node.js 22 con TypeScript y Go.
Las cuatro tienen licencia libre y ejecución en contenedores.
El \cref{tab:tec-ad50} presenta la matriz de decisión ponderada.
La alternativa elegida obtiene 4,90 y la segunda mejor, Java 21, obtiene 3,85.

\begingroup\small
\begin{longtable}{L{4.96cm} C{1.10cm} C{2.25cm} C{2.25cm} C{2.25cm} C{2.25cm}}
\cabecera{6}{tab:tec-ad50}{Matriz de decisión ponderada de AD-50}{\textbf{Criterio (driver)} & \textbf{Peso} & \hspace{0pt}\textbf{.NET 10} & \hspace{0pt}\textbf{Java 21} & \hspace{0pt}\textbf{Node.js 22} & \hspace{0pt}\textbf{Go}}
Rendimiento del camino síncrono (QA-06, QA-07) & 25\,\% & 5 & 4 & 3 & 5 \\
Bibliotecas maduras para bus, resiliencia, datos espaciales y JWT (AD-02, AD-05, AD-17, AD-22) & 25\,\% & 5 & 5 & 3 & 3 \\
Reutilización de \texttt{comun} y del prototipo de FP-06 (QA-14, AD-09) & 20\,\% & 5 & 2 & 2 & 2 \\
Memoria y arranque de cada réplica (QA-05, AD-18) & 10\,\% & 4 & 3 & 4 & 5 \\
Experiencia del equipo & 10\,\% & 5 & 4 & 3 & 1 \\
Licencia y costo & 10\,\% & 5 & 5 & 5 & 5 \\
\midrule
\textbf{Total ponderado} & 100\,\% & \textbf{4,90} & 3,85 & 3,10 & 3,50 \\
\end{longtable}
\endgroup

\paragraph{Alternativas descartadas.}
\begin{itemize}
  \item Java 21 pierde en la reutilización de \texttt{comun} y del prototipo de FP-06 (QA-14), que habría que reescribir.
  \item Node.js 22 pierde en el rendimiento del camino síncrono (QA-06, QA-07) y en la madurez de las bibliotecas para datos espaciales y bus.
  \item Go gana en memoria y arranque de cada réplica (QA-05), pero pierde en bibliotecas para los mecanismos del SAD y en experiencia del equipo.
\end{itemize}

\paragraph{Compromiso aceptado.} Una sola pila reduce el costo de mantener \texttt{comun} (AD-09) y acepta que todos los microservicios dependen del ciclo de soporte de .NET.
La heterogeneidad que exige CON-2 queda confinada en los adaptadores (AD-03).
El núcleo sigue en .NET aun con CON-9: Jakarta EE se confina a la frontera según AD-03 y AD-66.

\paragraph{Riesgo.} El fin del soporte de .NET 10 obligaría a migrar todos los microservicios a la vez.
El archivo \texttt{global.json} fija el SDK con avance solo de la versión de características, lo que acota este riesgo.
No corresponde a un riesgo del ATAM.

\paragraph{Dónde se ve.} Vista de implementación: proyectos \texttt{*.Api}, \texttt{*.Aplicacion}, \texttt{*.Dominio} e \texttt{*.Infraestructura} de cada microservicio.
Vista física de la implementación: artefactos \texttt{plataforma/<servicio>:local} sobre \texttt{mcr.microsoft.com/\br{}dotnet/\br{}aspnet:10.0}.

\paragraph{Evidencia en el código.} La suite de .NET suma 1\,008 pruebas en 19 proyectos en verde. QA-07 se midió con p95 de 1\,334\,ms en el perfil local (L-04).
Archivos de evidencia:
{\raggedright
\begin{itemize}
  \item \texttt{proyectoparte2/\br{}global.json}
  \item \texttt{proyectoparte2/\br{}Directory.Build.props}
\end{itemize}
\par}

%---------------------------------------------------------------------- AD-51
\subsection{AD-51. Marco de las API de los microservicios}
\label{sec:tec-ad51}

\paragraph{Contexto.} Esta decisión materializa la capa de acceso de cada microservicio (AD-08), el nivel 2 de la autorización (AD-22) y el sondeo de salud de la réplica activa (AD-24).
El marco atiende la latencia de QA-06 y QA-07 y el rechazo del 100\,\% de las modificaciones no autorizadas de QA-11.
La elección está condicionada por AD-50, porque el marco debe pertenecer a la pila elegida.

\paragraph{Alternativas.} Se consideran ASP.NET Core 10, Spring Boot 3 y NestJS.
Spring Boot 3 y NestJS se evalúan como el marco que se habría usado con Java y con Node.js.
El \cref{tab:tec-ad51} presenta la matriz de decisión ponderada.
La alternativa elegida obtiene 4,85 y la segunda mejor, Spring Boot 3, obtiene 3,75.

\begingroup\small
\begin{longtable}{L{6.74cm} C{1.10cm} C{2.50cm} C{2.50cm} C{2.50cm}}
\cabecera{5}{tab:tec-ad51}{Matriz de decisión ponderada de AD-51}{\textbf{Criterio (driver)} & \textbf{Peso} & \hspace{0pt}\textbf{ASP.NET Core 10} & \hspace{0pt}\textbf{Spring Boot 3} & \hspace{0pt}\textbf{NestJS}}
Latencia por petición (QA-06, QA-07) & 25\,\% & 5 & 4 & 3 \\
Validación JWT y autorización de nivel 2 en la tubería (QA-11, AD-22) & 20\,\% & 5 & 5 & 4 \\
Sondeos de salud vivo y listo (QA-02, QA-04, AD-24) & 15\,\% & 4 & 5 & 4 \\
Coherencia con AD-50 y con el gateway (QA-10, AD-57) & 20\,\% & 5 & 1 & 1 \\
Experiencia del equipo & 10\,\% & 5 & 3 & 3 \\
Licencia y costo & 10\,\% & 5 & 5 & 5 \\
\midrule
\textbf{Total ponderado} & 100\,\% & \textbf{4,85} & 3,75 & 3,15 \\
\end{longtable}
\endgroup

\paragraph{Alternativas descartadas.}
\begin{itemize}
  \item Spring Boot 3 gana en sondeos de salud con Actuator (QA-02), pero pierde en coherencia con AD-50, porque introduce una segunda pila.
  \item NestJS pierde en latencia por petición (QA-06, QA-07) y en coherencia con AD-50.
\end{itemize}

\paragraph{Compromiso aceptado.} Los sondeos de salud se componen con \texttt{Plataforma.Comun.Salud} en lugar de un módulo integrado como Actuator (AD-24).

\paragraph{Riesgo.} Una actualización mayor del marco afecta a todos los microservicios a la vez.
No corresponde a un riesgo del ATAM.

\paragraph{Dónde se ve.} Vista de implementación: proyectos \texttt{*.Api} con \texttt{Sdk.Web}.
Vista física de la implementación: artefactos \texttt{plataforma/<servicio>:local}.

\paragraph{Evidencia en el código.} Los proyectos \texttt{*.Api.csproj} usan \texttt{Microsoft.NET.Sdk.Web}. QA-07 se cumple con 100 de 100 registros respondidos con 201.
Archivos de evidencia:
{\raggedright
\begin{itemize}
  \item \texttt{proyectoparte2/\br{}servicios/\br{}gestion-desastres/\br{}src/\br{}GestionDesastres.Api}
  \item \texttt{proyectoparte2/\br{}comun/\br{}Plataforma.Comun.Salud}
\end{itemize}
\par}

%---------------------------------------------------------------------- AD-52
\subsection{AD-52. Motor de base de datos}
\label{sec:tec-ad52}

\paragraph{Contexto.} Esta decisión materializa el almacén propio de cada microservicio (AD-01), la geometría precalculada de las zonas (AD-17), los historiales versionados y de solo inserción (AD-20, AD-41), la réplica activa (AD-24) y la bandeja de salida (AD-45).
El motor atiende la consulta espacial de QA-06, la publicación sin pérdida de QA-04, la inmutabilidad de QA-13 y QA-15 y el RPO $\le$ 60\,s de QA-03.

\paragraph{Alternativas.} Se consideran PostgreSQL con PostGIS, MySQL 8, SQL Server y MongoDB.
El \cref{tab:tec-ad52} presenta la matriz de decisión ponderada.
La alternativa elegida obtiene 4,85 y la segunda mejor, SQL Server, obtiene 4,10.

\begingroup\small
\begin{longtable}{L{4.96cm} C{1.10cm} C{2.25cm} C{2.25cm} C{2.25cm} C{2.25cm}}
\cabecera{6}{tab:tec-ad52}{Matriz de decisión ponderada de AD-52}{\textbf{Criterio (driver)} & \textbf{Peso} & \hspace{0pt}\textbf{PostgreSQL + PostGIS} & \hspace{0pt}\textbf{MySQL 8} & \hspace{0pt}\textbf{SQL Server} & \hspace{0pt}\textbf{MongoDB}}
Consulta espacial con índice (QA-06, AD-17) & 25\,\% & 5 & 3 & 4 & 3 \\
Transacción ACID para la bandeja y el historial (QA-04, AD-20, AD-45) & 20\,\% & 5 & 4 & 5 & 3 \\
Inmutabilidad garantizada por permisos del almacén (QA-13, QA-15, AD-41) & 15\,\% & 5 & 4 & 5 & 3 \\
Réplica continua para la réplica activa y el RPO (QA-03, AD-24, AD-27) & 15\,\% & 4 & 4 & 5 & 5 \\
Experiencia del equipo & 10\,\% & 5 & 4 & 3 & 3 \\
Licencia y costo & 15\,\% & 5 & 4 & 2 & 3 \\
\midrule
\textbf{Total ponderado} & 100\,\% & \textbf{4,85} & 3,75 & 4,10 & 3,30 \\
\end{longtable}
\endgroup

\paragraph{Alternativas descartadas.}
\begin{itemize}
  \item MySQL 8 pierde en consulta espacial con índice (QA-06), porque sus funciones espaciales son menos completas que PostGIS.
  \item SQL Server gana en réplica continua (QA-03), pero pierde en licencia y costo, porque su licencia comercial se cobra por núcleo.
  \item MongoDB gana en réplica continua (QA-03), pero pierde en transacción ACID para la bandeja (QA-04) y en inmutabilidad por permisos (QA-13).
\end{itemize}

\paragraph{Compromiso aceptado.} En el entorno de un computador las nueve bases comparten la instancia \texttt{db}.
La separación por microservicio se conserva por base y por rol con mínimo privilegio (AD-01), no por servidor.

\paragraph{Riesgo.} La instancia \texttt{db} es un punto único de fallo en el entorno de un computador.
La réplica \texttt{db-replica} reduce la pérdida de datos, pero la promoción no es automática y no hay segunda región (L-01).
QA-02 y QA-03 no se midieron.

\paragraph{Dónde se ve.} Vista física de la implementación: nodos \texttt{db} y \texttt{db-replica}.
Vista de implementación: paquete \texttt{Infraestructura/Persistencia} de cada microservicio y migradores.

\paragraph{Evidencia en el código.} Los servicios \texttt{db} y \texttt{db-replica} usan la imagen \texttt{postgis/postgis:18-3.6}. QA-06 se midió con p95 de 34\,ms en el consolidado del desastre en el perfil local (L-04).
Archivos de evidencia:
{\raggedright
\begin{itemize}
  \item \texttt{proyectoparte2/\br{}despliegue/\br{}compose/\br{}docker-compose.yml}
\end{itemize}
\par}

%---------------------------------------------------------------------- AD-53
\subsection{AD-53. Acceso a datos y migraciones}
\label{sec:tec-ad53}

\paragraph{Contexto.} Esta decisión materializa la capa de infraestructura de cada microservicio (AD-08), el mapeo de las geometrías (AD-17), los permisos del historial de solo inserción (AD-41) y la escritura atómica con la bandeja (AD-45).
El acceso a datos afecta la latencia de QA-06, la evolución del esquema que pide QA-10 y la inmutabilidad de QA-13.

\paragraph{Alternativas.} Se consideran Entity Framework Core, Dapper, NHibernate y SQL directo con Npgsql sin mapeador.
El \cref{tab:tec-ad53} presenta la matriz de decisión ponderada.
La alternativa elegida obtiene 4,80 y la segunda mejor, Dapper, obtiene 3,60.

\begingroup\small
\begin{longtable}{L{4.96cm} C{1.10cm} C{2.25cm} C{2.25cm} C{2.25cm} C{2.25cm}}
\cabecera{6}{tab:tec-ad53}{Matriz de decisión ponderada de AD-53}{\textbf{Criterio (driver)} & \textbf{Peso} & \hspace{0pt}\textbf{EF Core 10} & \hspace{0pt}\textbf{Dapper} & \hspace{0pt}\textbf{NHibernate} & \hspace{0pt}\textbf{SQL directo}}
Migraciones versionadas con los permisos del historial (QA-10, QA-13, AD-41) & 25\,\% & 5 & 2 & 3 & 1 \\
Tipos espaciales mapeados (QA-06, AD-17) & 20\,\% & 5 & 3 & 2 & 3 \\
Control del SQL en las lecturas críticas (QA-06) & 20\,\% & 4 & 5 & 3 & 5 \\
Aislamiento en la capa de infraestructura y pruebas (QA-10, AD-08) & 15\,\% & 5 & 4 & 4 & 3 \\
Experiencia del equipo & 10\,\% & 5 & 4 & 2 & 4 \\
Licencia y costo & 10\,\% & 5 & 5 & 5 & 5 \\
\midrule
\textbf{Total ponderado} & 100\,\% & \textbf{4,80} & 3,60 & 3,05 & 3,20 \\
\end{longtable}
\endgroup

\paragraph{Alternativas descartadas.}
\begin{itemize}
  \item Dapper gana en control del SQL (QA-06), pero pierde en migraciones versionadas (QA-10, QA-13), que exigirían otra herramienta.
  \item NHibernate pierde en tipos espaciales mapeados (QA-06) y en experiencia del equipo.
  \item SQL directo gana en control del SQL (QA-06), pero pierde en migraciones versionadas y en aislamiento de la capa de infraestructura (QA-10).
\end{itemize}

\paragraph{Compromiso aceptado.} El mapeador agrega una capa sobre el SQL.
Las lecturas críticas del modelo de lectura (AD-15) pueden usar SQL explícito dentro del mismo contexto cuando QA-06 lo exija.
El esquema lo aplica un contenedor migrador con el rol dueño, y el microservicio corre con un rol sin permisos de DDL (AD-41).

\paragraph{Riesgo.} Una consulta generada por el mapeador puede degradar la latencia del mapa de daños, que se midió con p95 de 1\,628\,ms frente al límite de 2\,s de QA-06.
No corresponde a un riesgo del ATAM.

\paragraph{Dónde se ve.} Vista de implementación: paquete \texttt{Infraestructura/Persistencia} y migradores (\texttt{*-migrador}).
Vista física de la implementación: nodos \texttt{db} y \texttt{db-replica}.

\paragraph{Evidencia en el código.} El servicio \texttt{reconstruccion-migrador} del archivo de composición crea los roles y aplica las migraciones antes de arrancar el microservicio.
Archivos de evidencia:
{\raggedright
\begin{itemize}
  \item \texttt{proyectoparte2/\br{}Directory.Packages.props}
  \item \texttt{proyectoparte2/\br{}comun/\br{}Plataforma.Comun.Migraciones}
  \item \texttt{proyectoparte2/\br{}servicios/\br{}evaluacion-danos/\br{}src/\br{}EvaluacionDanos.Infraestructura/\br{}Persistencia}
\end{itemize}
\par}

%---------------------------------------------------------------------- AD-54
\subsection{AD-54. Bus de eventos}
\label{sec:tec-ad54}

\paragraph{Contexto.} Esta decisión materializa el bus de eventos con colas por consumidor (AD-02), las colas acotadas de la frontera (AD-05), la cola de salida con reintento (AD-36) y la entrega desde la bandeja de salida (AD-45).
El bus atiende QA-04, cuya respuesta exige 0 eventos perdidos ante el fallo de un componente interno, y QA-01, que exige procesar el 100\,\% de lo encolado cuando el sistema externo vuelve.

\paragraph{Alternativas.} Se consideran RabbitMQ, Apache Kafka, NATS JetStream y MassTransit sobre RabbitMQ.
MassTransit se evalúa como capa de abstracción sobre el mismo intermediario.
El \cref{tab:tec-ad54} presenta la matriz de decisión ponderada.
La alternativa elegida obtiene 4,90 y la segunda mejor, NATS JetStream, obtiene 3,85.

\begingroup\small
\begin{longtable}{L{4.96cm} C{1.10cm} C{2.25cm} C{2.25cm} C{2.25cm} C{2.25cm}}
\cabecera{6}{tab:tec-ad54}{Matriz de decisión ponderada de AD-54}{\textbf{Criterio (driver)} & \textbf{Peso} & \hspace{0pt}\textbf{RabbitMQ 4.4} & \hspace{0pt}\textbf{Kafka} & \hspace{0pt}\textbf{NATS JetStream} & \hspace{0pt}\textbf{MassTransit}}
Cola acotada por consumidor con rechazo al productor (QA-04, AD-05) & 25\,\% & 5 & 2 & 4 & 4 \\
Confirmación por mensaje, límite de entregas y DLQ (QA-01, QA-04, AD-36) & 20\,\% & 5 & 2 & 4 & 5 \\
Replicación de la cola ante la caída de un nodo (QA-02, QA-04) & 15\,\% & 5 & 5 & 5 & 5 \\
Topología explícita y declarada al desplegar (QA-09, QA-10) & 15\,\% & 5 & 3 & 3 & 2 \\
Experiencia del equipo & 10\,\% & 4 & 2 & 1 & 3 \\
Licencia y costo & 15\,\% & 5 & 5 & 5 & 2 \\
\midrule
\textbf{Total ponderado} & 100\,\% & \textbf{4,90} & 3,05 & 3,85 & 3,65 \\
\end{longtable}
\endgroup

\paragraph{Alternativas descartadas.}
\begin{itemize}
  \item Kafka iguala en replicación (QA-04), pero pierde en cola acotada por consumidor y en confirmación por mensaje con DLQ, que no ofrece de forma nativa.
  \item NATS JetStream pierde en topología declarada (QA-09) y en experiencia del equipo.
  \item MassTransit gana en confirmación y DLQ (QA-04), pero pierde en topología explícita, porque la oculta, y en licencia, porque sus autores anunciaron licencia comercial para la versión 9.
\end{itemize}

\paragraph{Compromiso aceptado.} La retención de un registro de eventos, que Kafka ofrece, no existe en RabbitMQ.
La recuperación de los mensajes pendientes descansa en la bandeja de salida de cada productor (AD-45).
La política \texttt{plataforma-colas} acota cada cola a 100\,000 mensajes y rechaza la publicación al llenarse.

\paragraph{Riesgo.} Una cola llena rechaza la publicación y el relevo de la bandeja reintenta.
Si la saturación persiste, la bandeja crece en la base del productor.
La alarma de la bandeja no se registra en Auditoría (L-02).

\paragraph{Dónde se ve.} Vista de procesos: exchange \texttt{plataforma.eventos} y colas \texttt{q.*} con su \texttt{dlq.*}.
Vista física de la implementación: nodos \texttt{bus1}, \texttt{bus2} y \texttt{bus3}.

\paragraph{Evidencia en el código.} Las 71 pruebas de \texttt{Plataforma.Comun.Pruebas} con RabbitMQ real verifican consumidor detenido sin pérdida, mensaje envenenado apartado y cola llena sin descarte (QA-04).
Archivos de evidencia:
{\raggedright
\begin{itemize}
  \item \texttt{proyectoparte2/\br{}despliegue/\br{}bus/\br{}definiciones.json}
  \item \texttt{proyectoparte2/\br{}comun/\br{}Plataforma.Comun.Mensajeria}
\end{itemize}
\par}

%---------------------------------------------------------------------- AD-55
\subsection{AD-55. Motor de caché}
\label{sec:tec-ad55}

\paragraph{Contexto.} Esta decisión materializa la caché de geometrías y de respuestas del mapa (AD-17) y la marca de tiempo del dato consultado (AD-16).
El motor atiende QA-06, que exige p95 $<$ 2\,s y antigüedad del dato $\le$ 2\,s en el p95.
También atiende QA-05, porque la caché separa la carga de lectura de la base de datos.

\paragraph{Alternativas.} Se consideran Redis, Valkey, Memcached y una vista materializada en PostgreSQL.
El \cref{tab:tec-ad55} presenta la matriz de decisión ponderada.
La alternativa elegida obtiene 4,70 y la segunda mejor, Valkey, obtiene 4,45.

\begingroup\small
\begin{longtable}{L{4.96cm} C{1.10cm} C{2.25cm} C{2.25cm} C{2.25cm} C{2.25cm}}
\cabecera{6}{tab:tec-ad55}{Matriz de decisión ponderada de AD-55}{\textbf{Criterio (driver)} & \textbf{Peso} & \hspace{0pt}\textbf{Redis 8.8} & \hspace{0pt}\textbf{Valkey} & \hspace{0pt}\textbf{Memcached} & \hspace{0pt}\textbf{Vista materializada}}
Lectura con TTL e invalidación por clave (QA-06, AD-17) & 30\,\% & 5 & 5 & 4 & 3 \\
Antigüedad del dato controlada (QA-06, AD-16) & 20\,\% & 4 & 4 & 4 & 2 \\
Componente separado de la base que escala aparte (QA-05, AD-18) & 15\,\% & 5 & 5 & 5 & 1 \\
Cliente maduro en la pila de AD-50 & 15\,\% & 5 & 4 & 3 & 5 \\
Experiencia del equipo & 10\,\% & 5 & 3 & 2 & 4 \\
Licencia y costo & 10\,\% & 4 & 5 & 5 & 5 \\
\midrule
\textbf{Total ponderado} & 100\,\% & \textbf{4,70} & 4,45 & 3,90 & 3,10 \\
\end{longtable}
\endgroup

\paragraph{Alternativas descartadas.}
\begin{itemize}
  \item Valkey gana en licencia, porque conserva la licencia BSD, pero pierde en experiencia del equipo y en madurez del cliente probado con la pila.
  \item Memcached pierde en invalidación por clave (QA-06), porque no ofrece estructuras ni invalidación por patrón.
  \item La vista materializada gana en cliente y licencia, pero pierde en componente separado (QA-05), porque concentra la lectura en la misma base.
\end{itemize}

\paragraph{Compromiso aceptado.} Redis 8 se distribuye con licencias RSALv2, SSPL y AGPLv3.
La AGPLv3 permite su uso en este proyecto académico.
Valkey es un sustituto compatible con el mismo protocolo si la licencia cambia.

\paragraph{Riesgo.} Riesgo R-14 del ATAM: la antigüedad del dato presentado depende de la carga.
En el perfil local se midió p95 de 1\,485\,ms con máximo de 11,8\,s.
Una corrida anterior dio p95 de 4,2\,s con 10 muestras.

\paragraph{Dónde se ve.} Vista física de la implementación: nodo \texttt{cache}.

\paragraph{Evidencia en el código.} QA-06 se midió con p95 de 142\,ms en el mapa de desastres y con la marca \texttt{datoAl} en el 100\,\% de las respuestas.
Archivos de evidencia:
{\raggedright
\begin{itemize}
  \item \texttt{proyectoparte2/\br{}servicios/\br{}consulta-danos/\br{}src/\br{}ConsultaDanos.Infraestructura/\br{}Cache.cs}
  \item \texttt{proyectoparte2/\br{}servicios/\br{}consulta/\br{}src/\br{}Consulta.Infraestructura/\br{}Adaptadores.cs}
\end{itemize}
\par}

%---------------------------------------------------------------------- AD-56
\subsection{AD-56. Almacén de archivos}
\label{sec:tec-ad56}

\paragraph{Contexto.} Esta decisión materializa el almacén de archivos de la vista física, donde se guardan las evidencias de Evaluación de Daños y los documentos de Reconstrucción.
El almacén debe mantener la evidencia disponible aunque falle una réplica del microservicio (AD-24, QA-02).
También debe conservar el 100\,\% de las fotos sincronizadas desde el cliente móvil (AD-31, QA-14).

\paragraph{Alternativas.} Se consideran MinIO, el disco local del contenedor, una columna binaria en PostgreSQL y SeaweedFS.
El \cref{tab:tec-ad56} presenta la matriz de decisión ponderada.
La alternativa elegida obtiene 4,20 y la segunda mejor, SeaweedFS, obtiene 4,05.

\begingroup\small
\begin{longtable}{L{4.96cm} C{1.10cm} C{2.25cm} C{2.25cm} C{2.25cm} C{2.25cm}}
\cabecera{6}{tab:tec-ad56}{Matriz de decisión ponderada de AD-56}{\textbf{Criterio (driver)} & \textbf{Peso} & \hspace{0pt}\textbf{MinIO} & \hspace{0pt}\textbf{Disco local} & \hspace{0pt}\textbf{Columna binaria} & \hspace{0pt}\textbf{SeaweedFS}}
Evidencia disponible ante la pérdida de una réplica (QA-02, AD-24) & 25\,\% & 4 & 1 & 4 & 4 \\
Integridad de la evidencia sincronizada (QA-14) & 20\,\% & 5 & 3 & 5 & 5 \\
Tamaño de la base y de su respaldo (QA-03, QA-06) & 15\,\% & 5 & 4 & 1 & 5 \\
Interfaz S3 estándar que permite sustituir el producto (QA-10) & 15\,\% & 5 & 1 & 1 & 4 \\
Experiencia del equipo & 10\,\% & 4 & 5 & 4 & 1 \\
Licencia y continuidad del producto & 15\,\% & 2 & 5 & 5 & 4 \\
\midrule
\textbf{Total ponderado} & 100\,\% & \textbf{4,20} & 2,85 & 3,45 & 4,05 \\
\end{longtable}
\endgroup

\paragraph{Alternativas descartadas.}
\begin{itemize}
  \item El disco local gana en experiencia y licencia, pero pierde en disponibilidad ante la pérdida de una réplica (QA-02), que AD-24 descarta.
  \item La columna binaria iguala en integridad (QA-14), pero pierde en tamaño de la base y de su respaldo (QA-03).
  \item SeaweedFS gana en continuidad del producto, pero pierde en experiencia del equipo y en compatibilidad S3 verificada.
\end{itemize}

\paragraph{Compromiso aceptado.} El proyecto original de MinIO dejó de publicar imágenes y el archivo de composición usa el fork comunitario \texttt{pgsty/minio} (DT-11).
El código depende solo de la API S3, lo que permite sustituir el producto sin modificar los microservicios.

\paragraph{Riesgo.} El fork comunitario puede dejar de mantenerse.
En el entorno de un computador \texttt{archivos} es un solo nodo, y su caída deja sin evidencias a las réplicas de \texttt{evaluacion-danos}.
No corresponde a un riesgo del ATAM.

\paragraph{Dónde se ve.} Vista física de la implementación: nodo \texttt{archivos}.

\paragraph{Evidencia en el código.} La corrida de QA-14 conservó 120 de 120 fotos a través del gateway.
Archivos de evidencia:
{\raggedright
\begin{itemize}
  \item \texttt{proyectoparte2/\br{}servicios/\br{}evaluacion-danos/\br{}src/\br{}EvaluacionDanos.Infraestructura/\br{}Archivos/\br{}AlmacenEvidenciasMinio.cs}
  \item \texttt{proyectoparte2/\br{}servicios/\br{}reconstruccion/\br{}src/\br{}Reconstruccion.Infraestructura/\br{}Archivos/\br{}AlmacenDocumentosMinio.cs}
\end{itemize}
\par}

%---------------------------------------------------------------------- AD-57
\subsection{AD-57. API Gateway}
\label{sec:tec-ad57}

\paragraph{Contexto.} Esta decisión materializa el punto único de entrada (AD-04), la priorización y degradación por tipo de cliente (AD-19) y el nivel 1 de la autorización (AD-22).
El gateway atiende QA-11, que exige el rechazo del 100\,\% de las modificaciones de un organismo no responsable.
También atiende QA-05 y QA-06, porque prioriza el tráfico crítico cuando la carga crece.

\paragraph{Alternativas.} Se consideran YARP, Kong, Traefik, Ocelot y nginx.
El \cref{tab:tec-ad57} presenta la matriz de decisión ponderada.
La alternativa elegida obtiene 4,85 y la segunda mejor, Kong, obtiene 3,45.

\begingroup\small
\begin{longtable}{L{3.93cm} C{1.10cm} C{1.95cm} C{1.95cm} C{1.95cm} C{1.95cm} C{1.95cm}}
\cabecera{7}{tab:tec-ad57}{Matriz de decisión ponderada de AD-57}{\textbf{Criterio (driver)} & \textbf{Peso} & \hspace{0pt}\textbf{YARP 2.3} & \hspace{0pt}\textbf{Kong} & \hspace{0pt}\textbf{Traefik} & \hspace{0pt}\textbf{Ocelot} & \hspace{0pt}\textbf{nginx}}
Priorización por clase de tráfico y degradación (QA-05, QA-06, AD-19) & 25\,\% & 5 & 3 & 2 & 3 & 2 \\
Validación JWT y regla de dominio por ruta (QA-11, AD-22) & 25\,\% & 5 & 4 & 3 & 4 & 2 \\
Catálogo de rutas validado al arrancar (QA-11) & 15\,\% & 5 & 4 & 3 & 3 & 2 \\
Latencia agregada por el intermediario (QA-06) & 15\,\% & 4 & 4 & 4 & 3 & 5 \\
Experiencia del equipo en la pila de AD-50 & 10\,\% & 5 & 2 & 3 & 4 & 3 \\
Licencia y mantenimiento & 10\,\% & 5 & 3 & 5 & 3 & 5 \\
\midrule
\textbf{Total ponderado} & 100\,\% & \textbf{4,85} & 3,45 & 3,10 & 3,35 & 2,85 \\
\end{longtable}
\endgroup

\paragraph{Alternativas descartadas.}
\begin{itemize}
  \item Kong pierde en priorización con degradación (QA-05), que exigiría complementos propios en otro lenguaje.
  \item Traefik pierde en priorización y en validación de la regla de dominio por ruta (QA-11).
  \item Ocelot pierde en catálogo validado al arrancar (QA-11) y en mantenimiento del proyecto.
  \item nginx gana en latencia agregada (QA-06), pero pierde en validación JWT con regla de dominio (QA-11), que no ofrece sin módulos adicionales.
\end{itemize}

\paragraph{Compromiso aceptado.} La priorización y la regla de dominio son código propio en el gateway.
Ese código se prueba con la misma suite de .NET, pero lo mantiene el equipo.

\paragraph{Riesgo.} Riesgo R-12 del ATAM: el gateway recibe tráfico sin TLS en el entorno de laboratorio (L-07).
Riesgo R-13 del ATAM: la validación depende de las llaves y del emisor de Keycloak (L-08).
Los rechazos del gateway no se registran en Auditoría (L-02).

\paragraph{Dónde se ve.} Vista lógica: componente API Gateway.
Vista física de la implementación: nodo \texttt{gateway} con dos réplicas y el artefacto \texttt{rutas.json}.

\paragraph{Evidencia en el código.} QA-11 se verificó con 218 de 218 intentos rechazados (100\,\%), 0 falsos negativos y 0 datos modificados sobre 12 rutas con regla de dominio.
Archivos de evidencia:
{\raggedright
\begin{itemize}
  \item \texttt{proyectoparte2/\br{}gateway/\br{}src/\br{}Plataforma.Gateway/\br{}rutas.json}
  \item \texttt{proyectoparte2/\br{}gateway/\br{}src/\br{}Plataforma.Gateway/\br{}Prioridad/\br{}Priorizacion.cs}
  \item \texttt{proyectoparte2/\br{}gateway/\br{}src/\br{}Plataforma.Gateway/\br{}Seguridad/\br{}AutorizacionNivel1.cs}
\end{itemize}
\par}

%---------------------------------------------------------------------- AD-58
\subsection{AD-58. Proveedor de identidad y clientes OIDC}
\label{sec:tec-ad58}

\paragraph{Contexto.} Esta decisión materializa el proveedor de identidad federado con OIDC y PKCE (AD-49) y la validación del token en el gateway y en cada microservicio (AD-04, AD-22).
Sus clientes son la consola, el cliente móvil y el tablero de daños. El tablero usa el cliente confidencial \texttt{consulta-\br{}danos-\br{}tablero} y una sesión con cookie que valida la propia página, porque su ruta es pública en el gateway (E-41).
El proveedor atiende QA-11, porque emite los roles y el atributo de organismo que usa la regla de dominio.
También atiende QA-14, porque el refresco del token sin conexión permite trabajar días en campo y sincronizar al volver.

\paragraph{Alternativas.} Se consideran Keycloak, Auth0, Duende IdentityServer y una identidad propia.
El \cref{tab:tec-ad58} presenta la matriz de decisión ponderada.
La alternativa elegida obtiene 4,90 y la segunda mejor, Duende IdentityServer, obtiene 3,75.

\begingroup\small
\begin{longtable}{L{4.96cm} C{1.10cm} C{2.25cm} C{2.25cm} C{2.25cm} C{2.25cm}}
\cabecera{6}{tab:tec-ad58}{Matriz de decisión ponderada de AD-58}{\textbf{Criterio (driver)} & \textbf{Peso} & \hspace{0pt}\textbf{Keycloak 26} & \hspace{0pt}\textbf{Auth0} & \hspace{0pt}\textbf{Duende IdentityServer} & \hspace{0pt}\textbf{Identidad propia}}
OIDC con PKCE para la consola y el móvil, con refresco sin conexión (QA-11, QA-14, AD-49) & 25\,\% & 5 & 5 & 4 & 2 \\
Roles y organismo como \emph{claims} para la regla de dominio (QA-11, AD-22) & 20\,\% & 5 & 4 & 4 & 3 \\
Despliegue dentro de la plataforma, sin servicio externo en la emergencia (QA-02, CON-3) & 15\,\% & 5 & 2 & 5 & 5 \\
Configuración declarada en un archivo versionado (QA-10) & 10\,\% & 5 & 3 & 4 & 3 \\
Experiencia del equipo & 10\,\% & 4 & 2 & 2 & 3 \\
Licencia y costo & 20\,\% & 5 & 2 & 3 & 5 \\
\midrule
\textbf{Total ponderado} & 100\,\% & \textbf{4,90} & 3,25 & 3,75 & 3,45 \\
\end{longtable}
\endgroup

\paragraph{Alternativas descartadas.}
\begin{itemize}
  \item Auth0 iguala en OIDC con PKCE (QA-11), pero pierde en despliegue dentro de la plataforma (CON-3) y en costo por usuario activo.
  \item Duende IdentityServer pierde en licencia, porque exige licencia comercial por encima de un umbral de ingresos, y en experiencia del equipo.
  \item La identidad propia gana en licencia y despliegue, pero pierde en OIDC con PKCE (QA-11), porque obligaría a construir la seguridad desde cero.
\end{itemize}

\paragraph{Compromiso aceptado.} El token se valida dos veces, en el gateway y en el microservicio, con un validador propio en \texttt{Plataforma.Comun.Seguridad}.
La doble validación agrega trabajo por petición a cambio de la defensa en profundidad de AD-22.

\paragraph{Riesgo.} Riesgo R-13 del ATAM: el emisor del token es \texttt{http://localhost:8180} y debe coincidir en el gateway y en los microservicios (L-08).
La recarga de las llaves JWT está pendiente (L-09).
Riesgo R-12 del ATAM: Keycloak corre en \texttt{start-dev} sin TLS (L-07).

\paragraph{Dónde se ve.} Vista lógica: componente externo Proveedor de identidad.
Vista física de la implementación: nodo \texttt{keycloak}.
Vista de implementación móvil: paquete \texttt{infraestructura/identidad}.

\paragraph{Evidencia en el código.} Siete pruebas instrumentadas de la app corren contra el gateway y Keycloak, incluido el rechazo por rol.
Archivos de evidencia:
{\raggedright
\begin{itemize}
  \item \texttt{proyectoparte2/\br{}despliegue/\br{}keycloak/\br{}realm-plataforma.json}
  \item \texttt{proyectoparte2/\br{}comun/\br{}Plataforma.Comun.Seguridad/\br{}IdentidadPorJwt.cs}
  \item \texttt{danos/\br{}app/\br{}src/\br{}main/\br{}java/\br{}com/\br{}example/\br{}danos/\br{}infraestructura/\br{}identidad/\br{}ConexionOidcBuilder.kt}
\end{itemize}
\par}

%---------------------------------------------------------------------- AD-59
\subsection{AD-59. Consola de coordinación}
\label{sec:tec-ad59}

\paragraph{Contexto.} Esta decisión materializa la capa de presentación web (AD-08) de la Consola de coordinación, con los módulos \texttt{mapa}, \texttt{evaluaciones} y \texttt{reconstruccion}.
La consola entra por el gateway (AD-04) y atiende QA-06 en la carga del mapa, QA-11 en el inicio de sesión y QA-10 en la incorporación de módulos.

\paragraph{Alternativas.} Se consideran Angular, React 19, Vue 3 y Blazor.
El \cref{tab:tec-ad59} presenta la matriz de decisión ponderada.
La alternativa elegida obtiene 4,80 y la segunda mejor, Vue 3, obtiene 4,10.

\begingroup\small
\begin{longtable}{L{4.96cm} C{1.10cm} C{2.25cm} C{2.25cm} C{2.25cm} C{2.25cm}}
\cabecera{6}{tab:tec-ad59}{Matriz de decisión ponderada de AD-59}{\textbf{Criterio (driver)} & \textbf{Peso} & \hspace{0pt}\textbf{Angular 22} & \hspace{0pt}\textbf{React 19} & \hspace{0pt}\textbf{Vue 3} & \hspace{0pt}\textbf{Blazor}}
Módulos con rutas, inyección y guardas (QA-10, AD-08) & 25\,\% & 5 & 3 & 4 & 4 \\
Cliente OIDC con PKCE certificado (QA-11, AD-49) & 20\,\% & 5 & 4 & 4 & 4 \\
Carga inicial y respuesta de la interfaz (QA-06) & 20\,\% & 4 & 4 & 5 & 2 \\
Archivos estáticos servidos detrás del gateway (QA-02, AD-04) & 10\,\% & 5 & 5 & 5 & 4 \\
Experiencia del equipo & 15\,\% & 5 & 3 & 2 & 3 \\
Licencia y costo & 10\,\% & 5 & 5 & 5 & 5 \\
\midrule
\textbf{Total ponderado} & 100\,\% & \textbf{4,80} & 3,80 & 4,10 & 3,55 \\
\end{longtable}
\endgroup

\paragraph{Alternativas descartadas.}
\begin{itemize}
  \item React 19 pierde en estructura por módulos con inyección y guardas (QA-10), que exige elegir y combinar bibliotecas.
  \item Vue 3 gana en carga inicial (QA-06), pero pierde en experiencia del equipo y en estructura por módulos.
  \item Blazor pierde en carga inicial (QA-06), por el tamaño del entorno WebAssembly que descarga el navegador.
\end{itemize}

\paragraph{Compromiso aceptado.} Angular produce un paquete inicial mayor que Vue.
La carga se mitiga con módulos cargados bajo demanda por ruta.

\paragraph{Riesgo.} Angular publica una versión mayor cada seis meses y obliga a actualizaciones frecuentes.
No corresponde a un riesgo del ATAM.

\paragraph{Dónde se ve.} Vista física de la implementación: nodo \texttt{web}.
Vista lógica: Consola de coordinación con los módulos \texttt{mapa}, \texttt{evaluaciones} y \texttt{reconstruccion}.

\paragraph{Evidencia en el código.} El \texttt{Dockerfile} compila con \texttt{node:22-alpine} y sirve los archivos con \texttt{nginx:1.29-alpine}.
Archivos de evidencia:
{\raggedright
\begin{itemize}
  \item \texttt{proyectoparte2/\br{}consola/\br{}package.json}
  \item \texttt{proyectoparte2/\br{}consola/\br{}Dockerfile}
  \item \texttt{proyectoparte2/\br{}consola/\br{}nginx.conf}
  \item \texttt{proyectoparte2/\br{}consola/\br{}src/\br{}app/\br{}modulos}
\end{itemize}
\par}

%---------------------------------------------------------------------- AD-60
\subsection{AD-60. Biblioteca de mapas}
\label{sec:tec-ad60}

\paragraph{Contexto.} Esta decisión materializa la presentación de las geometrías precalculadas (AD-17) y de su marca de tiempo (AD-16) en el módulo \texttt{mapa} de la consola.
La biblioteca atiende QA-06, porque el mapa es la consulta frecuente del coordinador.

\paragraph{Alternativas.} Se consideran Leaflet, OpenLayers, MapLibre GL y Google Maps.
El \cref{tab:tec-ad60} presenta la matriz de decisión ponderada.
La alternativa elegida obtiene 4,80 y la segunda mejor, OpenLayers, obtiene 3,90.

\begingroup\small
\begin{longtable}{L{4.96cm} C{1.10cm} C{2.25cm} C{2.25cm} C{2.25cm} C{2.25cm}}
\cabecera{6}{tab:tec-ad60}{Matriz de decisión ponderada de AD-60}{\textbf{Criterio (driver)} & \textbf{Peso} & \hspace{0pt}\textbf{Leaflet 1.9} & \hspace{0pt}\textbf{OpenLayers} & \hspace{0pt}\textbf{MapLibre GL} & \hspace{0pt}\textbf{Google Maps}}
Capas GeoJSON agregadas por zona (QA-06, AD-17) & 25\,\% & 5 & 5 & 5 & 4 \\
Peso de la biblioteca y tiempo de carga (QA-06) & 20\,\% & 5 & 3 & 2 & 3 \\
Independencia de un proveedor y de una llave (QA-02, CON-3) & 20\,\% & 4 & 4 & 4 & 1 \\
Integración con Angular y con el mapa del prototipo (QA-10) & 10\,\% & 5 & 3 & 3 & 3 \\
Experiencia del equipo & 10\,\% & 5 & 2 & 2 & 3 \\
Licencia y costo & 15\,\% & 5 & 5 & 5 & 2 \\
\midrule
\textbf{Total ponderado} & 100\,\% & \textbf{4,80} & 3,90 & 3,70 & 2,70 \\
\end{longtable}
\endgroup

\paragraph{Alternativas descartadas.}
\begin{itemize}
  \item OpenLayers iguala en capas GeoJSON (QA-06), pero pierde en peso de la biblioteca y en experiencia del equipo.
  \item MapLibre GL iguala en capas GeoJSON, pero pierde en peso de la biblioteca (QA-06), porque su ventaja con millones de puntos no hace falta con zonas agregadas.
  \item Google Maps pierde en independencia de un proveedor (CON-3) y en costo por carga de mapa.
\end{itemize}

\paragraph{Compromiso aceptado.} Leaflet dibuja en el navegador sin aceleración WebGL.
El volumen se controla con la agregación por zona del servidor (AD-17).

\paragraph{Riesgo.} Las teselas de OpenStreetMap no tienen acuerdo de servicio y su política de uso prohíbe el uso intensivo.
No corresponde a un riesgo del ATAM.

\paragraph{Dónde se ve.} Vista física de la implementación: nodo \texttt{web}.
Vista lógica: módulo \texttt{mapa} de la Consola de coordinación.

\paragraph{Evidencia en el código.} QA-06 se midió con p95 de 1\,628\,ms y p99 de 1\,935\,ms en el mapa de daños en el perfil local (L-04).
Archivos de evidencia:
{\raggedright
\begin{itemize}
  \item \texttt{proyectoparte2/\br{}consola/\br{}src/\br{}app/\br{}modulos/\br{}mapa/\br{}lienzo-mapa.ts}
  \item \texttt{proyectoparte2/\br{}consola/\br{}package.json}
\end{itemize}
\par}

%---------------------------------------------------------------------- AD-61
\subsection{AD-61. Plataforma del cliente móvil}
\label{sec:tec-ad61}

\paragraph{Contexto.} Esta decisión materializa la réplica local del cliente móvil (AD-31), la cola local de operaciones (AD-44) y el registro del desastre en campo (AD-48).
La plataforma atiende QA-14, cuya medida es 0 evaluaciones perdidas o sobrescritas en silencio al sincronizar 120 evaluaciones, y cuya respuesta exige 0 duplicados.

\paragraph{Alternativas.} Se consideran Android nativo con Kotlin, Flutter, React Native y Kotlin Multiplatform.
El \cref{tab:tec-ad61} presenta la matriz de decisión ponderada.
La alternativa elegida obtiene 4,60 y la segunda mejor, Kotlin Multiplatform, obtiene 4,00.

\begingroup\small
\begin{longtable}{L{4.96cm} C{1.10cm} C{2.25cm} C{2.25cm} C{2.25cm} C{2.25cm}}
\cabecera{6}{tab:tec-ad61}{Matriz de decisión ponderada de AD-61}{\textbf{Criterio (driver)} & \textbf{Peso} & \hspace{0pt}\textbf{Kotlin nativo} & \hspace{0pt}\textbf{Flutter} & \hspace{0pt}\textbf{React Native} & \hspace{0pt}\textbf{Kotlin Multiplatform}}
Trabajo en segundo plano y almacenamiento del sistema (QA-14, AD-44) & 25\,\% & 5 & 3 & 3 & 4 \\
Base local con transacciones y migraciones (QA-14, AD-31) & 20\,\% & 5 & 4 & 3 & 4 \\
Cámara y ubicación para las evidencias (QA-14) & 10\,\% & 5 & 4 & 4 & 4 \\
Reutilización del prototipo de FP-06 y de sus pruebas (QA-14) & 15\,\% & 5 & 1 & 1 & 4 \\
Portabilidad a otro sistema operativo móvil (QA-10) & 10\,\% & 1 & 5 & 5 & 4 \\
Experiencia del equipo & 10\,\% & 5 & 2 & 3 & 3 \\
Licencia y costo & 10\,\% & 5 & 5 & 5 & 5 \\
\midrule
\textbf{Total ponderado} & 100\,\% & \textbf{4,60} & 3,30 & 3,20 & 4,00 \\
\end{longtable}
\endgroup

\paragraph{Alternativas descartadas.}
\begin{itemize}
  \item Flutter gana en portabilidad (QA-10), pero pierde en reutilización del prototipo de FP-06 (QA-14), que habría que reescribir.
  \item React Native gana en portabilidad, pero pierde en base local transaccional (QA-14) y en reutilización del prototipo.
  \item Kotlin Multiplatform pierde en trabajo en segundo plano (QA-14), porque esa parte sigue siendo específica de cada plataforma.
\end{itemize}

\paragraph{Compromiso aceptado.} El cliente funciona solo en Android.
Un cliente para iOS exigiría otra implementación de la capa de presentación y de la sincronización.

\paragraph{Riesgo.} La app de entrega es la variante \texttt{debug}, que admite \texttt{http://} (L-07).
La evidencia solo admite fotos y no guarda un autor propio (L-10).

\paragraph{Dónde se ve.} Vista de implementación móvil: paquete \texttt{presentacion}.
Vista física de la implementación: \texttt{<<device>>} Teléfono Android con \texttt{DanosCampo.apk}.

\paragraph{Evidencia en el código.} La app \texttt{co.plataforma.danos} 2.1.0 tiene 81 pruebas unitarias y 16 instrumentadas en verde.
Archivos de evidencia:
{\raggedright
\begin{itemize}
  \item \texttt{danos/\br{}gradle/\br{}libs.versions.toml}
  \item \texttt{danos/\br{}app/\br{}build.gradle.kts}
\end{itemize}
\par}

%---------------------------------------------------------------------- AD-62
\subsection{AD-62. Base de datos local del cliente móvil}
\label{sec:tec-ad62}

\paragraph{Contexto.} Esta decisión materializa la réplica local (AD-31), la cola local de operaciones inmutables con copia completa del contenido (AD-44) y la tabla \texttt{registro\_desastre} del registro en campo (AD-48).
La base atiende QA-14, porque la escritura y su operación en la cola deben confirmarse en la misma transacción.

\paragraph{Alternativas.} Se consideran Room, Realm, SQLDelight y DataStore.
El \cref{tab:tec-ad62} presenta la matriz de decisión ponderada.
La alternativa elegida obtiene 4,90 y la segunda mejor, SQLDelight, obtiene 4,35.

\begingroup\small
\begin{longtable}{L{4.96cm} C{1.10cm} C{2.25cm} C{2.25cm} C{2.25cm} C{2.25cm}}
\cabecera{6}{tab:tec-ad62}{Matriz de decisión ponderada de AD-62}{\textbf{Criterio (driver)} & \textbf{Peso} & \hspace{0pt}\textbf{Room 2.8} & \hspace{0pt}\textbf{Realm} & \hspace{0pt}\textbf{SQLDelight} & \hspace{0pt}\textbf{DataStore}}
Transacción atómica entre la escritura y la cola local (QA-14, AD-44) & 30\,\% & 5 & 4 & 5 & 1 \\
Migraciones explícitas y verificables (QA-10, QA-14) & 20\,\% & 5 & 3 & 4 & 2 \\
Integración con Compose y corrutinas (AD-61) & 15\,\% & 5 & 4 & 4 & 4 \\
Continuidad del producto & 15\,\% & 5 & 1 & 5 & 5 \\
Experiencia del equipo & 10\,\% & 5 & 2 & 2 & 4 \\
Control del SQL verificado en compilación (QA-14) & 10\,\% & 4 & 2 & 5 & 1 \\
\midrule
\textbf{Total ponderado} & 100\,\% & \textbf{4,90} & 2,95 & 4,35 & 2,55 \\
\end{longtable}
\endgroup

\paragraph{Alternativas descartadas.}
\begin{itemize}
  \item Realm pierde en continuidad del producto, porque su fabricante retiró el SDK de dispositivos.
  \item SQLDelight gana en control del SQL verificado en compilación, pero pierde en migraciones con esquema exportado y en experiencia del equipo.
  \item DataStore pierde en transacción atómica con la cola (QA-14), porque guarda pares clave y valor sin relaciones.
\end{itemize}

\paragraph{Compromiso aceptado.} Cada cambio de esquema exige una migración escrita a mano y su esquema exportado.
Ese costo es el precio de no perder datos locales al actualizar la app (QA-14).

\paragraph{Riesgo.} Una migración defectuosa puede impedir abrir la base local con operaciones pendientes.
No corresponde a un riesgo del ATAM.

\paragraph{Dónde se ve.} Vista de implementación móvil: paquete \texttt{infraestructura/local}.
Vista física de la implementación: artefacto \texttt{danos.db} en el Teléfono Android.

\paragraph{Evidencia en el código.} \texttt{AppDatabase} declara la versión 3 con \texttt{MIGRACION\_1\_2} y \texttt{MIGRACION\_2\_3}, y la carpeta de esquemas guarda las versiones 1, 2 y 3.
Archivos de evidencia:
{\raggedright
\begin{itemize}
  \item \texttt{danos/\br{}app/\br{}src/\br{}main/\br{}java/\br{}com/\br{}example/\br{}danos/\br{}infraestructura/\br{}local/\br{}AppDatabase.kt}
  \item \texttt{danos/\br{}app/\br{}schemas}
\end{itemize}
\par}

%---------------------------------------------------------------------- AD-63
\subsection{AD-63. Sincronización del cliente móvil}
\label{sec:tec-ad63}

\paragraph{Contexto.} Esta decisión materializa la sincronización diferida (AD-31), el reenvío de la cola local (AD-44) y el envío idempotente del registro del desastre por el identificador del cliente (AD-48).
La sincronización atiende QA-14, cuya respuesta exige reenvío sin duplicados y 120 evaluaciones en $\le$ 10\,min.
También atiende QA-11, porque todo envío pasa por el gateway con el token OIDC.

\paragraph{Alternativas.} Se consideran WorkManager, SyncAdapter, un servicio en primer plano y Firebase.
El \cref{tab:tec-ad63} presenta la matriz de decisión ponderada.
La alternativa elegida obtiene 4,85 y la segunda mejor, SyncAdapter, obtiene 4,10.

\begingroup\small
\begin{longtable}{L{4.96cm} C{1.10cm} C{2.25cm} C{2.25cm} C{2.25cm} C{2.25cm}}
\cabecera{6}{tab:tec-ad63}{Matriz de decisión ponderada de AD-63}{\textbf{Criterio (driver)} & \textbf{Peso} & \hspace{0pt}\textbf{WorkManager} & \hspace{0pt}\textbf{SyncAdapter} & \hspace{0pt}\textbf{Servicio en primer plano} & \hspace{0pt}\textbf{Firebase}}
Ejecución garantizada con red disponible y reintento tras reinicio (QA-14, AD-44) & 30\,\% & 5 & 4 & 2 & 4 \\
Protocolo propio idempotente por identificador del cliente (QA-14, AD-48) & 25\,\% & 5 & 5 & 5 & 1 \\
Inmediatez del envío al recuperar la red (QA-14) & 15\,\% & 4 & 3 & 5 & 4 \\
Envío por el gateway con token OIDC (QA-11, AD-04) & 10\,\% & 5 & 5 & 5 & 1 \\
Experiencia del equipo & 10\,\% & 5 & 2 & 3 & 3 \\
Licencia y dependencia de un servicio externo & 10\,\% & 5 & 5 & 5 & 2 \\
\midrule
\textbf{Total ponderado} & 100\,\% & \textbf{4,85} & 4,10 & 3,90 & 2,65 \\
\end{longtable}
\endgroup

\paragraph{Alternativas descartadas.}
\begin{itemize}
  \item SyncAdapter pierde en ejecución garantizada (QA-14) y en experiencia del equipo, porque exige una cuenta del sistema y un proveedor de contenido.
  \item El servicio en primer plano gana en inmediatez del envío, pero pierde en ejecución garantizada tras un reinicio (QA-14).
  \item Firebase pierde en protocolo idempotente propio (QA-14) y en envío por el gateway (QA-11), porque sincroniza contra su propio servicio.
\end{itemize}

\paragraph{Compromiso aceptado.} El sistema operativo decide el instante de la ejecución de WorkManager.
El envío puede demorar unos segundos después de recuperar la red.

\paragraph{Riesgo.} Riesgo R-12 del ATAM: la variante \texttt{debug} admite tráfico sin TLS (L-07).
Riesgo R-13 del ATAM: desde un teléfono físico el emisor del token exige cambiar la configuración de Keycloak (L-08).

\paragraph{Dónde se ve.} Vista de implementación móvil: paquete \texttt{infraestructura/sync}.
Vista física de la implementación: \texttt{DanosCampo.apk} en el Teléfono Android.

\paragraph{Evidencia en el código.} QA-14 se verificó con 115 evaluaciones aplicadas, 5 en conflicto, 0 pérdidas, 0 sobrescrituras y 0 duplicados en 27,8\,s.
Archivos de evidencia:
{\raggedright
\begin{itemize}
  \item \texttt{danos/\br{}app/\br{}src/\br{}main/\br{}java/\br{}com/\br{}example/\br{}danos/\br{}infraestructura/\br{}sync/\br{}SincronizacionWorker.kt}
  \item \texttt{danos/\br{}app/\br{}src/\br{}main/\br{}java/\br{}com/\br{}example/\br{}danos/\br{}infraestructura/\br{}sync/\br{}EnviadorDesastres.kt}
\end{itemize}
\par}

%---------------------------------------------------------------------- AD-64
\subsection{AD-64. Biblioteca de resiliencia}
\label{sec:tec-ad64}

\paragraph{Contexto.} Esta decisión materializa el tiempo de espera, el reintento con espera creciente y el cortacircuitos de las fronteras (AD-05) y el cortacircuitos por adaptador (AD-37).
La biblioteca atiende QA-04, cuya respuesta exige cortar la propagación del fallo en $\le$ 60\,s.
También atiende QA-01, porque la caída de un sistema externo no debe degradar las demás funciones.

\paragraph{Alternativas.} Se consideran Microsoft.Extensions.Http.Resilience, una implementación propia y Steeltoe.
El \cref{tab:tec-ad64} presenta la matriz de decisión ponderada.
La alternativa elegida obtiene 4,90 y la segunda mejor, Implementación propia, obtiene 3,15.

\begingroup\small
\begin{longtable}{L{6.74cm} C{1.10cm} C{2.50cm} C{2.50cm} C{2.50cm}}
\cabecera{5}{tab:tec-ad64}{Matriz de decisión ponderada de AD-64}{\textbf{Criterio (driver)} & \textbf{Peso} & \hspace{0pt}\textbf{MS Resilience (Polly v8)} & \hspace{0pt}\textbf{Implementación propia} & \hspace{0pt}\textbf{Steeltoe}}
Tiempo de espera, reintento con variación aleatoria y cortacircuitos configurables (QA-01, AD-05) & 30\,\% & 5 & 3 & 2 \\
Corte de la propagación en $\le$ 60\,s verificable (QA-04) & 20\,\% & 5 & 3 & 3 \\
Integración con la fábrica de clientes HTTP de AD-50 & 15\,\% & 5 & 3 & 4 \\
Estado del circuito seguro ante concurrencia (QA-04) & 15\,\% & 5 & 2 & 4 \\
Experiencia del equipo & 10\,\% & 4 & 4 & 1 \\
Licencia y mantenimiento & 10\,\% & 5 & 5 & 4 \\
\midrule
\textbf{Total ponderado} & 100\,\% & \textbf{4,90} & 3,15 & 2,90 \\
\end{longtable}
\endgroup

\paragraph{Alternativas descartadas.}
\begin{itemize}
  \item La implementación propia iguala en experiencia y licencia, pero pierde en estado del circuito seguro ante concurrencia (QA-04).
  \item Steeltoe pierde en políticas configurables (QA-01) y en experiencia del equipo.
\end{itemize}

\paragraph{Compromiso aceptado.} Polly mide una proporción de fallos en una ventana, no un número de fallos consecutivos.
La regla de AD-05 se expresa como 3 llamadas mínimas con 100\,\% de fallos en 30\,s.

\paragraph{Riesgo.} Un reintento sobre una operación no idempotente duplicaría el efecto.
\texttt{Resiliencia.cs} reintenta solo operaciones idempotentes.
No corresponde a un riesgo del ATAM.

\paragraph{Dónde se ve.} Vista de implementación: proyecto \texttt{Plataforma.Comun.Adaptadores} (\texttt{Resiliencia.cs}).
Modelo de integración: cortacircuitos de cada adaptador.

\paragraph{Evidencia en el código.} La configuración usa 3\,s por intento, 2 reintentos, ventana de 30\,s y apertura de 15\,s. Las 12 pruebas de \texttt{AdaptadorSe6.Integracion} y las 14 de \texttt{AdaptadorSe7.Integracion} verifican el cortacircuitos y el reintento sin pérdida (QA-01).
Archivos de evidencia:
{\raggedright
\begin{itemize}
  \item \texttt{proyectoparte2/\br{}comun/\br{}Plataforma.Comun.Adaptadores/\br{}Resiliencia.cs}
  \item \texttt{proyectoparte2/\br{}Directory.Packages.props}
\end{itemize}
\par}

%---------------------------------------------------------------------- AD-65
\subsection{AD-65. Empaquetado, ejecución y escalado}
\label{sec:tec-ad65}

\paragraph{Contexto.} Esta decisión materializa las réplicas adicionales del camino de lectura (AD-18), la réplica activa con sondeo de salud (AD-24) y la ejecución híbrida con adaptadores efímeros (AD-25).
El empaquetado atiende QA-05, que exige escalar solo los componentes bajo presión sin modificar código.
También atiende QA-02 por el reemplazo de réplicas y QA-01 por el escalado de los adaptadores según su cola.

\paragraph{Alternativas.} Se consideran contenedores con Kubernetes, máquinas virtuales, Docker Swarm, Nomad y Knative.
El \cref{tab:tec-ad65} presenta la matriz de decisión ponderada.
La alternativa elegida obtiene 4,60 y la segunda mejor, Knative, obtiene 3,85.

\begingroup\small
\begin{longtable}{L{3.93cm} C{1.10cm} C{1.95cm} C{1.95cm} C{1.95cm} C{1.95cm} C{1.95cm}}
\cabecera{7}{tab:tec-ad65}{Matriz de decisión ponderada de AD-65}{\textbf{Criterio (driver)} & \textbf{Peso} & \hspace{0pt}\textbf{Docker + K8s (HPA, KEDA)} & \hspace{0pt}\textbf{Máquinas virtuales} & \hspace{0pt}\textbf{Docker Swarm} & \hspace{0pt}\textbf{Nomad} & \hspace{0pt}\textbf{Knative}}
Escalado automático selectivo sin modificar código (QA-05, AD-18) & 25\,\% & 5 & 2 & 2 & 4 & 5 \\
Escalado de 0 a N por largo de cola (QA-01, AD-25) & 15\,\% & 5 & 1 & 1 & 3 & 3 \\
Sondeo de salud y reemplazo de réplicas (QA-02, AD-24) & 20\,\% & 5 & 3 & 4 & 4 & 5 \\
Misma imagen en el computador y en el clúster (QA-09) & 15\,\% & 5 & 2 & 5 & 4 & 4 \\
Experiencia del equipo & 10\,\% & 4 & 3 & 3 & 1 & 1 \\
Costo de operación y licencia & 15\,\% & 3 & 2 & 5 & 3 & 3 \\
\midrule
\textbf{Total ponderado} & 100\,\% & \textbf{4,60} & 2,15 & 3,25 & 3,40 & 3,85 \\
\end{longtable}
\endgroup

\paragraph{Alternativas descartadas.}
\begin{itemize}
  \item Las máquinas virtuales pierden en escalado selectivo (QA-05) y en escalado por largo de cola (QA-01).
  \item Docker Swarm gana en costo de operación, pero pierde en escalado automático (QA-05), que no ofrece.
  \item Nomad pierde en escalado por cola y en licencia, porque cambió a BSL en 2023.
  \item Knative iguala en escalado selectivo y sondeo de salud, pero pierde en escalado por largo de cola (QA-01), porque está orientado a HTTP.
\end{itemize}

\paragraph{Compromiso aceptado.} En Docker Compose los adaptadores quedan siempre encendidos y no hay escalado automático.
El escalado de 0 a N solo existe en los manifiestos de Kubernetes (AD-25).

\paragraph{Riesgo.} El despliegue en Kubernetes no se verificó de punta a punta (L-05).
El escalado de 0 a N con KEDA no se midió (QA-01) y QA-05 no se midió con 10\,000 usuarios (L-03, L-04).

\paragraph{Dónde se ve.} Vista física de la implementación: nodo Docker Engine con los servicios del archivo de composición y anotación de Kubernetes con HPA y KEDA.

\paragraph{Evidencia en el código.} Los manifiestos declaran HPA de 2 a 6 réplicas para \texttt{consulta} y \texttt{consulta-\br{}danos}, de 2 a 5 para \texttt{consulta-\br{}programa}, y \texttt{ScaledObject} de KEDA de 0 a 4 réplicas para la salida de SE-6.
Archivos de evidencia:
{\raggedright
\begin{itemize}
  \item \texttt{proyectoparte2/\br{}despliegue/\br{}compose/\br{}docker-compose.yml}
  \item \texttt{proyectoparte2/\br{}despliegue/\br{}k8s/\br{}base/\br{}20-servicios.yaml}
  \item \texttt{proyectoparte2/\br{}despliegue/\br{}k8s/\br{}base/\br{}adaptador-se6.yaml}
\end{itemize}
\par}

%---------------------------------------------------------------------- AD-66
\subsection{AD-66. Plataforma del adaptador de SE-8}
\label{sec:tec-ad66}

\paragraph{Contexto.} CON-9, agregada el 19 de noviembre de 2026, exige que la solución final incluya Jakarta EE y .NET (E-44).
AD-03 ubica la heterogeneidad tecnológica (CON-2) en los adaptadores, no en los microservicios: un adaptador es el único punto de la arquitectura donde una pila distinta no toca el núcleo.
El adaptador de SE-8 estaba diseñado y no implementado (L-06).
QA-08 le exige una instancia encendida (AD-25), estado de entrega por alerta y \enquote{entregada} solo con acuse.
QA-09 pide que incorporar un sistema externo modifique 0 componentes existentes.
Esta decisión elige la plataforma de ese adaptador.

\paragraph{Alternativas.} Se consideran Jakarta EE 11 sobre WildFly 42.0.0.Final (imagen \texttt{quay.io/\br{}wildfly/\br{}wildfly:\br{}42.0.0.Final-\br{}jdk21}, Java 21), Spring Boot 3 sobre Java 21 y .NET 10 con \texttt{Plataforma.\br{}Comun.\br{}Adaptadores}, como los adaptadores de SE-6 y SE-7.
El \cref{tab:tec-ad66} presenta la matriz de decisión ponderada.
La alternativa elegida obtiene 4,35 y la segunda mejor, .NET 10, obtiene 4,00.

\begingroup\small
\begin{longtable}{L{6.0cm} C{1.10cm} C{2.6cm} C{2.6cm} C{2.6cm}}
\cabecera{5}{tab:tec-ad66}{Matriz de decisión ponderada de AD-66}{\textbf{Criterio (driver)} & \textbf{Peso} & \hspace{0pt}\textbf{Jakarta EE 11 (WildFly 42)} & \hspace{0pt}\textbf{Spring Boot 3} & \hspace{0pt}\textbf{.NET 10}}
Cumplimiento de CON-9 & 25\,\% & 5 & 2 & 1 \\
Adaptador nuevo sin modificar componentes existentes (QA-09, AD-03) & 20\,\% & 5 & 5 & 5 \\
Cortacircuitos del Cuadro 10 con aviso de apertura para EV-23 (QA-01, QA-08, AD-05, AD-37) & 15\,\% & 3 & 4 & 5 \\
Consumo transaccional: ack después del commit y DLQ a la quinta entrega (AD-02, AD-45) & 15\,\% & 4 & 4 & 5 \\
Instancia siempre encendida: memoria y arranque (QA-08, AD-25) & 10\,\% & 4 & 4 & 5 \\
Experiencia del equipo (especificaciones vistas en el curso) & 10\,\% & 4 & 3 & 5 \\
Licencia y costo & 5\,\% & 5 & 5 & 5 \\
\midrule
\textbf{Total ponderado} & 100\,\% & \textbf{4,35} & 3,65 & 4,00 \\
\end{longtable}
\endgroup

\paragraph{Alternativas descartadas.}
\begin{itemize}
  \item .NET 10 gana en resiliencia, mensajería y arranque, porque reutiliza \texttt{Plataforma.\br{}Comun.\br{}Adaptadores} y \texttt{ConsumidorIdempotente}. Pierde en el criterio que origina la decisión: no cumple CON-9.
  \item Spring Boot 3 corre sobre Java y usa algunas API de Jakarta, pero no es la plataforma Jakarta EE: cumple CON-9 solo en parte. Tiene un cortacircuitos con ventana de tiempo (Resilience4j) y es una tercera pila de marco que el equipo no vio en el curso.
\end{itemize}

\paragraph{Compromiso aceptado.} El adaptador usa las especificaciones estándar: Jakarta Messaging, EJB (\texttt{@Stateless}, \texttt{@Singleton}, Timer Service), Jakarta Persistence, Jakarta REST, CDI, Jakarta Concurrency y MicroProfile Config.
Dos piezas son propias y no vienen del estándar:
\begin{itemize}
  \item \textbf{Cortacircuitos propio en lugar de MicroProfile Fault Tolerance.} El \texttt{@CircuitBreaker} de MicroProfile mide una ventana de N llamadas, no la ventana de \qty{30}{\second} del Cuadro 10, y no avisa cuando el circuito se abre o se cierra. Ese aviso hace falta para publicar EV-23 por la bandeja (E-15). La clase \texttt{Cortacircuitos} aplica la regla de Polly (al menos 3 llamadas en \qty{30}{\second} con \qty{100}{\percent} de fallo, \qty{15}{\second} abierto y una sola llamada de prueba) y es segura ante concurrencia. Tiene 9 pruebas unitarias propias, una de ellas con 200 hilos.
  \item \textbf{EJB \texttt{@Singleton @Startup} en lugar de MDB.} Un MDB necesita un resource adapter JCA para el proveedor de mensajería, y RabbitMQ no tiene uno mantenido. El consumo usa igual la API estándar de Jakarta Messaging 3.1 (\texttt{Connection}, \texttt{Session} transaccional, \texttt{MessageConsumer}) con el cliente JMS de RabbitMQ en modo AMQP sobre la cola existente. El hilo lo da un \texttt{ManagedExecutorService}. Con \texttt{nackOnRollback}, \texttt{rollback()} manda el mensaje a la DLQ y cerrar la sesión lo reencola.
\end{itemize}
Se acepta además una segunda pila que mantener en la frontera. Esa pila duplica en Java la lectura del sobre y la parte del catálogo que usa el adaptador. Los esquemas JSON siguen siendo una sola fuente: el WAR los copia de \texttt{Plataforma.\br{}Comun.\br{}Eventos} al compilar.

\paragraph{Riesgo.} Un cambio del sobre canónico o de las claves de ruteo obliga a cambiar dos implementaciones. Lo detecta la validación contra los mismos esquemas, en las pruebas y al consumir. Corresponde al riesgo R-15 del ATAM.
Mientras Notificación no exista, el adaptador recibe EV-03 y EV-04 como sustituto temporal (L-11).

\paragraph{Dónde se ve.} Vista lógica: componente Adaptador de SE-8 [\texttt{adaptador-se8}] «implementado».
Vista de procesos: proceso \texttt{adaptador-se8} ×1, siempre encendido.
Vista física de la implementación: entorno de ejecución WildFly 42 dentro del contenedor \texttt{adaptador-se8} (\cref{fig:vista-despliegue-impl}, \cref{tab:nodos-impl}).
Vista de implementación: paquete \texttt{adaptador-\br{}se8-\br{}jee} (\cref{fig:vista-implementacion}).
Modelo de integración: fila de SE-8 del \cref{tab:integracion}.

\paragraph{Evidencia en el código.} 52 pruebas unitarias y 7 de integración en verde. Las de integración usan Testcontainers con la topología real del bus (cola quorum con \texttt{delivery-limit} 5) y la base del migrador. Demuestran el ack después del commit JPA, la DLQ a la quinta entrega fallida, el duplicado sin segundo envío y el circuito con EV-23.
La prueba \texttt{prueba-se8.sh} mide contra el compose: \qty{1,8}{\second} del registro a la alerta entregada con acuse, circuito abierto a los \qty{25,6}{\second} de la caída, y cola en 0 a los \qty{12,9}{\second} de la vuelta, con las 20 alertas entregadas una sola vez.
El \cref{tab:que-se-modifica} registra los 0 componentes existentes modificados.
Archivos de evidencia:
{\raggedright
\begin{itemize}
  \item \texttt{proyectoparte2/\br{}adaptadores/\br{}adaptador-se8-jee/\br{}adaptador/\br{}src/\br{}main/\br{}java/\br{}plataforma/\br{}adaptador/\br{}se8/\br{}jee/\br{}AdaptadorSe8.java}
  \item \texttt{proyectoparte2/\br{}adaptadores/\br{}adaptador-se8-jee/\br{}adaptador/\br{}src/\br{}main/\br{}java/\br{}plataforma/\br{}adaptador/\br{}se8/\br{}consumo/\br{}BucleConsumo.java}
  \item \texttt{proyectoparte2/\br{}adaptadores/\br{}adaptador-se8-jee/\br{}adaptador/\br{}src/\br{}main/\br{}java/\br{}plataforma/\br{}adaptador/\br{}se8/\br{}resiliencia/\br{}Cortacircuitos.java}
  \item \texttt{proyectoparte2/\br{}adaptadores/\br{}adaptador-se8-jee/\br{}Dockerfile}
  \item \texttt{proyectoparte2/\br{}pruebas/\br{}humo/\br{}prueba-se8.sh}
\end{itemize}
\par}

%------------------------------------------------------------------------------
\subsection{Restricciones de las tecnologías elegidas}
\label{sec:tec-restricciones}

Las tecnologías elegidas imponen límites al entorno implementado.
El \cref{tab:tec-restricciones} los relaciona con la decisión que los origina y con el escenario afectado.
Ninguno modifica una decisión estructural.

\begingroup\small
\begin{longtable}{L{2.6cm} L{1.3cm} L{6.4cm} L{1.7cm} L{3.3cm}}
\cabecera{5}{tab:tec-restricciones}{Restricciones que imponen las tecnologías elegidas}{\textbf{Restricción} &
\textbf{Origen} & \textbf{Efecto} & \textbf{QA afectado} & \textbf{Qué haría falta}}
Variante \texttt{debug} y Keycloak sin TLS (L-07) & AD-58, AD-61, AD-63 &
La app de entrega admite \texttt{http://}. Keycloak corre sin TLS. El tráfico entre el teléfono, el gateway y el proveedor de identidad viaja sin cifrar. Corresponde al riesgo R-12 del ATAM. &
QA-11, QA-12 & TLS en el borde y URL HTTPS en la variante \texttt{release} \\
Emisor del token en \texttt{localhost} (L-08) & AD-57, AD-58 &
El \texttt{iss} del token es \texttt{http://localhost:8180}. El emulador exige \texttt{adb reverse}. Un teléfono físico exige cambiar \texttt{KC\_HOSTNAME} e \texttt{Identidad\_\_Emisor} a la vez. Corresponde al riesgo R-13 del ATAM. &
QA-11, QA-14 & Un nombre DNS común para Keycloak \\
Imagen comunitaria de MinIO (DT-11) & AD-56 &
El proyecto original dejó de publicar imágenes. El archivo de composición usa el fork \texttt{pgsty/minio}, cuya continuidad no está garantizada. &
QA-02, QA-14 & Sustituir por otro almacén con API S3 sin cambiar el código \\
Política de uso de las teselas de OpenStreetMap & AD-60 &
Las teselas públicas no tienen acuerdo de servicio, exigen atribución y prohíben el uso intensivo. Durante una emergencia con 10\,000 usuarios la carga del mapa dependería de un tercero. &
QA-05, QA-06 & Servidor de teselas propio o proveedor con acuerdo \\
Keycloak en \texttt{start-dev} & AD-58 &
El modo de desarrollo usa una base embebida, no exige TLS y no está pensado para producción. El realm se importa desde \texttt{realm-plataforma.json} en cada arranque. Es una sola instancia. &
QA-02, QA-11 & Modo \texttt{start} con base PostgreSQL, TLS y dos instancias \\
Sustituto temporal de Notificación (L-11) & AD-66 &
EV-03 y EV-04 llegan a \texttt{q.salida.se8} por dos enlaces temporales, y el adaptador alerta a un único destino de difusión. No hay suscripciones ni resolución por ámbito, y el estado de entrega se consulta en el adaptador. &
QA-08 & Implementar Notificación (AD-13), quitar los enlaces y poner \texttt{Sustituto\_\_Activo=false} \\
Dos pilas en la frontera & AD-66 &
El adaptador de SE-8 duplica en Java la lectura del sobre y la parte del catálogo que usa. Un cambio del contrato se hace en .NET y en Java. Corresponde al riesgo R-15 del ATAM. &
QA-09, QA-10 & Mantener los esquemas JSON como única fuente y validarlos en las pruebas de las dos pilas \\
\end{longtable}
\endgroup
